\IfFileExists{stacks-project.cls}{%
\documentclass{stacks-project}
}{%
\documentclass{amsart}
}
\usepackage{amssymb}

% For dealing with references we use the comment environment
\usepackage{verbatim}
\usepackage{fancyvrb}
\newenvironment{reference}{\comment}{\endcomment}
%\newenvironment{reference}{}{}
\newenvironment{slogan}{\comment}{\endcomment}
\newenvironment{history}{\comment}{\endcomment}

% For source figures retained as assets
\usepackage{graphicx}

% For commutative diagrams we use Xy-pic
\usepackage[all]{xy}

% We use 2cell for 2-commutative diagrams.
\xyoption{2cell}
\UseAllTwocells

% We use multicol for the list of chapters between chapters
\usepackage{multicol}

% This is generally recommended for better output
\usepackage{lmodern}
\usepackage[T1]{fontenc}
\usepackage{textalpha}

% For cross-file-references
% We prefer the snapshot tmp/xref/#2.aux,
% which latexmkrc replaces atomically, so
% that documents compiled simultaneously
% never read a half-written .aux file.
% The last argument keeps links to #2.pdf.
\usepackage{xr-hyper}
\newcommand{\maybeexternaldocument}[2]{%
  \IfFileExists{#2.aux}{%
    \externaldocument[#1]{#2}%
  }{%
    \PackageWarning{stack}{Missing #2.aux; external references with prefix #1 unavailable}%
  }%
}
\let\externaldocumentorig\externaldocument
\renewcommand{\externaldocument}[2][]{%
  \IfFileExists{tmp/xref/#2.aux}{%
    \externaldocumentorig[#1]{tmp/xref/#2}[#2.pdf]%
  }{%
  \IfFileExists{#2.aux}{%
    \externaldocumentorig[#1]{#2}%
  }{%
    \PackageWarning{stack}{Missing #2.aux; external references with prefix #1 unavailable}%
  }%
  }%
}

% Package for hypertext links:
\usepackage{hyperref}

% For any local file, say "hello.tex" you want to link to please
% use \externaldocument[hello-]{hello}
\externaldocument[introduction-]{introduction}
\externaldocument[conventions-]{conventions}
\externaldocument[sets-]{sets}
\externaldocument[categories-]{categories}
\externaldocument[topology-]{topology}
\externaldocument[sheaves-]{sheaves}
\externaldocument[sites-]{sites}
\externaldocument[stacks-]{stacks}
\externaldocument[fields-]{fields}
\externaldocument[algebra-]{algebra}
\externaldocument[brauer-]{brauer}
\externaldocument[homology-]{homology}
\externaldocument[derived-]{derived}
\externaldocument[simplicial-]{simplicial}
\externaldocument[more-algebra-]{more-algebra}
\externaldocument[smoothing-]{smoothing}
\externaldocument[modules-]{modules}
\externaldocument[sites-modules-]{sites-modules}
\externaldocument[injectives-]{injectives}
\externaldocument[cohomology-]{cohomology}
\externaldocument[sites-cohomology-]{sites-cohomology}
\externaldocument[dga-]{dga}
\externaldocument[dpa-]{dpa}
\externaldocument[sdga-]{sdga}
\externaldocument[hypercovering-]{hypercovering}
\externaldocument[schemes-]{schemes}
\externaldocument[constructions-]{constructions}
\externaldocument[properties-]{properties}
\externaldocument[morphisms-]{morphisms}
\externaldocument[coherent-]{coherent}
\externaldocument[divisors-]{divisors}
\externaldocument[limits-]{limits}
\externaldocument[varieties-]{varieties}
\externaldocument[topologies-]{topologies}
\externaldocument[descent-]{descent}
\externaldocument[perfect-]{perfect}
\externaldocument[more-morphisms-]{more-morphisms}
\externaldocument[flat-]{flat}
\externaldocument[groupoids-]{groupoids}
\externaldocument[more-groupoids-]{more-groupoids}
\externaldocument[etale-]{etale}
\externaldocument[chow-]{chow}
\externaldocument[intersection-]{intersection}
\externaldocument[pic-]{pic}
\externaldocument[weil-]{weil}
\externaldocument[adequate-]{adequate}
\externaldocument[dualizing-]{dualizing}
\externaldocument[duality-]{duality}
\externaldocument[discriminant-]{discriminant}
\externaldocument[derham-]{derham}
\externaldocument[local-cohomology-]{local-cohomology}
\externaldocument[algebraization-]{algebraization}
\externaldocument[curves-]{curves}
\externaldocument[resolve-]{resolve}
\externaldocument[models-]{models}
\externaldocument[functors-]{functors}
\externaldocument[equiv-]{equiv}
\externaldocument[pione-]{pione}
\externaldocument[etale-cohomology-]{etale-cohomology}
\externaldocument[proetale-]{proetale}
\externaldocument[relative-cycles-]{relative-cycles}
\externaldocument[more-etale-]{more-etale}
\externaldocument[trace-]{trace}
\externaldocument[crystalline-]{crystalline}
\externaldocument[spaces-]{spaces}
\externaldocument[spaces-properties-]{spaces-properties}
\externaldocument[spaces-morphisms-]{spaces-morphisms}
\externaldocument[decent-spaces-]{decent-spaces}
\externaldocument[spaces-cohomology-]{spaces-cohomology}
\externaldocument[spaces-limits-]{spaces-limits}
\externaldocument[spaces-divisors-]{spaces-divisors}
\externaldocument[spaces-over-fields-]{spaces-over-fields}
\externaldocument[spaces-topologies-]{spaces-topologies}
\externaldocument[spaces-descent-]{spaces-descent}
\externaldocument[spaces-perfect-]{spaces-perfect}
\externaldocument[spaces-more-morphisms-]{spaces-more-morphisms}
\externaldocument[spaces-flat-]{spaces-flat}
\externaldocument[spaces-groupoids-]{spaces-groupoids}
\externaldocument[spaces-more-groupoids-]{spaces-more-groupoids}
\externaldocument[bootstrap-]{bootstrap}
\externaldocument[spaces-pushouts-]{spaces-pushouts}
\externaldocument[spaces-chow-]{spaces-chow}
\externaldocument[groupoids-quotients-]{groupoids-quotients}
\externaldocument[spaces-more-cohomology-]{spaces-more-cohomology}
\externaldocument[spaces-simplicial-]{spaces-simplicial}
\externaldocument[spaces-duality-]{spaces-duality}
\externaldocument[formal-spaces-]{formal-spaces}
\externaldocument[restricted-]{restricted}
\externaldocument[spaces-resolve-]{spaces-resolve}
\externaldocument[formal-defos-]{formal-defos}
\externaldocument[defos-]{defos}
\externaldocument[cotangent-]{cotangent}
\externaldocument[examples-defos-]{examples-defos}
\externaldocument[algebraic-]{algebraic}
\externaldocument[examples-stacks-]{examples-stacks}
\externaldocument[stacks-sheaves-]{stacks-sheaves}
\externaldocument[criteria-]{criteria}
\externaldocument[artin-]{artin}
\externaldocument[quot-]{quot}
\externaldocument[stacks-properties-]{stacks-properties}
\externaldocument[stacks-morphisms-]{stacks-morphisms}
\externaldocument[stacks-limits-]{stacks-limits}
\externaldocument[stacks-cohomology-]{stacks-cohomology}
\externaldocument[stacks-perfect-]{stacks-perfect}
\externaldocument[stacks-introduction-]{stacks-introduction}
\externaldocument[stacks-more-morphisms-]{stacks-more-morphisms}
\externaldocument[stacks-geometry-]{stacks-geometry}
\externaldocument[moduli-]{moduli}
\externaldocument[moduli-curves-]{moduli-curves}
\externaldocument[examples-]{examples}
\externaldocument[exercises-]{exercises}
\externaldocument[guide-]{guide}
\externaldocument[desirables-]{desirables}
\externaldocument[coding-]{coding}
\externaldocument[obsolete-]{obsolete}
\externaldocument[fdl-]{fdl}
\externaldocument[index-]{index}

% Heap Project documents
\externaldocument[k3-]{k3}
\externaldocument[torelli-]{torelli}
\externaldocument[linear-algebra-]{linear-algebra}
\externaldocument[groups-]{groups}
\externaldocument[complex-analysis-]{complex-analysis}
\externaldocument[functional-analysis-]{functional-analysis}
\externaldocument[categories-]{categories}
\externaldocument[commutative-algebra-]{commutative-algebra}
\externaldocument[ringed-spaces-]{ringed-spaces}
\externaldocument[homological-algebra-]{homological-algebra}
\externaldocument[lie-algebras-]{lie-algebras}
\externaldocument[representation-theory-]{representation-theory}
\externaldocument[smooth-manifolds-]{smooth-manifolds}
\externaldocument[differential-topology-]{differential-topology}
\externaldocument[algebraic-topology-]{algebraic-topology}
\externaldocument[differential-geometry-]{differential-geometry}
\externaldocument[algebraic-geometry-]{algebraic-geometry}
\externaldocument[symplectic-geometry-]{symplectic-geometry}
\externaldocument[complex-geometry-]{complex-geometry}
\externaldocument[hodge-theory-]{hodge-theory}
\externaldocument[%
variations-of-hodge-structures-]{%
variations-of-hodge-structures}
\externaldocument[%
homological-mirror-symmetry-]{%
homological-mirror-symmetry}
\maybeexternaldocument{deformations-}{deformations}
\externaldocument[vertex-algebras-]{vertex-algebras}
\externaldocument[springer-]{springer}
\externaldocument[infty-categories-]{infty-categories}
\externaldocument[seminars-differential-geometry-]{%
seminars-differential-geometry}
\externaldocument[seminars-complex-geometry-]{%
seminars-complex-geometry}
\externaldocument[seminars-algebraic-geometry-]{%
seminars-algebraic-geometry}
\externaldocument[seminars-others-]{%
seminars-others}
\externaldocument[%
geometric-prequantization-]{%
geometric-prequantization}
\externaldocument[surfaces-]{surfaces}
\externaldocument[birational-maps-commutative-algebra-]{birational-maps-commutative-algebra}
\externaldocument[cremona-transformations-]{cremona-transformations}
\externaldocument[derived-categories-of-sheaves-]{derived-categories-of-sheaves}
\externaldocument[geometric-invariant-theory-]{geometric-invariant-theory}
\externaldocument[geometric-stability-conditions-and-group-actions-]{geometric-stability-conditions-and-group-actions}
\externaldocument[moduli-spaces-of-sheaves-]{moduli-spaces-of-sheaves}
\externaldocument[bridgeland-stability-general-theory-]{bridgeland-stability-general-theory}
\externaldocument[hilbert-schemes-]{hilbert-schemes}
\externaldocument[noncommutative-abelian-surfaces-and-kummer-type-hyperkahler-varieties-]{noncommutative-abelian-surfaces-and-kummer-type-hyperkahler-varieties}
\externaldocument[mumford-tate-groups-in-hodge-theory-]{mumford-tate-groups-in-hodge-theory}
\externaldocument[hodge-birational-atoms-2026-]{hodge-birational-atoms-2026}
\externaldocument[xxii-egd-2026-]{%
xxii-egd-2026}
\externaldocument[pde-]{pde}
\externaldocument[probability-]{probability}

\externaldocument[linguistics-]{linguistics}
\externaldocument[genealogia-familiar-]{%
genealogia-familiar}

\externaldocument[%
16th-alga-meeting-2026-]{%
16th-alga-meeting-2026}

\externaldocument[%
iv-jornadas-lefschetz-2026-]{%
iv-jornadas-lefschetz-2026}

\externaldocument[reading-the-masters-]{%
reading-the-masters}

\externaldocument[real-analysis-]{%
real-analysis}

\externaldocument[optimal-transport-]{%
optimal-transport}

\externaldocument[history-]{history}

\externaldocument[seiberg-witten-]{seiberg-witten}

\externaldocument[history-philosophy-biology-]{history-philosophy-biology}

\externaldocument[chemistry-]{chemistry}

\externaldocument[molecular-cell-biology-]{molecular-cell-biology}

\externaldocument[microbiology-]{microbiology}

\externaldocument[protistology-]{protistology}

\externaldocument[mycology-]{mycology}

\externaldocument[botany-]{botany}

\externaldocument[zoology-]{zoology}

\externaldocument[animal-histology-]{animal-histology}

\externaldocument[animal-physiology-]{animal-physiology}

\externaldocument[systematics-biogeography-]{systematics-biogeography}

\externaldocument[ecology-natural-resources-]{ecology-natural-resources}

\externaldocument[earth-sciences-]{earth-sciences}

\externaldocument[biostatistics-]{biostatistics}

\externaldocument[mathematical-systems-biology-]{mathematical-systems-biology}

\externaldocument[cimpa-el-salvador-2024-]{cimpa-el-salvador-2024}

% Theorem environments.
%
\theoremstyle{plain}
\newtheorem{theorem}[subsection]{Theorem}
\newtheorem{proposition}[subsection]{Proposition}
\newtheorem{lemma}[subsection]{Lemma}
\newtheorem{corollary}[subsection]{Corollary}

\theoremstyle{definition}
\newtheorem{definition}[subsection]{Definition}
\newtheorem{example}[subsection]{Example}
\newtheorem{exercise}[subsection]{Exercise}
\newtheorem{situation}[subsection]{Situation}

\theoremstyle{remark}
\newtheorem{remark}[subsection]{Remark}
\newtheorem{remarks}[subsection]{Remarks}

\numberwithin{equation}{subsection}

% Macros
%
\def\lim{\mathop{\mathrm{lim}}\nolimits}
\def\colim{\mathop{\mathrm{colim}}\nolimits}
\def\Spec{\mathop{\mathrm{Spec}}}
\def\Hom{\mathop{\mathrm{Hom}}\nolimits}
\def\Ext{\mathop{\mathrm{Ext}}\nolimits}
\def\SheafHom{\mathop{\mathcal{H}\!\mathit{om}}\nolimits}
\def\SheafExt{\mathop{\mathcal{E}\!\mathit{xt}}\nolimits}
\def\Sch{\mathit{Sch}}
\def\Mor{\mathop{\mathrm{Mor}}\nolimits}
\def\Ob{\mathop{\mathrm{Ob}}\nolimits}
\def\Sh{\mathop{\mathit{Sh}}\nolimits}
\def\NL{\mathop{N\!L}\nolimits}
\def\CH{\mathop{\mathrm{CH}}\nolimits}
\def\proetale{{pro\text{-}\acute{e}tale}}
\def\etale{{\acute{e}tale}}
\def\QCoh{\mathit{QCoh}}
\def\Ker{\mathop{\mathrm{Ker}}}
\def\Im{\mathop{\mathrm{Im}}}
\def\Coker{\mathop{\mathrm{Coker}}}
\def\Coim{\mathop{\mathrm{Coim}}}

% Boxtimes
%
\DeclareMathSymbol{\boxtimes}{\mathbin}{AMSa}{"02}

%
% Macros for moduli stacks/spaces
%
\def\QCohstack{\mathcal{QC}\!\mathit{oh}}
\def\Cohstack{\mathcal{C}\!\mathit{oh}}
\def\Spacesstack{\mathcal{S}\!\mathit{paces}}
\def\Quotfunctor{\mathrm{Quot}}
\def\Hilbfunctor{\mathrm{Hilb}}
\def\Curvesstack{\mathcal{C}\!\mathit{urves}}
\def\Polarizedstack{\mathcal{P}\!\mathit{olarized}}
\def\Complexesstack{\mathcal{C}\!\mathit{omplexes}}
% \Pic is the operator that assigns to X its picard group, usage \Pic(X)
% \Picardstack_{X/B} denotes the Picard stack of X over B
% \Picardfunctor_{X/B} denotes the Picard functor of X over B
\def\Pic{\mathop{\mathrm{Pic}}\nolimits}
\def\Picardstack{\mathcal{P}\!\mathit{ic}}
\def\Picardfunctor{\mathrm{Pic}}
\def\Deformationcategory{\mathcal{D}\!\mathit{ef}}


\begin{document}

\title{CIMPA School El Salvador 2024}
\maketitle

\phantomsection
\label{section-phantom}

\tableofcontents

\medskip\noindent
{\bf Conference information.}

\href{%
https://naturales.ues.edu.sv/%
ecimpasv2024/}{%
Escuela CIMPA El Salvador 2024:
Teoría de Subvariedades y Análisis
Geométrico}.
On the
\href{https://cimpa.info/en/node/7295}{%
CIMPA page}: Submanifold Theory and
Geometric Analysis.

8 to 19 July 2024.

Auditorio de la Facultad de CCNN y
Matemática ``Dr. Fabio Castillo'',
Universidad de El Salvador,
San Salvador, El Salvador.

Information accessed November 1, 2026.

\medskip\noindent
{\bf Organizers.}

Organizing Committee: Dr. Riquelmi
Cardona, Dr. Dimas Tejada, M.Sc. Martín
Guerra, M.Sc. Ingrid Martínez, Lic.
Ernesto Hidalgo, Lic. Mirna Galdámez.
Local Coordinator: Dr. Riquelmi Cardona.
External Coordinator: Dr. Luis Alías.

\medskip\noindent
{\bf Programme.}

The official
\href{%
https://naturales.ues.edu.sv/%
ecimpasv2024/schedule/}{%
schedule} lists three introductory
courses in the first week and three
advanced courses in the second, each in
four sessions of 90 minutes. It also lists
two talks (``Charla 1'' and ``Charla 2''),
student talks and workshops, whose titles
and speakers are not listed on the
official site as of November 1, 2026; they
have no sections here. The official
language of the school was Spanish. These
notes are in English; the abstracts are
quoted in Spanish as published.

\medskip\noindent
{\bf Related material in HP.}

Riemannian background, including the
Cartan--Hadamard theorem, is in the
Differential Geometry chapter, Section
\ref{differential-geometry-%
section-Hadamard}.
Alexandrov's theorem on constant mean
curvature surfaces, with an anisotropic
version, appears in Section
\ref{xxii-egd-2026-section-%
alexandrov-crystalline}
of the XXII Escola de Geometria
Diferencial notes. A Finsler metric
appears in the notes on magnetic
geodesics, Section
\ref{seminars-differential-geometry-%
section-magnetic-geodesics-low-energy}.
The Physics notebook (not part of the
book) has short sections on Lorentz
transformations and general relativity.

\section{Introducción a la Geometría
Lorentziana de Curvas y Superficies}
\label{section-palmas-lorentzian-curves}

\medskip\noindent
{\bf Curso Introductorio 1.}

Oscar Palmas,
UNAM-México.

\medskip\noindent
{\bf Schedule.}

Week 1, Monday to Thursday, 11:00--12:30
(8--11 July 2024).

\medskip\noindent
{\bf Official listing.}

\href{%
https://naturales.ues.edu.sv/%
ecimpasv2024/introductory-curse-1/}{%
Programme entry.}

\medskip\noindent
{\bf Abstract.}

El objetivo de este minicurso es motivar
geométricamente a la Geometría Lorentziana.
Primero daremos algunos preliminares sobre
Álgebra Lineal Indefinida en espacios
vectoriales reales, analizando varios temas
sobre carácter causal, grupo de isometrías y
formas canónicas de operadores lineales
auto-adjuntos. Luego revisaremos rápidamente
la teoría de curvas en el espacio
tridimensional de Lorenz-Minkowski.
Concluimos con el estudio geométrico de las
superficies, también en el caso
tridimensional.

\medskip\noindent
{\bf Sources of these notes.}

These notes are an English version of
Daniel's notes taken during the course,
in the order of the lectures. No lecture
notes or slides of the course were
available. Statements were checked
against López's survey
\cite{lopez2014differential} (theorem
numbers as in its arXiv version) and
against O'Neill's book \cite{oneill}.
Several statements of the notes taken in
class were incomplete or garbled; they
were completed or corrected, and some
short proofs and computations were added.
Where this changes a statement
substantially, the text says so. Points
that could not be verified are marked
``(not checked)''. The notes use the
Lorentzian cross product, Frenet frames
and the shape operator; the Riemannian
theory of isometric immersions, with
the shape operator and the Gauss equation,
is in Section
\ref{differential-geometry-%
section-isometric-immersions}
of the Differential Geometry chapter.

\medskip\noindent
{\bf Lorentzian linear algebra.}

Throughout, $V$ is a real vector space of
finite dimension.

\begin{definition}[Scalar product]
\label{definition-palmas-scalar-product}
A \emph{scalar product} on $V$ is a
symmetric bilinear form
$\left\langle{} \cdot,\cdot \right\rangle{}$
on $V$ which is \emph{nondegenerate}: if
$\left\langle{} v,w \right\rangle{}=0$ for
all $w\in V$, then $v=0$.
\end{definition}

\begin{definition}[Lorentz--Minkowski space]
\label{definition-palmas-lorentz-minkowski}
The \emph{Lorentz--Minkowski space}
$\mathbb{R}^{n+1}_1$ is $\mathbb{R}^{n+1}$,
with coordinates $(x_0,x_1,\dots,x_n)$,
equipped with the scalar product
$$
\left\langle{} u,v \right\rangle{}
=-u_0v_0+u_1v_1+\dots+u_nv_n.
$$
\end{definition}

More generally, $\mathbb{R}^m_s$ denotes
$\mathbb{R}^m$ with the scalar product in
which the first $s$ squares carry a minus
sign. A \emph{Lorentzian vector space} is
a vector space with a scalar product of
index $1$ (the index is defined below) and
dimension at least $2$. In
$\mathbb{R}^3_1$ we also write the
coordinates as $(t,x,y)$.

\begin{proposition}
\label{proposition-palmas-nondegenerate}
Let $V$ be a vector space with a scalar
product.
\begin{enumerate}
\item If
$\left\langle{} u,w \right\rangle{}
=\left\langle{} v,w \right\rangle{}$
for all $w\in V$, then $u=v$.
\item If $V\neq 0$, then $V$ has an
orthonormal basis.
\end{enumerate}
\end{proposition}

\begin{proof}
For the first item, $u-v$ is orthogonal to
every vector, so $u-v=0$ by
nondegeneracy. The second item is standard
\cite{oneill}.
\end{proof}

The second item was cut off in the notes
taken in class; we completed it with the
standard statement.

\begin{definition}[Causal character]
\label{definition-palmas-causal-character}
Let $V$ be a vector space with a scalar
product. A vector $v\in V$ is
\begin{enumerate}
\item \emph{spacelike} if $v=0$ or
$\left\langle{} v,v \right\rangle{}>0$;
\item \emph{timelike} if
$\left\langle{} v,v \right\rangle{}<0$;
\item \emph{lightlike} or \emph{null} if
$v\neq 0$ and
$\left\langle{} v,v \right\rangle{}=0$;
\item \emph{causal} if $v\neq 0$ and
$\left\langle{} v,v \right\rangle{}\leq 0$.
\end{enumerate}
The \emph{light cone} or \emph{null cone}
of $V$ is the set of null vectors of $V$.
The \emph{time cone} of $V$ is the set of
timelike vectors of $V$. For a null vector
$u$, the \emph{null cone of $u$},
denoted $CN(u)$, is the connected component
of the null cone containing $u$.
Analogously, for a timelike vector $u$ the
\emph{time cone of $u$}, $CT(u)$, is the
component of the time cone containing $u$,
and for a causal vector $u$ the
\emph{causal cone of $u$}, $CC(u)$, is the
component of the set of causal vectors
containing $u$.
\end{definition}

In $\mathbb{R}^{n+1}_1$ the time cone and
the set of causal vectors each have two
components, given by $x_0>0$ (the
\emph{future}) and $x_0<0$ (the
\emph{past}). The null cone has two
components if $n\geq 2$, and four (open
half-lines) if $n=1$. Two timelike vectors
$u,v$ lie in the same time cone if and
only if
$\left\langle{} u,v \right\rangle{}<0$
\cite[Section 1.2]{lopez2014differential}.

\begin{proposition}
\label{proposition-palmas-cones-convex}
In $\mathbb{R}^{n+1}_1$, the cones $CT(u)$
and $CC(u)$ are convex. If $n\geq 2$, the
cone $CN(u)$ is not convex.
\end{proposition}

\begin{proof}
Write a vector as $(x_0,x')$ with
$x'\in\mathbb{R}^n$, and let $|x'|$ be the
Euclidean norm of $x'$. The future time
cone is the set where $x_0>|x'|$. Since
$x'\mapsto |x'|$ is a convex function, this
set is convex. The future causal cone is the
set where $x_0\geq |x'|$ and $x_0>0$; it is
convex for the same reason. The past cones
are their images under $v\mapsto -v$.
If $n\geq 2$, the vectors $(1,1,0,\dots,0)$
and $(1,-1,0,\dots,0)$ are null and in the
same component, but their sum
$(2,0,\dots,0)$ is timelike.
\end{proof}

The notes taken in class stated that all
three cones $CN(u)$, $CT(u)$ and $CC(u)$
are convex; this is false for $CN(u)$, as
the proof shows.

\begin{example}
\label{example-palmas-orthogonal-plane}
In $\mathbb{R}^2_1$ with coordinates
$(t,x)$, the vectors orthogonal to
$(1,1)$ form the line $x=t$, which
contains $(1,1)$ itself. The vectors
orthogonal to $(1,0)$ form the line $t=0$,
as in the Euclidean plane. In general the
orthogonal line to $(a,b)$ is spanned by
$(b,a)$: it is the reflection of the line
spanned by $(a,b)$ in the diagonal $x=t$.
\end{example}

\begin{definition}[Lorentzian cross product]
\label{definition-palmas-cross-product}
Let $e_0,e_1,e_2$ be the canonical basis of
$\mathbb{R}^3_1$ and let
$u=(u_0,u_1,u_2)$ and $v=(v_0,v_1,v_2)$.
The \emph{Lorentzian cross product} of $u$
and $v$ is the formal determinant
$$
u\times v=\det
\begin{pmatrix}
-e_0 & e_1 & e_2\\
u_0 & u_1 & u_2\\
v_0 & v_1 & v_2
\end{pmatrix}.
$$
\end{definition}

From now on this is simply called the
cross product. Explicitly,
$$
\begin{aligned}
u\times v=\big(
&u_2v_1-u_1v_2,\,
u_2v_0-u_0v_2,\\
&u_0v_1-u_1v_0\big).
\end{aligned}
$$

\begin{proposition}[Properties of the
cross product]
\label{proposition-palmas-cross-product}
For $u,v,w\in\mathbb{R}^3_1$:
\begin{enumerate}
\item
$\left\langle{} u\times v,w \right\rangle{}
=\left\langle{} u,v\times w \right\rangle{}
=\det(u,v,w)$, where $\det(u,v,w)$ is the
determinant of the matrix with rows $u$,
$v$, $w$;
\item $v\times u=-u\times v$, and
$u\times v$ is orthogonal to $u$ and $v$;
\item (Lagrange identity)
$$
\begin{aligned}
&\left\langle{} u\times v,u\times v
\right\rangle{}\\
&\quad=\left\langle{} u,v \right\rangle{}^2
-\left\langle{} u,u \right\rangle{}
\left\langle{} v,v \right\rangle{}.
\end{aligned}
$$
\end{enumerate}
\end{proposition}

\begin{proof}
The first item follows by expanding both
sides in coordinates (the minus sign in
front of $e_0$ is compensated by the minus
sign in the scalar product). The second
item follows from the first, since a
determinant with two equal rows vanishes.
The Lagrange identity is checked by direct
expansion in coordinates.
\end{proof}

In the notes taken in class the
determinant in the first item was left as
``$\det(?)$''. This agrees with López's
cross product \cite[Section
1.2]{lopez2014differential}, who puts the
timelike coordinate last.

\begin{definition}
\label{definition-palmas-nondeg-subspace}
Let $V$ be a vector space with a scalar
product. A subspace $W$ of $V$ is
\emph{nondegenerate} if the restriction of
the scalar product of $V$ to $W$ is again a
scalar product.
\end{definition}

\begin{proposition}[Subspaces]
\label{proposition-palmas-subspaces}
Let $V$ be a vector space with a scalar
product and $W$ a subspace of $V$.
\begin{enumerate}
\item $\dim W+\dim W^{\perp}=\dim V$.
\item $(W^{\perp})^{\perp}=W$.
\item $W$ is nondegenerate if and only if
$V=W\oplus W^{\perp}$.
\item $W$ is nondegenerate if and only if
$W^{\perp}$ is nondegenerate.
\end{enumerate}
\end{proposition}

These are standard \cite{oneill}; for a
line in $\mathbb{R}^3_1$ compare
\cite[Proposition 1.1]{%
lopez2014differential}.
The last item was cut off in the notes
taken in class; we completed it with the
standard statement.

\begin{definition}[Orthonormal set]
\label{definition-palmas-orthonormal-set}
Let $V$ be a vector space with a scalar
product. A set of vectors $\{v_i\}$ of $V$
is \emph{orthonormal} if
$\left\langle{} v_i,v_j \right\rangle{}=0$
for $i\neq j$ and
$\left\langle{} v_i,v_i \right\rangle{}
=\pm 1$ for all $i$.
\end{definition}

\begin{definition}[Pseudo-orthonormal basis]
\label{definition-palmas-pseudo-orthonormal}
Let $V$ be a Lorentzian vector space of
dimension $n+1\geq 2$. A basis
$\{v_1,\dots,v_{n+1}\}$ of $V$ is
\emph{pseudo-orthonormal} if
$v_1$ and $v_2$ are null,
$$
\left\langle{} v_1,v_2 \right\rangle{}=1,
$$
$v_1$ and $v_2$ are orthogonal to
$v_3,\dots,v_{n+1}$, and
$$
\left\langle{} v_i,v_j \right\rangle{}
=\delta_{ij}
$$
for $i,j=3,\dots,n+1$.
\end{definition}

\begin{remark}
\label{remark-palmas-null-pair-sign}
If we put $-1$ instead of $1$ in the
condition on
$\left\langle{} v_1,v_2 \right\rangle{}$,
then $v_1$ and $v_2$ lie both in the future
or both in the past. With $1$, as above,
one of them is in the future and the other
in the past. For instance, in
$\mathbb{R}^2_1$ the future null vectors
$(1,1)$ and $(1,-1)$ have scalar product
$-2$.
\end{remark}

\begin{proposition}
\label{proposition-palmas-null-expansion}
Let $V$ be a Lorentzian vector space and
$\{v_i\}$ a pseudo-orthonormal basis of
$V$. Every $v\in V$ satisfies
$$
\begin{aligned}
v={}&\left\langle{} v,v_2 \right\rangle{}v_1
+\left\langle{} v,v_1 \right\rangle{}v_2\\
&+\sum_{i=3}^{n+1}
\left\langle{} v,v_i \right\rangle{}v_i.
\end{aligned}
$$
\end{proposition}

\begin{proof}
Write
$v=a v_1+bv_2+\sum_{i\geq 3}c_iv_i$.
Taking the scalar product with $v_1$
gives
$\left\langle{} v,v_1 \right\rangle{}=b$,
with $v_2$ gives
$\left\langle{} v,v_2 \right\rangle{}=a$,
and with $v_i$, $i\geq 3$, gives
$\left\langle{} v,v_i \right\rangle{}=c_i$.
\end{proof}

So, compared with the usual expansion in an
orthonormal basis, the roles of $v_1$ and
$v_2$ are exchanged.

\begin{definition}[Index]
\label{definition-palmas-index}
Let $V$ be a vector space with a scalar
product and $\{e_i\}$ an orthonormal basis
of $V$. The number of timelike vectors in
$\{e_i\}$ is the \emph{index} of the scalar
product.
\end{definition}

\begin{proposition}[Sylvester's law of
inertia]
\label{proposition-palmas-sylvester}
The number of timelike vectors in an
orthonormal basis of $V$ does not depend on
the basis. It equals the largest dimension
of a subspace of $V$ on which the scalar
product is negative definite.
\end{proposition}

This is standard \cite{oneill}. The
statement in the notes taken in class was
garbled; we replaced it by this standard
statement.

\begin{definition}[Antinormal set]
\label{definition-palmas-antinormal}
Let $V$ be a vector space with a scalar
product and $\varepsilon=\pm 1$. A set of
vectors $\{v_1,\dots,v_p\}$ in $V$ is
\emph{antinormal} with sign $\varepsilon$
if
$\left\langle{} v_i,v_j \right\rangle{}
=\varepsilon$ when $i+j=p+1$ and
$\left\langle{} v_i,v_j \right\rangle{}=0$
otherwise.
\end{definition}

So the Gram matrix of an antinormal set
has $\varepsilon$ on the antidiagonal and
zeros elsewhere.

\begin{theorem}
\label{theorem-palmas-antinormal}
Let $\{v_1,\dots,v_p\}$ be an antinormal
set with sign $\varepsilon$ in a vector
space $V$ with a scalar product. Then:
\begin{enumerate}
\item $\{v_1,\dots,v_p\}$ is linearly
independent;
\item its span $W$ is nondegenerate;
\item the index of $W$ is
$\lfloor p/2\rfloor$ if $\varepsilon=1$
and $\lceil p/2\rceil$ if
$\varepsilon=-1$.
\end{enumerate}
\end{theorem}

\begin{proof}
If $\sum_i a_iv_i=0$, then taking the
scalar product with $v_j$ gives
$\varepsilon a_{p+1-j}=0$, so all $a_i$
vanish. The Gram matrix $\varepsilon J$ of
the basis $\{v_i\}$ of $W$, where $J$ is
the antidiagonal matrix of ones, is
invertible, so $W$ is nondegenerate. The
vectors $e_i\pm e_{p+1-i}$ for
$i<p+1-i$, where the $e_i$ are the standard
basis vectors of $\mathbb{R}^p$, together
with the middle vector
$e_{(p+1)/2}$ when $p$ is odd, form an
eigenbasis of $J$, with eigenvalue $+1$
for the sign $+$ and for the middle
vector, and $-1$ for the sign $-$. So $J$
has $\lfloor p/2\rfloor$ negative
eigenvalues, and $-J$ has
$\lceil p/2\rceil$.
\end{proof}

The second item was left unfinished in the
notes taken in class; the second and third
items, with their proof, were supplied
here. In a Lorentzian space an antinormal
set therefore has at most three vectors if
$\varepsilon=1$ and at most two if
$\varepsilon=-1$; such sets appear as bases
in types II and III of the canonical forms
below.

\begin{definition}[Linear isometry]
\label{definition-palmas-linear-isometry}
A linear map $T\colon V\to W$ between
vector spaces with scalar products is a
\emph{linear isometry} if it is bijective
and
$\left\langle{} Tu,Tv \right\rangle{}
=\left\langle{} u,v \right\rangle{}$
for all $u,v\in V$. If such a map exists,
$V$ and $W$ are \emph{isometric}.
\end{definition}

\begin{proposition}
\label{proposition-palmas-polarization}
Let $T\colon V\to W$ be a linear map
between vector spaces with scalar products.
\begin{enumerate}
\item $T$ preserves scalar products if and
only if
$\left\langle{} Tv,Tv \right\rangle{}
=\left\langle{} v,v \right\rangle{}$
for all $v\in V$.
\item If $T$ preserves scalar products,
then $T$ is injective. In particular, if
$\dim V=\dim W$, then $T$ is a linear
isometry.
\end{enumerate}
\end{proposition}

\begin{proof}
The first item follows from the
polarization identity
$$
\left\langle{} u,v \right\rangle{}
=\tfrac{1}{2}\big(q(u+v)-q(u)-q(v)\big),
$$
where
$q(v)=\left\langle{} v,v \right\rangle{}$.
If $Tv=0$, then
$\left\langle{} v,u \right\rangle{}
=\left\langle{} Tv,Tu \right\rangle{}=0$
for all $u$, so $v=0$.
\end{proof}

The proposition recorded in class read ``a
bijective linear map preserves the scalar
product if and only if it is a linear
map'', which is garbled; the statement
above is the standard one.

\begin{theorem}
\label{theorem-palmas-isometric-spaces}
Two vector spaces with scalar products are
isometric if and only if they have the same
dimension and the same index.
\end{theorem}

\begin{proof}
A linear isometry maps orthonormal bases to
orthonormal bases and preserves the sign of
each vector, so dimension and index are
preserved
(Proposition~\ref{%
proposition-palmas-sylvester}).
Conversely, order orthonormal bases of both
spaces with the timelike vectors first;
the linear map sending one basis to the
other is a linear isometry.
\end{proof}

\begin{corollary}
\label{corollary-palmas-standard-model}
Every vector space with a scalar product of
index $s$ and dimension $m$ is isometric to
$\mathbb{R}^m_s$.
\end{corollary}

\begin{proposition}[Reverse Cauchy--Schwarz
inequality]
\label{proposition-palmas-reverse-cs}
Let $V$ be a Lorentzian vector space and
$v,w\in V$ timelike. Then
$$
\left\langle{} v,w \right\rangle{}^2
\geq\left\langle{} v,v \right\rangle{}
\left\langle{} w,w \right\rangle{},
$$
with equality if and only if $v$ and $w$
are linearly dependent.
\end{proposition}

\begin{proof}
Since $v$ is timelike, $v^{\perp}$ is
spacelike and $V=\mathbb{R}v\oplus
v^{\perp}$. Write $w=av+w'$ with
$w'\in v^{\perp}$, so that
$\left\langle{} w',w' \right\rangle{}\geq 0$.
Then
$\left\langle{} v,w \right\rangle{}
=a\left\langle{} v,v \right\rangle{}$
and
$$
\left\langle{} w,w \right\rangle{}
=a^2\left\langle{} v,v \right\rangle{}
+\left\langle{} w',w' \right\rangle{}.
$$
Hence
$$
\begin{aligned}
\left\langle{} v,w \right\rangle{}^2
&=\left\langle{} v,v \right\rangle{}
\big(\left\langle{} w,w \right\rangle{}
-\left\langle{} w',w' \right\rangle{}\big)\\
&\geq\left\langle{} v,v \right\rangle{}
\left\langle{} w,w \right\rangle{},
\end{aligned}
$$
because
$\left\langle{} v,v \right\rangle{}<0$.
Equality holds if and only if
$\left\langle{} w',w' \right\rangle{}=0$,
that is, $w'=0$.
\end{proof}

In terms of the norm
$|v|=\sqrt{|\left\langle{} v,v
\right\rangle{}|}$, the inequality reads
$|\left\langle{} v,w \right\rangle{}|\geq
|v|\,|w|$. The notes taken in class had the
inequality in the wrong direction and
attributed it to negative definite scalar
products; the statement above is the
correct one \cite[Theorem
1.1]{lopez2014differential}. If $v$ and
$w$ lie in the same time cone, there is a
unique $\varphi\geq 0$, the
\emph{hyperbolic angle} between them, with
$$
\left\langle{} v,w \right\rangle{}
=-|v|\,|w|\cosh\varphi .
$$

\begin{proposition}[Reverse triangle
inequality]
\label{proposition-palmas-reverse-triangle}
Let $V$ be a Lorentzian vector space and
let $v,w$ be timelike vectors in the same
time cone. Then
$$
|v|+|w|\leq |v+w|,
$$
with equality if and only if $v$ and $w$
are linearly dependent.
\end{proposition}

\begin{proof}
By Corollary~\ref{%
corollary-palmas-standard-model} we may
assume $V=\mathbb{R}^{n+1}_1$, since an
isometry preserves time cones and null
vectors; the same applies in the next
proposition.
The vector $v+w$ is timelike and in the
same time cone, by convexity
(Proposition~\ref{%
proposition-palmas-cones-convex}).
Since
$\left\langle{} v,w \right\rangle{}<0$,
Proposition~\ref{%
proposition-palmas-reverse-cs}
gives
$-\left\langle{} v,w \right\rangle{}
\geq |v|\,|w|$. Therefore
$$
\begin{aligned}
|v+w|^2
&=|v|^2-2\left\langle{} v,w
\right\rangle{}+|w|^2\\
&\geq\big(|v|+|w|\big)^2,
\end{aligned}
$$
with equality exactly in the equality case
of Proposition~\ref{%
proposition-palmas-reverse-cs}.
\end{proof}

The notes taken in class have only the
title of this proposition; the statement
and proof were supplied here.

\begin{definition}[Causal character of a
subspace]
\label{definition-palmas-causal-subspace}
Let $V$ be a Lorentzian vector space and
$W$ a subspace of $V$.
\begin{enumerate}
\item $W$ is \emph{spacelike} if $W$ is
nondegenerate of index $0$, like a plane
that meets the light cone only at the
origin.
\item $W$ is \emph{timelike} if $W$ is
nondegenerate of index $1$, like a plane
that meets the light cone in two lines.
\item $W$ is \emph{lightlike} or
\emph{null} if $W$ is degenerate, like a
plane tangent to the light cone along one
line.
\end{enumerate}
\end{definition}

\begin{proposition}
\label{proposition-palmas-subspace-criteria}
Let $V$ be a Lorentzian vector space and
$W$ a subspace of $V$.
\begin{enumerate}
\item $W$ is timelike if and only if $W$
contains a timelike vector.
\item $W$ is lightlike if and only if $W$
contains a null vector but no timelike
vector.
\item If $\dim W\geq 2$, then $W$ is
timelike if and only if $W$ contains two
linearly independent null vectors.
\item If $\dim W\geq 2$ and $W$ contains no
null vectors, then $W$ is spacelike.
\end{enumerate}
\end{proposition}

\begin{proof}
(1) If $t\in W$ is timelike, then
$W=\mathbb{R}t\oplus(t^{\perp}\cap W)$
and $t^{\perp}$ is spacelike, so $W$ is
nondegenerate of index $1$. The converse
holds by definition.
(2) If $W$ is degenerate, a nonzero $u$ in
$W\cap W^{\perp}$ is null. If $W$ also
contained a timelike $t$, then $u$ would
be a nonzero vector of the spacelike
space $t^{\perp}$, a contradiction.
Conversely, if $W$ contains a null vector
but no timelike vector, it is neither
spacelike nor timelike, so it is
degenerate.
(3) If $t\in W$ is a unit timelike vector,
choose a unit vector $e\in t^{\perp}\cap W$;
then $t+e$ and $t-e$ are independent and
null. Conversely, two independent null
vectors $u,v$ are not orthogonal
(Exercise~\ref{%
exercise-palmas-orthogonal-null}),
and $u+\lambda v$ is timelike for a
suitable sign of $\lambda$, since its
square is
$2\lambda\left\langle{} u,v
\right\rangle{}$.
(4) By (2) $W$ is not lightlike, and by (3)
it is not timelike.
\end{proof}

The proposition in the notes taken in class
was incomplete; the four statements above
are the standard ones (compare
\cite[Section 1.1]{lopez2014differential}
for $\mathbb{R}^3_1$), with short proofs
supplied here.

\begin{definition}[Boost]
\label{definition-palmas-boost}
Let $\theta\in\mathbb{R}$. The
\emph{boost} $E_{\theta}$ of
$\mathbb{R}^2_1$ is the operator whose
matrix in the canonical basis is
$$
\begin{pmatrix}
\cosh\theta & \sinh\theta\\
\sinh\theta & \cosh\theta
\end{pmatrix}.
$$
\end{definition}

A boost is a ``hyperbolic rotation'': it
is a linear isometry of $\mathbb{R}^2_1$
with determinant $1$ that maps each branch
of the hyperbolas $-t^2+x^2=c$ to itself.
Conversely, if $T$ is a linear isometry of
$\mathbb{R}^2_1$ with determinant $1$ that
maps the future time cone to itself, then
its first column is a future unit timelike
vector, hence equal to
$(\cosh\theta,\sinh\theta)$; its second
column is a unit spacelike vector
orthogonal to it, hence
$\pm(\sinh\theta,\cosh\theta)$, and the
determinant forces the sign $+$. So $T$ is
a boost.

\begin{definition}[Orthogonal group]
\label{definition-palmas-orthogonal-group}
The \emph{orthogonal group}
$\operatorname{O}(V)$ of a vector space
$V$ with a scalar product is the group of
linear isometries of $V$.
\end{definition}

\begin{definition}[Orthochronous]
\label{definition-palmas-orthochronous}
Let $V$ be a Lorentzian vector space. A
linear isometry $T\colon V\to V$ is
\emph{orthochronous} if it preserves the
time orientation, that is, it maps each
time cone to itself (equivalently, it maps
a future-pointing timelike vector to a
future-pointing one). The orthochronous
isometries form the \emph{orthochronous
orthogonal group} $\operatorname{O}^+(V)$,
and those which moreover preserve the
orientation form the \emph{special
orthochronous orthogonal group}
$\operatorname{SO}^+(V)$.
\end{definition}

For $V=\mathbb{R}^3_1$ the group
$\operatorname{O}(V)$ has four connected
components, distinguished by the
determinant and by whether the isometry is
orthochronous
\cite[Theorem 1.2]{lopez2014differential}.

\begin{definition}[Self-adjoint operator]
\label{definition-palmas-self-adjoint}
Let $V$ be a vector space with a scalar
product. A linear operator
$T\colon V\to V$ is \emph{self-adjoint} if
$$
\left\langle{} Tu,v \right\rangle{}
=\left\langle{} u,Tv \right\rangle{}
$$
for all $u,v\in V$.
\end{definition}

\begin{proposition}[Properties of
self-adjoint operators]
\label{proposition-palmas-self-adjoint}
Let $T$ be a self-adjoint operator on a
vector space $V$ with a scalar product.
Matrices act on columns: the $j$-th column
of $M$ holds the coordinates of $Tv_j$.
\begin{enumerate}
\item If $M$ is the matrix of $T$ in a
basis and $G$ is the Gram matrix of that
basis, then $M^{t}G=GM$.
\item Eigenvectors of $T$ with different
eigenvalues are orthogonal.
\item If $W$ is a $T$-invariant subspace,
then so is $W^{\perp}$.
\end{enumerate}
\end{proposition}

\begin{proof}
(1) The entries of $M^tG$ and $GM$ are
$\left\langle{} Tv_i,v_j \right\rangle{}$
and $\left\langle{} v_i,Tv_j
\right\rangle{}$.
(2) If $Tu=\lambda u$ and $Tv=\mu v$, then
$\lambda\left\langle{} u,v \right\rangle{}
=\mu\left\langle{} u,v \right\rangle{}$.
(3) If $v\in W^{\perp}$ and $w\in W$, then
$\left\langle{} Tv,w \right\rangle{}
=\left\langle{} v,Tw \right\rangle{}=0$.
\end{proof}

The notes taken in class have only the
title of this proposition; the three
properties above were supplied here.

\begin{example}
\label{example-palmas-self-adjoint-plane}
In an orthonormal basis $\{e_0,e_1\}$ of
$\mathbb{R}^2_1$ with $e_0$ timelike, the
self-adjoint operators are exactly those
with matrix
$$
\begin{pmatrix}
a & b\\
-b & c
\end{pmatrix}.
$$
Their characteristic polynomial has
discriminant $(a-c)^2-4b^2$. For example,
if $a=c$ and $b\neq 0$, the eigenvalues
$a\pm ib$ are not real, so the operator is
not diagonalizable over $\mathbb{R}$.
\end{example}

The matrix in the notes taken in class was
garbled; the one above follows from
Proposition~\ref{%
proposition-palmas-self-adjoint}.
In contrast with the Euclidean case,
self-adjoint operators are not always
diagonalizable.

\begin{theorem}[Canonical forms]
\label{theorem-palmas-canonical-forms}
Let $T$ be a self-adjoint operator on a
Lorentzian vector space $V$ of dimension
$n$. Then $V$ has a basis in which $T$
takes one of the following four forms.
\begin{enumerate}
\item[I.] $T$ is diagonal in an
orthonormal basis.
\item[II.] There is a basis
$\{u,v,e_3,\dots,e_n\}$ with $u,v$ null,
$\left\langle{} u,v \right\rangle{}=1$,
and $e_3,\dots,e_n$ orthonormal, spacelike,
orthogonal to $u,v$, and eigenvectors of
$T$, such that the matrix of $T$ on the
span of $u,v$ is
$$
\begin{pmatrix}
a_0 & 0\\
\delta & a_0
\end{pmatrix},
$$
with $\delta=\pm 1$.
\item[III.] There is a basis
$\{u,v,e_3,\dots,e_n\}$ as in II, with
$n\geq 3$, such that the matrix of $T$ on
the span of $u,v,e_3$ is
$$
\begin{pmatrix}
a_0 & 0 & 0\\
0 & a_0 & 1\\
1 & 0 & a_0
\end{pmatrix},
$$
and $e_4,\dots,e_n$ are eigenvectors.
\item[IV.] There is an orthonormal basis
$\{e_1,\dots,e_n\}$ with $e_1$ timelike,
whose last $n-2$ vectors are
eigenvectors of $T$, such that the matrix
of $T$ on the span of $e_1,e_2$ is
$$
\begin{pmatrix}
a_0 & b_0\\
-b_0 & a_0
\end{pmatrix},
$$
with $b_0\neq 0$.
\end{enumerate}
\end{theorem}

In type II the eigenvector $v$ is null and
$T-a_0$ has a Jordan block of size $2$; in
type III a Jordan block of size $3$; in
type IV there is a pair of non-real
eigenvalues $a_0\pm ib_0$. The notes taken
in class only say that there are four
cases. We have checked that the four
matrices are self-adjoint for the given
bases and that the sign $\delta$ in type II
cannot be normalized: the sign of
$\left\langle{} (T-a_0)x,x \right\rangle{}$
does not depend on the basis. The full
statement in dimension $n$ was not checked
against a source (not checked); the case
of dimension $2$, which is the one used for
surfaces, is proved next.

\begin{proposition}
\label{proposition-palmas-canonical-plane}
Let $T$ be a self-adjoint operator on a
Lorentzian vector space of dimension $2$.
Then $T$ has one of the forms I, II or IV
of
Theorem~\ref{theorem-palmas-canonical-forms}.
\end{proposition}

\begin{proof}
In an orthonormal basis $\{e_1,e_2\}$ with
$e_1$ timelike, $T$ has a matrix as in
Example~\ref{%
example-palmas-self-adjoint-plane};
let $D=(a-c)^2-4b^2$.
If $D>0$, the two eigenvalues are real and
different, and the eigenvectors are
orthogonal
(Proposition~\ref{%
proposition-palmas-self-adjoint}).
Neither is null: the vectors orthogonal to
a null vector of a Lorentzian plane are its
multiples. Normalizing them gives type I.
If $D=0$ and $T=a_0\operatorname{id}$, we
are in type I. If $D=0$ and $T\neq
a_0\operatorname{id}$, then $N=T-a_0$ is
nonzero with $N^2=0$, so
$\operatorname{im}N=\ker N$ is spanned by
an eigenvector $v$. For every $x$,
$\left\langle{} Nx,v \right\rangle{}
=\left\langle{} x,Nv \right\rangle{}=0$,
so $v$ is orthogonal to
$\operatorname{im}N$, which contains $v$;
hence $v$ is null. Choose a null $u$ with
$\left\langle{} u,v \right\rangle{}=1$;
then $Nu=cv$ with $c\neq 0$, and replacing
$u,v$ by $\lambda u,\lambda^{-1}v$ with
$\lambda^2=|c|^{-1}$ gives type II.
If $D<0$, then $b\neq 0$. Replace $e_1,e_2$
by $\cosh\theta\,e_1+\sinh\theta\,e_2$ and
$\sinh\theta\,e_1+\cosh\theta\,e_2$. A
computation shows that the difference of
the diagonal entries becomes
$$
(a-c)\cosh 2\theta+2b\sinh 2\theta,
$$
which vanishes when
$\tanh 2\theta=-(a-c)/(2b)$; this is
possible because $|a-c|<2|b|$. The new
matrix is self-adjoint with equal diagonal
entries, so it has the form IV, and
$b_0\neq 0$ because the eigenvalues are
not real.
\end{proof}

\medskip\noindent
{\bf Curves in $\mathbb{R}^3_1$.}

\begin{definition}[Regular curve]
\label{definition-palmas-regular-curve}
A \emph{differentiable parametrized curve}
is a differentiable map
$\alpha\colon I\to\mathbb{R}^{n+1}_1$,
where $I$ is an open interval. It is
\emph{regular} if $\alpha'(s)\neq 0$ for
all $s\in I$.
\end{definition}

\begin{definition}[Causal character of a
curve]
\label{definition-palmas-causal-curve}
A regular curve $\alpha$ is
\emph{timelike} (resp.\
\emph{spacelike}, \emph{null} or
\emph{lightlike}) if $\alpha'(s)$ is
timelike (resp.\ spacelike, null) for all
$s\in I$.
\end{definition}

The notes taken in class defined null
curves by $\alpha'(s)=0$, which is not
possible for a regular curve; the
condition is that $\alpha'(s)$ is a null
vector.

\begin{remark}
\label{remark-palmas-null-curves-exist}
The notes taken in class contain the
remark ``there are no regular parametrized
null curves in $\mathbb{R}^3_1$''. This is
false: the helix
$$
\alpha(s)=(s,\cos s,\sin s)
$$
has $\alpha'(s)=(1,-\sin s,\cos s)$, which
is null, and null straight lines are
further examples. What is true is that a
null curve has length zero, so it cannot
be parametrized by arc length (see below).
Also, a regular curve need not have a
causal character: for
$\alpha(s)=(s,s^2/2,0)$ we get
$\left\langle{} \alpha'(s),\alpha'(s)
\right\rangle{}=s^2-1$, which changes sign.
\end{remark}

For timelike and spacelike curves we can
measure length with the usual definition,
using the norm $|v|=\sqrt{|\left\langle{}
v,v \right\rangle{}|}$, and reparametrize
by arc length, so that
$\left\langle{} \alpha',\alpha'
\right\rangle{}=\pm 1$.

What do we do with null curves? We use the
second derivative.

\begin{proposition}
\label{proposition-palmas-null-curves}
Let $\alpha$ be a null curve in
$\mathbb{R}^3_1$. For each $s$, either
$\alpha''(s)$ is a multiple of
$\alpha'(s)$, or $\alpha''(s)$ is
spacelike and nonzero.
\begin{enumerate}
\item If $\alpha''(s)$ is a multiple of
$\alpha'(s)$ for all $s$ (for instance if
$\alpha''(s)$ is null or zero), then
$\alpha$ is a \emph{pregeodesic}: its
image lies on a null straight line.
\item If $\alpha''(s)$ is spacelike and
nonzero for all $s$, then $\alpha$ can be
reparametrized to a curve $\beta$ with
$\left\langle{} \beta'',\beta''
\right\rangle{}=1$.
\end{enumerate}
\end{proposition}

\begin{proof}
Differentiating
$\left\langle{} \alpha',\alpha'
\right\rangle{}=0$ gives
$\left\langle{} \alpha',\alpha''
\right\rangle{}=0$. The plane
$\alpha'(s)^{\perp}$ is lightlike and
contains $\alpha'(s)$; its vectors are
spacelike or multiples of $\alpha'(s)$.
(1) If $\alpha''=f\alpha'$, then
$\alpha'(s)=e^{F(s)}\alpha'(s_0)$ with
$F'=f$, so $\alpha$ moves along the line
through $\alpha(s_0)$ in the direction
$\alpha'(s_0)$.
(2) If $\beta(t)=\alpha(\varphi(t))$, then
$\beta''=\varphi''\alpha'
+(\varphi')^2\alpha''$, and since $\alpha'$
is null and orthogonal to $\alpha''$,
$$
\left\langle{} \beta'',\beta''
\right\rangle{}
=(\varphi')^4\left\langle{}
\alpha'',\alpha'' \right\rangle{}.
$$
Choose $\varphi$ with
$\varphi'=\left\langle{} \alpha'',\alpha''
\right\rangle{}^{-1/4}$.
\end{proof}

\begin{definition}[Pseudo-arc length]
\label{definition-palmas-pseudo-arc-length}
A null curve $\beta$ with
$\left\langle{} \beta'',\beta''
\right\rangle{}=1$ is
\emph{parametrized by pseudo-arc length}.
\end{definition}

In the notes taken in class the first case
was stated as ``$\alpha''$ null implies
$\alpha'$ and $\alpha''$ dependent
(exercise)''; the exercise is
Exercise~\ref{%
exercise-palmas-orthogonal-null}.
The proof above was supplied here. The
helix of
Remark~\ref{remark-palmas-null-curves-exist}
is already parametrized by pseudo-arc
length.

\begin{definition}[Frenet--Serret frame]
\label{definition-palmas-frenet-frame}
Let $\alpha\colon I\to\mathbb{R}^3_1$ be a
spacelike or timelike curve parametrized
by arc length, such that $\alpha''(s)$ is
nonzero and not null for all $s$. The
\emph{curvature} of $\alpha$ is
$\kappa(s)=|\alpha''(s)|$. The
\emph{Frenet--Serret frame} at $s$ is the
basis $\{T(s),N(s),B(s)\}$ of
$\mathbb{R}^3_1$ where $T=\alpha'$ is the
tangent, $N=\alpha''/\kappa$ is the
normal, and the binormal $B$ is the unique
vector making $\{T,N,B\}$ an orthonormal,
positively oriented basis.
\end{definition}

\begin{remark}
\label{remark-palmas-binormal}
Let $\varepsilon_T$, $\varepsilon_N$ and
$\varepsilon_B$ be the values
$\left\langle{} T,T \right\rangle{}$,
$\left\langle{} N,N \right\rangle{}$ and
$\left\langle{} B,B \right\rangle{}$.
Exactly one of $T,N,B$ is timelike, so
$\varepsilon_B=-\varepsilon_T
\varepsilon_N$. By the Lagrange identity
(Proposition~\ref{%
proposition-palmas-cross-product}),
$$
\left\langle{} T\times N,T\times N
\right\rangle{}
=-\varepsilon_T\varepsilon_N
=\varepsilon_B,
$$
and $\det(T,N,T\times N)=\varepsilon_B$.
Hence the positively oriented binormal is
$$
B=\varepsilon_B\,T\times N.
$$
\end{remark}

\begin{definition}[Torsion]
\label{definition-palmas-torsion}
The \emph{torsion} of $\alpha$ is
$$
\tau(s)=-\left\langle{} N'(s),B(s)
\right\rangle{}
=\left\langle{} B'(s),N(s)
\right\rangle{}.
$$
\end{definition}

The two expressions agree because
$\left\langle{} N,B \right\rangle{}=0$.
The notes taken in class do not write down
the Frenet equations; with the
conventions above they are as follows
(derived here).

\begin{proposition}[Frenet equations]
\label{proposition-palmas-frenet-equations}
For $\alpha$ as in
Definition~\ref{%
definition-palmas-frenet-frame},
$$
\begin{aligned}
T'&=\kappa N,\\
N'&=-\varepsilon_T\varepsilon_N\kappa T
-\varepsilon_B\tau B,\\
B'&=\varepsilon_N\tau N.
\end{aligned}
$$
Explicitly: if $\alpha$ is timelike,
$$
N'=\kappa T-\tau B,\qquad B'=\tau N;
$$
if $\alpha$ and $N$ are spacelike,
$$
N'=-\kappa T+\tau B,\qquad B'=\tau N;
$$
if $\alpha$ is spacelike and $N$ is
timelike,
$$
N'=\kappa T-\tau B,\qquad B'=-\tau N.
$$
\end{proposition}

\begin{proof}
Expand $N'$ and $B'$ in the orthonormal
basis: the coefficient of a vector $E$ of
the frame in $X$ is
$\varepsilon_E\left\langle{} X,E
\right\rangle{}$. Differentiating the
relations
$\left\langle{} N,N \right\rangle{}
=\varepsilon_N$,
$\left\langle{} N,T \right\rangle{}=0$,
$\left\langle{} B,B \right\rangle{}
=\varepsilon_B$ and
$\left\langle{} B,T \right\rangle{}=0$
gives
$\left\langle{} N',N \right\rangle{}=0$,
$\left\langle{} N',T \right\rangle{}
=-\kappa\varepsilon_N$,
$\left\langle{} B',B \right\rangle{}=0$ and
$\left\langle{} B',T \right\rangle{}
=-\kappa\left\langle{} B,N
\right\rangle{}=0$. Together with
$\left\langle{} N',B \right\rangle{}
=-\tau$ and
$\left\langle{} B',N \right\rangle{}=\tau$
this gives the formulas. The explicit cases
use $\varepsilon_B=-\varepsilon_T
\varepsilon_N$.
\end{proof}

López \cite[Section
2.2]{lopez2014differential} uses $B=T\times
N$, which need not be positively oriented,
and other sign conventions for the torsion.
His curvature is the same, and in all three
cases above his torsion is $-\tau$; his
Frenet equations agree with the ones above
after this change.

\begin{example}
\label{example-palmas-pseudo-null-curve}
The curve $\alpha(s)=(s^2,s^2,s)$ has
$\alpha'(s)=(2s,2s,1)$ with
$\left\langle{} \alpha',\alpha'
\right\rangle{}=1$, so it is spacelike and
parametrized by arc length, while
$\alpha''(s)=(2,2,0)$ is null and
independent of $\alpha'$. For curves of
this kind we put $T=\alpha'$; there is no
curvature, since $|T'|=0$, and we only set
$N=T'$. For the binormal we take the
unique null vector $B$ orthogonal to $T$
with
$$
\left\langle{} N,B \right\rangle{}=-1.
$$
Here
$B(s)=(s^2+\tfrac14,\,s^2-\tfrac14,\,s)$.
\end{example}

For such curves (``pseudo-null'' curves),
define again
$\tau=-\left\langle{} N',B \right\rangle{}$.
The same method as in the proof of
Proposition~\ref{%
proposition-palmas-frenet-equations}
gives
$$
T'=N,\qquad N'=\tau N,\qquad
B'=T-\tau B,
$$
and in the example $\tau=0$ and $B'=T$.
For a null curve parametrized by pseudo-arc
length one takes $T=\alpha'$, $N=T'$ and
$B$ the unique null vector orthogonal to $N$
with
$\left\langle{} T,B \right\rangle{}=-1$;
with
$\tau=-\left\langle{} N',B \right\rangle{}$
one gets
$$
T'=N,\qquad N'=\tau T+B,\qquad B'=\tau N.
$$
For the helix
$\alpha(s)=(s,\cos s,\sin s)$ this gives
$B=\tfrac12(1,\sin s,-\cos s)$ and
$\tau=-\tfrac12$. These two systems were
derived here; they agree with
\cite[Section 2.2]{lopez2014differential}.

\begin{remark}
\label{remark-palmas-no-uniqueness}
Recall that in Euclidean $\mathbb{R}^3$,
given smooth functions $\kappa>0$ and
$\tau$, there is a curve parametrized by
arc length with curvature $\kappa$ and
torsion $\tau$, unique up to rigid
motions. In $\mathbb{R}^3_1$ there are
several types of Frenet equations, so
curvature and torsion alone do not
determine the curve up to rigid motions.
For instance, $(0,\cos s,\sin s)$,
$(\sinh s,0,\cosh s)$ and
$(\cosh s,0,\sinh s)$ all have
$\kappa=1$ and $\tau=0$, but the first is
spacelike with spacelike normal, the
second is timelike, and the third is
spacelike with timelike normal
\cite[Section 2.2]{lopez2014differential}.
\end{remark}

\begin{theorem}[Existence and uniqueness
for non-null curves]
\label{theorem-palmas-fundamental-curves}
Let $\kappa,\tau\colon I\to\mathbb{R}$ be
differentiable with $\kappa>0$, let
$s_0\in I$, $p_0\in\mathbb{R}^3_1$, and
let $\{T_0,N_0,B_0\}$ be a positively
oriented orthonormal basis of
$\mathbb{R}^3_1$. Then there is a unique
curve $\alpha\colon I\to\mathbb{R}^3_1$
such that
\begin{enumerate}
\item $\alpha(s_0)=p_0$;
\item $\alpha$ is parametrized by arc
length and $\alpha''$ is nonzero and not
null;
\item $\kappa$ and $\tau$ are the
curvature and the torsion of $\alpha$;
\item the Frenet--Serret frame of $\alpha$
at $s_0$ is $\{T_0,N_0,B_0\}$.
\end{enumerate}
\end{theorem}

The causal characters of $T$, $N$, $B$
along $\alpha$ are those of $T_0$, $N_0$,
$B_0$. The proof is as in the Euclidean
case: one solves the linear system of
Proposition~\ref{%
proposition-palmas-frenet-equations}
with the given initial frame and integrates
$T$; compare \cite[Theorems 2.6 and
2.8]{lopez2014differential}. In the notes
taken in class the initial point was
written $s\in I$ instead of $s_0\in I$.

\medskip\noindent
{\bf Surfaces in $\mathbb{R}^3_1$.}

\begin{definition}[Regular surface]
\label{definition-palmas-regular-surface}
A \emph{regular surface} in
$\mathbb{R}^3_1$ is a subset
$S\subset\mathbb{R}^3$ such that every
point of $S$ has a neighbourhood $V$ in
$S$ which is the image of a differentiable
map $X\colon U\to\mathbb{R}^3$, $U$ open in
$\mathbb{R}^2$, that is a homeomorphism
onto $V$ and whose differential is
injective at every point.
\end{definition}

This definition does not use the scalar
product; the notes taken in class stop
after ``is'', and we completed it with the
usual definition.

\begin{definition}[Causal character of a
surface]
\label{definition-palmas-causal-surface}
Let $S$ be a regular surface in
$\mathbb{R}^3_1$.
\begin{enumerate}
\item $S$ is \emph{spacelike} if $T_pS$ is
spacelike for every $p\in S$.
\item $S$ is \emph{timelike} if $T_pS$ is
timelike for every $p\in S$.
\item $S$ is \emph{nondegenerate} if it is
spacelike or timelike.
\item $S$ is \emph{degenerate} or
\emph{lightlike} if $T_pS$ is lightlike
for every $p\in S$.
\end{enumerate}
\end{definition}

\begin{proposition}
\label{proposition-palmas-compact-surface}
A compact regular surface in
$\mathbb{R}^3_1$ has spacelike, timelike
and lightlike tangent planes.
\end{proposition}

\begin{proof}
Each connected component of $S$ is
compact, so we may assume $S$ connected.
At a maximum $p$ of the coordinate $t$ on
$S$, the tangent plane $T_pS$ is the plane
$t=0$, which is spacelike. At a maximum
$q$ of the coordinate $x$ on $S$, the
tangent plane is the plane $x=0$, which
contains $(1,0,0)$ and is timelike. The
sets of points with spacelike and with
timelike tangent planes are open and
nonempty, so by connectedness they cannot
cover $S$.
\end{proof}

The notes taken in class state this as a
remark without proof; the proof was
supplied here. In particular there are no
compact spacelike, timelike or lightlike
surfaces without boundary
\cite[Section 3.1]{lopez2014differential}.

\begin{example}[De Sitter space of
dimension $2$]
\label{example-palmas-de-sitter}
The set
$$
\mathbb{S}^2_1=\{(t,x,y)\in\mathbb{R}^3_1
: -t^2+x^2+y^2=1\}
$$
is the ``unit sphere'' of
$\mathbb{R}^3_1$ (a pseudo-sphere). At
$p\in\mathbb{S}^2_1$ the position vector
$p$ is a spacelike unit normal, so
$T_p\mathbb{S}^2_1=p^{\perp}$ is timelike:
$\mathbb{S}^2_1$ is a timelike surface.
\end{example}

\begin{example}[The hyperbolic plane]
\label{example-palmas-hyperbolic-plane}
The set
$$
\mathbb{H}^2_0=\{(t,x,y)\in\mathbb{R}^3_1
: -t^2+x^2+y^2=-1,\ t>0\}
$$
is a spacelike surface: its normal $p$ is
timelike, so $p^{\perp}$ is spacelike. With
the induced metric it is the hyperbolic
plane. The subscript $0$ records that the
induced metric has index $0$.
\end{example}

\begin{definition}[Surface of revolution]
\label{definition-palmas-revolution}
Let $\ell$ be a line in $\mathbb{R}^3_1$,
$C$ a curve contained in a plane that
contains $\ell$, and $G$ the group of
orientation preserving isometries of
$\mathbb{R}^3_1$ that fix $\ell$ pointwise
(and, when $\ell$ is spacelike, preserve
the time orientation). The \emph{surface
of revolution} generated by $C$ around
$\ell$ is the orbit of $C$ under $G$.
\end{definition}

The notes taken in class stop at ``the
orbit of a curve $C$ is \dots''; the
definition was completed as above,
following the three cases discussed next
(compare the boosts with axis $\ell$ in
\cite[Section 1.3]{lopez2014differential}).
The extra condition when $\ell$ is
spacelike selects the connected group of
boosts and was added here.

\begin{remark}
\label{remark-palmas-rotation-groups}
Take $\ell$ through the origin.
\begin{enumerate}
\item If $\ell$ is timelike, take an
orthonormal basis $\{e_0,e_1,e_2\}$ with
$\ell=\mathbb{R}e_0$. An isometry fixing
$\ell$ preserves the spacelike plane
$\ell^{\perp}$ and acts on it by an
isometry of determinant $1$, that is, a
rotation:
$$
\begin{pmatrix}
1 & 0 & 0\\
0 & \cos\theta & -\sin\theta\\
0 & \sin\theta & \cos\theta
\end{pmatrix}.
$$
Here the other two basis vectors are
spacelike.
\item If $\ell=\mathbb{R}e_2$ is spacelike,
the plane $\ell^{\perp}$ is timelike and
the isometry acts on it by a boost, so the
matrix is
$$
\begin{pmatrix}
\cosh\theta & \sinh\theta & 0\\
\sinh\theta & \cosh\theta & 0\\
0 & 0 & 1
\end{pmatrix}.
$$
\item If $\ell$ is null, take a
pseudo-orthonormal basis
$\{v_1,v_2,v_3\}$ with $\ell=\mathbb{R}v_1$.
A computation with the scalar products
shows that the matrix is
$$
\begin{pmatrix}
1 & -\theta^2/2 & \theta\\
0 & 1 & 0\\
0 & -\theta & 1
\end{pmatrix},
$$
and in the lecture the generating curve
was taken in the plane spanned by $v_1$
and $v_2$.
\end{enumerate}
\end{remark}

In the notes taken in class, the first case
says that the $2\times 2$ block is ``a
rotation (or do we assume so, to fix
ideas?)''; it is necessarily a rotation, by
the argument above. The matrix of the third
case was derived here and replaces an
incomplete matrix in the notes.

Spoiler: spacelike surfaces are analogous
to Euclidean surfaces.

\begin{proposition}
\label{proposition-palmas-orientable}
Every spacelike surface in
$\mathbb{R}^3_1$ is orientable.
\end{proposition}

\begin{proof}
At each $p\in S$ the line $(T_pS)^{\perp}$
is timelike, so it contains exactly two
unit vectors, one in each time cone. Let
$N(p)$ be the one in the future cone, that
is, with
$\left\langle{} N(p),e_0 \right\rangle{}<0$.
Locally $N$ agrees, up to sign, with the
normalized cross product $X_u\times X_v$ of
a parametrization, and the sign is constant
on connected parameter domains, so $N$ is a
continuous unit normal field.
\end{proof}

This is \cite[Theorem
3.1]{lopez2014differential}. The notes
taken in class assume $S$ connected, which
is not needed.

\begin{definition}[Gauss map]
\label{definition-palmas-gauss-map}
The map $N\colon S\to\mathbb{H}^2_0$
constructed in the proof of
Proposition~\ref{%
proposition-palmas-orientable}
is the \emph{Gauss map} of the spacelike
surface $S$.
\end{definition}

\begin{lemma}
\label{lemma-palmas-dn-self-adjoint}
Let $S$ be a nondegenerate surface with a
unit normal field $N$. For each $p\in S$,
$dN_p$ maps $T_pS$ to $T_pS$ and is
self-adjoint.
\end{lemma}

\begin{proof}
Since
$\left\langle{} N,N \right\rangle{}$ is
constant,
$\left\langle{} dN_p(v),N(p)
\right\rangle{}=0$, so $dN_p(v)\in T_pS$.
For a parametrization $X$, differentiating
$\left\langle{} N,X_u \right\rangle{}=0$
with respect to $v$ gives
$\left\langle{} N_v,X_u \right\rangle{}
=-\left\langle{} N,X_{uv} \right\rangle{}$,
which is symmetric in $u$ and $v$.
\end{proof}

\begin{corollary}
\label{corollary-palmas-dn-diagonalizable}
If $S$ is spacelike, then $dN_p$ is
diagonalizable, with an orthonormal basis
of eigenvectors.
\end{corollary}

\begin{proof}
$T_pS$ is a Euclidean plane, so this is the
spectral theorem.
\end{proof}

\begin{definition}[Shape operator]
\label{definition-palmas-shape-operator}
For an orientable nondegenerate surface
with unit normal $N$, the \emph{shape
operator} at $p$ is $A=-dN_p$. It is a
self-adjoint operator on $T_pS$. When $A$
is diagonalizable its eigenvalues
$\kappa_1,\kappa_2$ are the
\emph{principal curvatures} at $p$.
\end{definition}

In the notes taken in class the shape
operator is used in the next definition
(with a question mark) and defined only
afterwards, for timelike surfaces; we have
moved the definition here.

\begin{definition}[Isoparametric surface]
\label{definition-palmas-isoparametric}
Let $S$ be a spacelike surface, $N\colon
S\to\mathbb{H}^2_0$ a unit normal field
and $A$ the shape operator. The surface $S$
is \emph{isoparametric} if its principal
curvatures $\kappa_1$ and $\kappa_2$ are
constant on $S$.
\end{definition}

Now we pass to timelike surfaces. For an
orientable timelike surface, $A=-dN_p$ is
again self-adjoint, now with respect to the
Lorentzian scalar product of $T_pS$. Since
$\dim T_pS=2$, by
Proposition~\ref{%
proposition-palmas-canonical-plane}
the matrix of $A$ has one of the forms
$$
\begin{pmatrix}
a_0 & 0\\
0 & a_1
\end{pmatrix},
\quad
\begin{pmatrix}
a_0 & 0\\
\delta & a_0
\end{pmatrix},
\quad
\begin{pmatrix}
a_0 & b_0\\
-b_0 & a_0
\end{pmatrix},
$$
of types I, II and IV respectively, with
$\delta=\pm 1$ and $b_0\neq 0$. Replacing
$N$ by $-N$ replaces $A$ by $-A$, which can
be used to make $\delta=1$. The notes taken
in class give type II with $1$ in place of
$\delta$ and do not specify the bases.

\begin{theorem}[Magid, 1985]
\label{theorem-palmas-magid}
There are no isoparametric timelike
surfaces in $\mathbb{R}^3_1$ whose shape
operators are of type IV.
\end{theorem}

Here isoparametric is meant in Magid's
sense: the minimal polynomial of the shape
operator is constant
\cite{magid1985lorentzian}; for type IV
this means that $a_0$ and $b_0$ are
constant. We have not checked the
statement in Magid's paper (not checked),
but in dimension two it follows from
López's theorem below: then $H$ and $K$ are
constant with $H^2-K=-b_0^2<0$, which is
excluded in the proof of \cite[Theorem
3.3]{lopez2014differential}.

\begin{definition}[Second fundamental
form]
\label{definition-palmas-second-form}
The \emph{second fundamental form} of $S$
at $p$ is the bilinear form
$$
\begin{aligned}
h_p\colon T_pS\times T_pS
&\longrightarrow\mathbb{R},\\
(u,v)&\longmapsto
\left\langle{} Au,v \right\rangle{}.
\end{aligned}
$$
\end{definition}

\begin{definition}[Umbilic point]
\label{definition-palmas-umbilic}
A point $p$ of a nondegenerate surface $S$
is \emph{umbilic} if there is a number
$\kappa(p)\in\mathbb{R}$ such that
$$
h_p(u,v)=\kappa(p)
\left\langle{} u,v \right\rangle{}
$$
for all $u,v\in T_pS$, that is,
$A=\kappa(p)\operatorname{id}$. The surface
$S$ is \emph{totally umbilic} if every
point of $S$ is umbilic.
\end{definition}

In the notes taken in class the left-hand
side is
$\left\langle{} \operatorname{II}(u,v),N(p)
\right\rangle{}$, where
$\operatorname{II}$ is the vector valued
second fundamental form; this is
$h_p(u,v)$.

\begin{proposition}
\label{proposition-palmas-totally-umbilic}
A connected totally umbilic nondegenerate
surface in $\mathbb{R}^3_1$ is, up to rigid
motions, an open subset of a plane, of a
hyperbolic plane
$$
\{p : \left\langle{} p,p
\right\rangle{}=-r^2,\ t>0\}
$$
(if $S$ is spacelike), or of a
pseudo-sphere
$$
\{p : \left\langle{} p,p
\right\rangle{}=r^2\}
$$
(if $S$ is timelike).
\end{proposition}

\begin{proof}
In a parametrization $X$ we have
$N_u=-\kappa X_u$ and $N_v=-\kappa X_v$.
Comparing $N_{uv}$ with $N_{vu}$ gives
$\kappa_vX_u=\kappa_uX_v$, so $\kappa$ is
locally constant, hence constant. If
$\kappa=0$, then $N$ is constant and $S$ is
in a plane. If $\kappa\neq 0$, then
$X+\kappa^{-1}N$ has zero derivatives, so
it is a constant point $c$, and
$$
\left\langle{} X-c,X-c \right\rangle{}
=\kappa^{-2}\left\langle{} N,N
\right\rangle{}.
$$
This is $-\kappa^{-2}$ if $S$ is spacelike
and $\kappa^{-2}$ if $S$ is timelike.
\end{proof}

The notes taken in class list only the
plane and the hyperbolic plane, which is
the spacelike case; the pseudo-spheres were
added (compare \cite[Theorem
3.2]{lopez2014differential}). The idea
given in class was that all points are at
the same distance from a fixed point, so
the surface is a ``sphere''; the proof above
makes this precise.

\begin{definition}[Mean curvature]
\label{definition-palmas-mean-curvature}
Let $S$ be a nondegenerate surface with
unit normal $N$ and
$\varepsilon=\left\langle{} N,N
\right\rangle{}$, so $\varepsilon=-1$ if
$S$ is spacelike and $\varepsilon=1$ if
$S$ is timelike. The \emph{mean curvature}
of $S$ at $p$ is
$$
H(p)=\frac{\varepsilon}{2}
\operatorname{tr}A.
$$
\end{definition}

\begin{definition}[Gaussian curvature]
\label{definition-palmas-gauss-curvature}
The \emph{Gaussian curvature} of $S$ at $p$
is
$$
K(p)=\varepsilon\det A.
$$
\end{definition}

The factor $\varepsilon$ in $H$ makes the
mean curvature vector
$\tfrac12\operatorname{tr}\operatorname{II}$
equal to $HN$, and with the factor
$\varepsilon$ in $K$ the Gauss equation
says that $K$ is the intrinsic curvature of
the induced metric \cite[Section
3.2]{lopez2014differential}. For example,
for the hyperbolic plane of radius $r$ with
the future normal $N=p/r$ we get $A=-r^{-1}
\operatorname{id}$, $H=1/r$ and
$K=-1/r^2$; for the pseudo-sphere of radius
$r$ with $N=p/r$ we get $H=-1/r$ and
$K=1/r^2$.

\begin{remark}
\label{remark-palmas-principal-curvatures}
If $A$ is diagonalizable, then
$$
H=\frac{\varepsilon}{2}(\kappa_1+\kappa_2),
\qquad
K=\varepsilon\kappa_1\kappa_2.
$$
\end{remark}

\begin{proposition}
\label{proposition-palmas-h2-minus-k}
Let $S$ be a nondegenerate surface in
$\mathbb{R}^3_1$ and $p\in S$ a point where
the shape operator $A$ is diagonalizable.
Then
$$
H(p)^2-\varepsilon K(p)\geq 0,
$$
with equality if and only if $p$ is
umbilic.
\end{proposition}

\begin{proof}
Since $\varepsilon^2=1$,
$$
H^2-\varepsilon K
=\frac{(\kappa_1+\kappa_2)^2}{4}
-\kappa_1\kappa_2
=\frac{(\kappa_1-\kappa_2)^2}{4},
$$
which vanishes exactly when
$\kappa_1=\kappa_2$, that is, when the
diagonalizable $A$ is a multiple of the
identity.
\end{proof}

For a timelike surface the hypothesis is
needed: at a point of type II,
$H^2-K=0$ although the point is not
umbilic, and at a point of type IV,
$H^2-K=-b_0^2<0$, as one computes from the
matrices above (compare
\cite[Proposition 3.4]{%
lopez2014differential}).

\begin{theorem}[López, 2014]
\label{theorem-palmas-lopez-constant-hk}
A connected nondegenerate surface in
$\mathbb{R}^3_1$ with constant mean
curvature and constant Gaussian curvature
is an open subset of a totally umbilic
surface, of one of the circular cylinders
$\mathbb{S}^1(r)\times\mathbb{R}$,
$\mathbb{S}^1_1(r)\times\mathbb{R}$,
$\mathbb{H}^1(r)\times\mathbb{R}$, or of a
timelike ruled
surface whose directrix and rulings are
lightlike; in the last case $H^2-K=0$.
\end{theorem}

This is \cite[Theorem
3.3]{lopez2014differential}. The notes
taken in class say only ``a ruled
surface''; the condition on the rulings
was added from López's statement.

\begin{definition}[Time orientation]
\label{definition-palmas-time-orientation}
A \emph{time orientation} on a timelike
surface $S$ is a continuous (smooth) choice
of a timelike tangent vector at each point
of $S$; it makes sense because each
tangent plane carries a Lorentzian scalar
product.
\end{definition}

\medskip\noindent
{\bf Lorentzian manifolds.}

The notes taken in class have a heading
for Lorentzian manifolds but no content;
they stop there.

\medskip\noindent
{\bf Exercises.}

\begin{exercise}
\label{exercise-palmas-orthogonal-null}
Show that in $\mathbb{R}^{n+1}_1$ two
orthogonal null vectors are linearly
dependent.
\end{exercise}

\begin{proof}[Solution]
Write $v=(v_0,\hat v)$ and $w=(w_0,\hat w)$
with $\hat v,\hat w\in\mathbb{R}^n$, and
let $\cdot$ be the Euclidean dot product
of $\mathbb{R}^n$. The hypotheses say
$$
\hat v\cdot\hat v=v_0^2,\qquad
\hat w\cdot\hat w=w_0^2,\qquad
\hat v\cdot\hat w=v_0w_0.
$$
Hence
$$
(\hat v\cdot\hat w)^2
=(\hat v\cdot\hat v)(\hat w\cdot\hat w),
$$
which is the equality case of the
Euclidean Cauchy--Schwarz inequality. Since
$v\neq 0$ is null, $\hat v\neq 0$, so
$\hat w=\lambda\hat v$ for some
$\lambda\in\mathbb{R}$. Then
$v_0w_0=\lambda v_0^2$ with $v_0\neq 0$,
so $w_0=\lambda v_0$ and $w=\lambda v$.
\end{proof}

\begin{exercise}
\label{exercise-palmas-light-cones}
\begin{enumerate}
\item What do the light cones of
$\mathbb{R}^3_1$, $\mathbb{R}^4_1$ and
$\mathbb{R}^4_2$ look like?
\item Are they connected?
\item What do the unit spheres of these
spaces look like?
\end{enumerate}
\end{exercise}

\begin{proof}[Solution]
In class we treated the light cone of
$\mathbb{R}^4_2$, written as the set of
nonzero $(a,b,c,d)$ with
$$
a^2+b^2-c^2-d^2=0 .
$$
Consider the map
$$
\begin{aligned}
f\colon (0,\infty)\times S^1\times S^1
&\longrightarrow\mathbb{R}^4,\\
(r,z_1,z_2)&\longmapsto(a,b,c,d),
\end{aligned}
$$
where $a+ib=rz_1$ and $c+id=rz_2$. Then
$a^2+b^2=r^2=c^2+d^2$, so the image lies in
the light cone, and every point of the
light cone is in the image, with
$r=\sqrt{a^2+b^2}>0$. Since the image of a
connected set under a continuous map is
connected, the light cone of
$\mathbb{R}^4_2$ is connected; in fact $f$
is a homeomorphism onto it. Similarly, the
light cone of $\mathbb{R}^3_1$ is
homeomorphic to two copies of
$(0,\infty)\times S^1$ and that of
$\mathbb{R}^4_1$ to two copies of
$(0,\infty)\times S^2$, one for each sign
of the time coordinate; they are not
connected. The unit spheres
$\left\langle{} v,v \right\rangle{}=1$ are
connected: in $\mathbb{R}^3_1$ it is a
hyperboloid of one sheet, homeomorphic to
$\mathbb{R}\times S^1$; in
$\mathbb{R}^4_1$ it is homeomorphic to
$\mathbb{R}\times S^2$; and in
$\mathbb{R}^4_2$, with the convention above,
it is homeomorphic to
$S^1\times\mathbb{R}^2$.
\end{proof}

In the notes taken in class $r$ ranged
over $[0,\infty)$, which also covers the
origin (not a null vector), and the image
was said to satisfy ``the equation of the
hyperboloid''. The answers for
$\mathbb{R}^3_1$, $\mathbb{R}^4_1$ and the
unit spheres were added here.

\begin{exercise}[Homework]
\label{exercise-palmas-cauchy-schwarz}
\begin{enumerate}
\item Give an example of a degenerate
symmetric bilinear form.
\item Find a proof of the classical
Cauchy--Schwarz inequality.
\item Can this proof be adapted to the
Lorentzian setting?
\end{enumerate}
\end{exercise}

\section{Introducción a la geometría
diferencial de superficies}
\label{section-garcia-surfaces}

\medskip\noindent
{\bf Curso Introductorio 2.}

Sandra García,
Universidad Nacional de Colombia --
Colombia.

\medskip\noindent
{\bf Schedule.}

Week 1, Monday to Thursday, 14:30--16:00
(8--11 July 2024).

\medskip\noindent
{\bf Official listing.}

\href{%
https://naturales.ues.edu.sv/%
ecimpasv2024/introductory-curse-2/}{%
Programme entry.}

\medskip\noindent
{\bf Abstract.}

El objetivo de este curso es presentar
algunas nociones básicas sobre la geometría
de curvas y superficies en el espacio
Euclídeo. Iniciaremos con un tratamiento de
curvas regulares en R3, centrándose en el
estudio de sus propiedades locales que
aparecen de forma natural al estudiar las
superficies. En concreto, introduciremos los
conceptos esenciales de curvas
parametrizadas, longitud de arco, curvatura,
torsión y el triedro de Frenet.

Comenzaremos el estudio de superficies
desarrollando el concepto de superficie
regular en R3 y algunas nociones
fundamentales como el plano tangente y la
primera forma fundamental de una superficie.
Posteriormente, introduciremos uno de los
conceptos más importantes para el estudio de
la geometría extrínseca de superficies: la
aplicación de Gauss. La aplicación de Gauss
nos permitirá definir la segunda forma
fundamental, que será esencial para
introducir dos de las formas más comunes de
medir la curvatura de una superficie, la
media y la curvatura de Gauss. Demostraremos
que la curvatura de Gauss es, de hecho, un
concepto intrínseco, ya que es invariante por
isometrías locales, mientras que la curvatura
media es extrínseca.

Se introducirán algunas familias importantes
de superficies regulares, como superficies
regladas y mínimas. A lo largo del curso se
presentarán muchos ejemplos de curvas y
superficies conocidas.

Este curso representa la base para una
introducción al ``Análisis Geométrico en
Superficies'', que se estudiará en un curso
diferente.

\medskip\noindent
{\bf Sources of these notes.}
Daniel's notes taken during the course
record only four exercises proposed in the
first session, and nothing of the lectures
themselves. The exercises are translated
below. The notes contain no solutions, so
the short solutions given here were written
afterwards for these notes; they are
computations, checked by hand, and those
for the curves also numerically. For
background on curves and surfaces see
\cite{docsu}. Minimal
immersions in general are defined in
Definition~%
\ref{differential-geometry-definition-%
minimal-immersion}.

\medskip\noindent
{\bf Conventions.}
For a regular curve $\alpha$ in
$\mathbb{R}^3$ with $\alpha' \times \alpha''$
nowhere zero, the curvature and the torsion
are
$$
\kappa =
\frac{|\alpha' \times \alpha''|}
{|\alpha'|^3},
\qquad
\tau =
\frac{\left\langle{} \alpha' \times \alpha'',
\alpha''' \right\rangle{}}
{|\alpha' \times \alpha''|^2}.
$$
The Frenet frame is $T = \alpha'/|\alpha'|$,
$N = B \times T$ and
$B = (\alpha' \times \alpha'') /
|\alpha' \times \alpha''|$.
With this sign of $\tau$ the Frenet
equations in arc length read $T' = \kappa N$,
$N' = -\kappa T + \tau B$ and
$B' = -\tau N$; some books use the
opposite sign for the torsion.

\begin{exercise}
\label{exercise-garcia-catenary}
Find the curvature of the catenary.
\end{exercise}

\begin{proof}[Solution]
Take the catenary as the graph of
$y = a\cosh(x/a)$ with $a > 0$, that is
$\alpha(t) = (t, a\cosh(t/a))$. For a plane
graph $y = f(x)$ the curvature is
$$
\kappa = \frac{|f''|}{(1 + f'^2)^{3/2}}.
$$
Here $f' = \sinh(t/a)$ and
$f'' = \cosh(t/a)/a$, and
$1 + \sinh^2 = \cosh^2$, so
$$
\kappa(t) =
\frac{\cosh(t/a)}{a \cosh^3(t/a)} =
\frac{1}{a \cosh^2(t/a)}.
$$
The curvature is largest, equal to $1/a$, at
the vertex $t = 0$, and tends to $0$ as
$|t| \to \infty$.
\end{proof}

\begin{exercise}
\label{exercise-garcia-helix}
Parametrize the helix and find its
curvature.
\end{exercise}

\begin{proof}[Solution]
The circular helix of radius $a > 0$ and
pitch $2\pi b$, with $b \neq 0$, is
$$
\alpha(t) = (a\cos t, a\sin t, bt),
\qquad t \in \mathbb{R}.
$$
Then
$$
\begin{aligned}
\alpha'(t) &= (-a\sin t, a\cos t, b), \\
\alpha''(t) &= (-a\cos t, -a\sin t, 0), \\
\alpha'''(t) &= (a\sin t, -a\cos t, 0).
\end{aligned}
$$
So $|\alpha'| = \sqrt{a^2 + b^2}$ is
constant and the arc length is
$s = \sqrt{a^2 + b^2}\, t$. Moreover
$$
\alpha' \times \alpha'' =
(ab\sin t, -ab\cos t, a^2),
$$
whose norm is $a\sqrt{a^2 + b^2}$. Hence
$$
\kappa = \frac{a}{a^2 + b^2}.
$$
Also
$\left\langle{} \alpha' \times \alpha'',
\alpha''' \right\rangle{} = a^2 b$, so the
torsion is the constant
$$
\tau = \frac{b}{a^2 + b^2}.
$$
For $b \to 0$ we recover the circle of
radius $a$, with $\kappa = 1/a$ and
$\tau = 0$.
\end{proof}

\begin{exercise}
\label{exercise-garcia-frenet-frame}
Find the Frenet frame of the curve
$$
\alpha(t) =
\left(
\frac{(1+t)^{3/2}}{3},
\frac{(1-t)^{3/2}}{3},
\frac{t}{\sqrt{2}}
\right),
\qquad t \in (-1,1),
$$
at the point $(1/3, 1/3, 0)$.
\end{exercise}

\begin{proof}[Solution]
The point is $\alpha(0)$. We have
$$
\begin{aligned}
\alpha'(t) &=
\left(
\frac{(1+t)^{1/2}}{2},
-\frac{(1-t)^{1/2}}{2},
\frac{1}{\sqrt{2}}
\right), \\
\alpha''(t) &=
\left(
\frac{1}{4(1+t)^{1/2}},
\frac{1}{4(1-t)^{1/2}},
0
\right).
\end{aligned}
$$
Since
$|\alpha'|^2 = (1+t)/4 + (1-t)/4 + 1/2 = 1$,
the curve is parametrized by arc length, so
$T = \alpha'$ and $\alpha'' = \kappa N$.
From
$$
|\alpha''|^2 =
\frac{1}{16}
\left(\frac{1}{1+t} + \frac{1}{1-t}\right)
= \frac{1}{8(1-t^2)}
$$
we get
$$
\kappa(t) = \frac{1}{2\sqrt{2}\sqrt{1-t^2}},
\qquad
\kappa(0) = \frac{\sqrt{2}}{4}.
$$
At $t = 0$,
$$
\begin{aligned}
T(0) &=
\left(\frac{1}{2}, -\frac{1}{2},
\frac{1}{\sqrt{2}}\right), \\
N(0) &= \frac{\alpha''(0)}{\kappa(0)} =
\left(\frac{1}{\sqrt{2}},
\frac{1}{\sqrt{2}}, 0\right), \\
B(0) &= T(0) \times N(0) =
\left(-\frac{1}{2}, \frac{1}{2},
\frac{1}{\sqrt{2}}\right).
\end{aligned}
$$
One checks directly that these three
vectors are orthonormal. For the torsion,
$\alpha'''(0) = (-1/8, 1/8, 0)$ and
$\alpha'(0) \times \alpha''(0) =
\kappa(0) B(0)$, so
$$
\tau(0) =
\frac{\left\langle{} \kappa(0) B(0),
\alpha'''(0) \right\rangle{}}
{\kappa(0)^2}
= \frac{\sqrt{2}}{4}.
$$
In fact the same computation for general
$t$ gives $\tau(t) = \kappa(t)$ for all
$t \in (-1,1)$.
\end{proof}

\begin{exercise}
\label{exercise-garcia-helicoid-minimal}
Show that the helicoid is a minimal
surface.
\end{exercise}

\begin{proof}[Solution]
Parametrize the helicoid, for a constant
$a \neq 0$, by
$$
X(u,v) = (v\cos u, v\sin u, au),
\qquad (u,v) \in \mathbb{R}^2.
$$
Then
$$
\begin{aligned}
X_u &= (-v\sin u, v\cos u, a), \\
X_v &= (\cos u, \sin u, 0),
\end{aligned}
$$
so the coefficients of the first
fundamental form are
$E = v^2 + a^2$, $F = 0$ and $G = 1$. Since
$$
X_u \times X_v =
(-a\sin u, a\cos u, -v)
$$
never vanishes, $X$ is regular, with unit
normal
$$
\nu = \frac{(-a\sin u, a\cos u, -v)}
{\sqrt{a^2 + v^2}}.
$$
The second derivatives are
$$
\begin{aligned}
X_{uu} &= (-v\cos u, -v\sin u, 0), \\
X_{uv} &= (-\sin u, \cos u, 0), \\
X_{vv} &= 0,
\end{aligned}
$$
so the coefficients of the second
fundamental form are
$$
e = \left\langle{} X_{uu}, \nu
\right\rangle{} = 0,
\qquad
f = \frac{a}{\sqrt{a^2 + v^2}},
\qquad
g = 0.
$$
Therefore the mean curvature vanishes:
$$
H = \frac{eG - 2fF + gE}{2(EG - F^2)}
= 0.
$$
For the record, the Gaussian curvature is
$$
K = \frac{eg - f^2}{EG - F^2}
= -\frac{a^2}{(a^2 + v^2)^2}.
$$
\end{proof}

\section{Operadores Diferenciales en
Variedades de Riemann}
\label{section-santos-lobos-operators}

\medskip\noindent
{\bf Curso Introductorio 3.}

Walcy Santos,
UFRJ -- Brasil \& Guillermo Lobos,
UFSCar -- Brasil.

\medskip\noindent
{\bf Schedule.}

Week 1, Tuesday to Friday, 09:00--10:30
(9--12 July 2024).

\medskip\noindent
{\bf Official listing.}

\href{%
https://naturales.ues.edu.sv/%
ecimpasv2024/introductory-curse-3/}{%
Programme entry.}

\medskip\noindent
{\bf Abstract.}

La extensión del análisis de operadores
diferenciales en dominios del espacio
euclidiano al contexto de variedades suaves
es un tema moderno y relevante que involucra
conceptos geométricos y analíticos, y ha
recibido mucha atención en la literatura
reciente. En este curso, estableceremos las
herramientas básicas del análisis geométrico
como primer contacto para estudiantes e
investigadores interesados en trabajar en
temas en esta área. Cubriremos los siguientes
temas: una revisión de la geometría
riemanniana: variedades suaves, métrica
riemanniana, conexión de Levi-Civita y
curvaturas; operadores diferenciales en
variedades riemannianas: gradiente,
divergente, Laplaciano y Hessiano; tensores y
los operadores diferenciales: identidad
contraída de Bianchi, fórmula de Bochner;
inmersiones isométricas.

\medskip\noindent
{\bf Sources of these notes.}
Daniel's notes taken during the course
record only four exercises, two of them
with attempted solutions, and nothing of
the lectures themselves. The course
material includes the lecture notes
\cite{gomes2015operadores} by José Nazareno
Vieira Gomes, written at UFAM in 2015
(the school's page lists Dr.\ José Gomes,
Universidade Federal de São Carlos, in
the teaching team), which cover
the topics of the abstract; we cite them
for the definitions used below. The
solutions of Exercises~%
\ref{exercise-santos-lobos-constant-len}
and~%
\ref{exercise-santos-lobos-distance-fn}
were corrected, and the other solutions
were written afterwards for these notes.
The statement of
Exercise~%
\ref{exercise-santos-lobos-distance-fn}
as recorded in class is false, and has been
corrected; see
Remark~%
\ref{remark-santos-lobos-distance-false}.
The Frobenius theorem used there is
Theorem~\ref{hodge-theory-theorem-frobenius}.

\medskip\noindent
{\bf Conventions.}
Throughout, $(M,g)$ is a Riemannian
manifold, we write
$\left\langle{} X,Y \right\rangle{}
= g(X,Y)$,
and $\nabla$ is the Levi-Civita connection,
the unique affine connection that is
torsion free and compatible with $g$
\cite[Section 1.3, Theorem 1.2]%
{gomes2015operadores}. The gradient
$\nabla f$ of a function $f$ is defined by
$\left\langle{} \nabla f, Y \right\rangle{}
= df(Y) = Y(f)$. Following
\cite[Section 2.1, Definition 2.10]%
{gomes2015operadores}, the Hessian of $f$
is the $(0,2)$-tensor
$$
\operatorname{Hess} f(X,Y) =
\left\langle{} \nabla_X \nabla f, Y
\right\rangle{},
$$
also written $\nabla^2 f$. The Lie
derivative of a $(0,2)$-tensor $T$ along a
vector field $X$ satisfies
\cite[Chapter 2, p.~22]{gomes2015operadores}
$$
(\mathcal{L}_X T)(Y,Z) =
X(T(Y,Z)) - T([X,Y],Z) - T(Y,[X,Z]).
$$

\begin{exercise}
\label{exercise-santos-lobos-constant-len}
Show that if $X$ is a vector field of
constant length on a Riemannian manifold,
then $\nabla_v X$ is orthogonal to $X$ for
every tangent vector $v$.
\end{exercise}

\begin{proof}[Solution]
The function
$\left\langle{} X,X \right\rangle{}$ is
constant, so its derivative in the
direction of any $v \in T_pM$ vanishes.
By compatibility of $\nabla$ with $g$,
$$
0 = v\left\langle{} X,X \right\rangle{}
= 2\left\langle{} \nabla_v X, X(p)
\right\rangle{}.
$$
Hence $\nabla_v X \perp X(p)$.
\end{proof}

\begin{exercise}
\label{exercise-santos-lobos-euclidean}
Show that the Riemannian connection of
Euclidean space $\mathbb{R}^n$ is the only
affine connection $D$ such that
$D_Y E = 0$ for every constant vector field
$E$ and every vector field $Y$.
\end{exercise}

\begin{proof}[Solution]
Let $e_1, \ldots, e_n$ be the standard
constant vector fields. Write a vector
field as $Y = \sum_j b^j e_j$ with smooth
functions $b^j$. If $D$ is an affine
connection with $D e_j = 0$ for all $j$,
then the Leibniz rule gives
$$
D_X Y = \sum_j X(b^j)\, e_j
+ \sum_j b^j D_X e_j
= \sum_j X(b^j)\, e_j.
$$
So $D_X Y$ is the ordinary directional
derivative of $Y$ in the direction $X$, and
$D$ is uniquely determined. On the other
hand, the Euclidean metric has constant
coefficients $g_{ij} = \delta_{ij}$, so its
Christoffel symbols vanish by the Koszul
formula, and its Levi-Civita connection is
also the directional derivative
\cite[Section 1.3, equations (1.4) and
(1.5)]{gomes2015operadores}. In particular
it kills the constant vector fields, which
proves existence, and by the uniqueness just
shown it is the only such connection.
\end{proof}

\begin{exercise}
\label{exercise-santos-lobos-hessian}
Let $f$ be a smooth function on $M$. Show
that
$$
\operatorname{Hess} f = \nabla df,
\qquad
\mathcal{L}_{\nabla f} g =
2 \operatorname{Hess} f.
$$
\end{exercise}

\begin{proof}[Solution]
Both identities are stated in
\cite[Section 2.1, (2.25)]%
{gomes2015operadores}. For the first, the
covariant derivative of the $1$-form $df$ is
$$
\begin{aligned}
(\nabla df)(X,Y)
&= (\nabla_X df)(Y) \\
&= X(df(Y)) - df(\nabla_X Y) \\
&= X(Y(f)) - (\nabla_X Y)(f).
\end{aligned}
$$
By compatibility with the metric,
$$
\begin{aligned}
\operatorname{Hess} f(X,Y)
&= X\left\langle{} \nabla f, Y
\right\rangle{}
- \left\langle{} \nabla f, \nabla_X Y
\right\rangle{} \\
&= X(Y(f)) - (\nabla_X Y)(f),
\end{aligned}
$$
which is the same. Since $\nabla$ is torsion
free,
$$
\operatorname{Hess} f(X,Y)
- \operatorname{Hess} f(Y,X)
= [X,Y](f) - (\nabla_X Y - \nabla_Y X)(f)
= 0,
$$
so the Hessian is symmetric.

For the second identity, first let $X$ be
any vector field. Using
$[X,Y] = \nabla_X Y - \nabla_Y X$ and
compatibility with the metric in the
formula for $\mathcal{L}_X g$,
$$
\begin{aligned}
(\mathcal{L}_X g)(Y,Z)
&= \left\langle{} \nabla_X Y, Z
\right\rangle{}
+ \left\langle{} Y, \nabla_X Z
\right\rangle{} \\
&\quad - \left\langle{} \nabla_X Y
- \nabla_Y X, Z \right\rangle{} \\
&\quad - \left\langle{} Y, \nabla_X Z
- \nabla_Z X \right\rangle{} \\
&= \left\langle{} \nabla_Y X, Z
\right\rangle{}
+ \left\langle{} Y, \nabla_Z X
\right\rangle{},
\end{aligned}
$$
as in \cite[(2.10)]{gomes2015operadores}.
For $X = \nabla f$ the right-hand side is
$\operatorname{Hess} f(Y,Z) +
\operatorname{Hess} f(Z,Y)$, which equals
$2\operatorname{Hess} f(Y,Z)$ by symmetry.
\end{proof}

\medskip\noindent
{\bf Distance functions.}
Let $U \subset M$ be open. A function
$r : U \to \mathbb{R}$ is called a
\emph{distance function} if
$|\nabla r| = 1$ on $U$. For a vector field
$X$ let $\omega_X$ be the $1$-form
$\omega_X(Y) = \left\langle{} X,Y
\right\rangle{}$, and let $X^\perp$ be the
distribution of vectors orthogonal to $X$.
The exercise recorded in class asked to
show that $X$ is locally the gradient of a
distance function if and only if $X^\perp$
is integrable, with the hint to use
$\omega_X$. As explained in
Remark~%
\ref{remark-santos-lobos-distance-false},
this is false as stated; the following is
a corrected version.

\begin{exercise}
\label{exercise-santos-lobos-distance-fn}
Let $X$ be a vector field with $|X| = 1$ on
$M$. Show that the following are equivalent.
\begin{enumerate}
\item Every point of $M$ has a neighbourhood
$U$ and a distance function
$r : U \to \mathbb{R}$ with $X = \nabla r$
on $U$.
\item The form $\omega_X$ is closed.
\item The distribution $X^\perp$ is
integrable and $\nabla_X X = 0$, that is,
the integral curves of $X$ are geodesics.
\end{enumerate}
\end{exercise}

\begin{proof}[Solution]
Write $\omega = \omega_X$. Recall that
$$
d\omega(Y,Z) = Y(\omega(Z)) - Z(\omega(Y))
- \omega([Y,Z]).
$$
Expanding with the compatibility of
$\nabla$ with $g$ and its torsion freeness,
the terms
$\left\langle{} X, \nabla_Y Z
\right\rangle{}$ and
$\left\langle{} X, \nabla_Z Y
\right\rangle{}$ cancel and we get
\begin{equation}
\label{equation-santos-lobos-domega}
d\omega(Y,Z) =
\left\langle{} \nabla_Y X, Z \right\rangle{}
- \left\langle{} \nabla_Z X, Y
\right\rangle{}.
\end{equation}
Since $|X| = 1$, Exercise~%
\ref{exercise-santos-lobos-constant-len}
gives $\nabla_Y X \perp X$ for every $Y$.

(1) implies (2): on $U$ we have
$\omega = dr$, which is closed.

(2) implies (1): by the Poincaré lemma,
on a coordinate ball $U$ around any point
there is $r$ with $dr = \omega$. Then
$\nabla r = X$ and $|\nabla r| = |X| = 1$,
so $r$ is a distance function.

(2) is equivalent to (3): since
$TM = \mathbb{R}X \oplus X^\perp$ and
$d\omega$ is alternating, $d\omega = 0$ if
and only if $d\omega(Y,Z) = 0$ and
$d\omega(X,Z) = 0$ for all sections $Y, Z$
of $X^\perp$. For the first condition,
$\omega(Y) = \omega(Z) = 0$, so the formula
for $d\omega$ gives
$$
d\omega(Y,Z) = -\left\langle{} X, [Y,Z]
\right\rangle{}.
$$
Thus $d\omega$ vanishes on $X^\perp$ if and
only if $X^\perp$ is involutive, which by
the Frobenius theorem means that it is
integrable. For the second condition,
\eqref{equation-santos-lobos-domega} and
$\nabla_Z X \perp X$ give
$$
d\omega(X,Z) =
\left\langle{} \nabla_X X, Z \right\rangle{}.
$$
As $\nabla_X X \perp X$, this vanishes for
all $Z \in X^\perp$ if and only if
$\nabla_X X = 0$.
\end{proof}

\noindent
Directly, if $X = \nabla r$ with
$|\nabla r| = 1$, then $X^\perp = \ker dr$
is tangent to the level sets of $r$, which
are hypersurfaces because $dr \neq 0$; this
is the argument of the notes. Moreover
$\nabla_X X = 0$ follows from the identity
$\frac{1}{2} d|\nabla r|^2 =
\operatorname{Hess} r(\nabla r, \cdot)$
\cite[(2.26)]{gomes2015operadores}, since
the left-hand side vanishes. The attempt in
the notes to prove the converse took $r$ to
be the distance to a leaf of $X^\perp$ and
needed $X(r) = |X|^2$, a step Daniel
doubted. The doubt was justified: the
converse needs the extra hypotheses
$|X| = 1$ and $\nabla_X X = 0$.

\begin{remark}
\label{remark-santos-lobos-distance-false}
The statement recorded in class, that $X$
is locally the gradient of a distance
function if and only if $X^\perp$ is
integrable, fails in two ways.
\begin{enumerate}
\item The length matters: on
$\mathbb{R}^n$ the field
$X = 2\,\partial/\partial x_1$ has
integrable $X^\perp$, whose leaves are the
hyperplanes $x_1 = c$, but $|X| = 2$, so it
is not the gradient of a distance function.
\item Even for $|X| = 1$ integrability is
not enough. On $\mathbb{R}^2 \setminus
\{0\}$ take
$$
X = \frac{1}{\sqrt{x^2 + y^2}}
\left(-y \frac{\partial}{\partial x}
+ x \frac{\partial}{\partial y}\right).
$$
Every line field on a surface is
integrable; here the leaves of $X^\perp$
are the open rays from the origin. But the
integral curves of $X$ are circles
traversed with unit speed, so
$\nabla_X X$ at a point $p$ is
$-p/|p|^2 \neq 0$. Equivalently, in polar
coordinates $\omega_X = \rho\, d\theta$ and
$d\omega_X = d\rho \wedge d\theta \neq 0$.
By Exercise~%
\ref{exercise-santos-lobos-distance-fn}
$X$ is not locally a gradient of a distance
function. Nor is $X$ locally parallel to
the gradient of a distance function $r$:
the level sets of such an $r$ would be the
rays, so $r$ would be a function of
$\theta$ alone, and then
$|\nabla r| = |r'(\theta)|/\rho$ is not
constant.
\end{enumerate}
\end{remark}

\begin{remark}
\label{remark-santos-lobos-frobenius-form}
What integrability of $X^\perp$ does give,
for a nowhere vanishing $X$ of any length,
is that locally $X = h \nabla f$ for smooth
functions $h > 0$ and $f$ with
$df \neq 0$. Indeed, if $X = h\nabla f$ then
$X^\perp = \ker df$ is integrable.
Conversely, by the Frobenius theorem there
is locally a submersion $f$ with
$\ker df = X^\perp = \ker \omega_X$, so
$\omega_X = h\, df$ for a nowhere vanishing
function $h$; replacing $f$ by $-f$ if
necessary, $h > 0$ and then
$X = h \nabla f$. In terms of forms,
$X^\perp$ is integrable if and only if
$\omega_X \wedge d\omega_X = 0$, while
Exercise~%
\ref{exercise-santos-lobos-distance-fn}
asks for the stronger condition
$d\omega_X = 0$.
\end{remark}

\section{Avances Recientes en Geometría de
Lorentz}
\label{section-solis-lorentzian-geometry}

\medskip\noindent
{\bf Curso Avanzado 1.}

Didier Solis,
Universidad Autónoma de Yucatán-Mexico.

\medskip\noindent
{\bf Schedule.}

Week 2, Monday to Thursday, 09:00--10:30
(15--18 July 2024).

\medskip\noindent
{\bf Official listing.}

\href{%
https://naturales.ues.edu.sv/%
ecimpasv2024/advanced-curse-1/}{%
Programme entry.}

\medskip\noindent
{\bf Abstract.}

Históricamente, el desarrollo de la geometría
de Lorentz ha estado íntimamente relacionado
con el de la teoría de la Reletividad
General. En este mini-curso exploraremos
algunas áreas activas de investigación en
geometría de Lorentz, abarcando desde temas
propios de la matemática abstracta hasta
aplicaciones a la Relatividad. Nos
enfocaremos en tres aspectos principales:
geometría nula, teoría causal y una novedosa
perspectiva sintética de la geometría de
Lorentz. En la primera sesión describiremos
la geometría de hipersuperficies nulas usando
en lo general el formalismo propuesto por K.
Duggal, que permite definir ecuaciones del
tipo Gauss-Weingarten con el propósito de
desarrollas conceptos geométricos clásicos
(como el de umbilicidad). Además discutiremos
algunos teoremas recientes de clasificación
en el caso de los espacios de
Robertson-Walker. En la segunda sesión
abordaremos los principios básicos de la
teoría causal, que es uno de los aspectos
fundamentales en Relatividad Matemática. La
tercera sesión está dedicada al estudio de
las subavariedades nulas desde un punto de
vista intrínseco, lo que permite interpretar
en términos geométricos algunas de las
propiedades más conocidas de la teoría de la
Relatividad, como es el caso de los agujeros
negros y la estructura asintótica del
espacio-tiempo. Finalmente, en la última
sesión exploraremos la noción reciente de
espacios Lorentzianos de longitud,
introducida apenas en 2018 por M. Kunzinger y
C. Sämann. Este camopo de investigación de
vertiginoso crecimiento nos permite
establecer conceptos geométricos
fundamentales (como el de cotas de curvatura)
en contextos donde no es posible garantizar
la regularidad de tipo $C^2$ del tensor de
curvatura, o incluso en escenarios donde no
existe ningún tensor métrico. En la sesión de
problemas se analizarán los célebres teoremas
de singularidad de S. Hawking y R. Penrose a
la luz de las herramientas desarrolladas a lo
largo del curso.

\medskip\noindent
{\bf Sources of these notes.}
These notes are an English version of
Daniel's notes taken during the course
(lectures and problem session), in two
almost identical versions; the later one
was used. The notes cover the first three
sessions and the problem session; the
fourth session, on Lorentzian length
spaces, is not recorded. Missing pieces of
formulas, wrong signs and wrong hypotheses
were completed or corrected using
\cite{oneill} (Chapters 13 and 14),
\cite{navarro2020screen} (Sections 2, 3
and 5), and short arguments written out
here. The classification results on null
hypersurfaces are taken from the lecture;
where their hypotheses could not be
checked against a source this is marked
``(not checked)''. The general framework
for null hypersurfaces is
\cite{duggal1996lightlike}; for causality
theory see also \cite{minguzzi2019causality}
and \cite{hawking1973large}.

Lorentzian curves and surfaces in
Minkowski space are the subject of Oscar
Palmas's course in this school,
Section~\ref{section-palmas-%
lorentzian-curves}.

\medskip\noindent
{\bf Session 1. Geometry of null
hypersurfaces.}
If we restrict a Riemannian metric to a
submanifold we obtain a Riemannian metric
on the submanifold. This fails in the
Lorentzian case: the restriction of a
Lorentzian metric to a submanifold may be
Lorentzian, Riemannian, degenerate, or may
change type from point to point.

\begin{example}
\label{example-solis-induced-metrics}
Let $\mathbb{R}^{n+2}_1$ be Minkowski
space with the scalar product
$$
\left\langle{} x,y \right\rangle{}
=-x_0y_0+x_1y_1+\cdots+x_{n+1}y_{n+1}.
$$
\begin{itemize}
\item On the pseudosphere
$\left\langle{} x,x \right\rangle{}=1$,
de Sitter space $\mathbb{S}^{n+1}_1$,
the restriction is again Lorentzian.
\item On the hyperboloid of two sheets
$\left\langle{} x,x \right\rangle{}=-1$
the restriction is Riemannian (each sheet
is a copy of hyperbolic space).
\item A hyperplane tangent to the light
cone contains no timelike vectors, and the
restriction is degenerate: it is
positive semidefinite with a
one-dimensional kernel.
\item For a ``horizontal'' cylinder, that
is, one whose axis is spacelike, such as
$x_0^2+x_1^2=1$ in $\mathbb{R}^3_1$, the
induced metric is Lorentzian on some
regions, Riemannian on others, and
degenerate along four lines.
\end{itemize}
\end{example}

\noindent
In what follows $(\widetilde{M},
\widetilde{g})$ is a Lorentzian manifold
of dimension $n+2$ with Levi-Civita
connection $\widetilde{\nabla}$, and $M$
is an immersed hypersurface, of dimension
$n+1$, with first fundamental form
$g=\widetilde{g}|_{TM}$.

\begin{definition}
\label{definition-solis-null-hypersurface}
The hypersurface $M$ is a \emph{null
hypersurface} (or lightlike hypersurface)
if its first fundamental form $g$ is
degenerate at every point.
\end{definition}

\noindent
The \emph{radical} of $M$ at $p$ is
$$
\operatorname{rad}(T_pM)
=T_pM\cap T_pM^{\perp}.
$$

\begin{lemma}
\label{lemma-solis-radical-rank-one}
If $M$ is a null hypersurface, then
$\operatorname{rad}(TM)=TM^{\perp}$. In
particular $\operatorname{rad}(TM)$ is a
smooth vector bundle of rank one spanned
by null vectors, and
$T_pM=\operatorname{rad}(T_pM)^{\perp}$.
\end{lemma}

\begin{proof}
Since $\widetilde{g}$ is non-degenerate
and $T_pM$ is a hyperplane of
$T_p\widetilde{M}$, the orthogonal
$T_pM^{\perp}$ is a line. The radical is a
subspace of this line, and it is non-zero
because $g$ is degenerate at $p$. Hence
$\operatorname{rad}(T_pM)=T_pM^{\perp}$,
and taking orthogonals gives
$T_pM=\operatorname{rad}(T_pM)^{\perp}$.
A non-zero $\xi$ in the radical lies in
$T_pM$ and is orthogonal to $T_pM$, so
$\widetilde{g}(\xi,\xi)=0$. Finally
$TM^{\perp}$ is a smooth line bundle, as
for any hypersurface.
\end{proof}

\noindent
In class the rank one property was
justified by the following exercise.

\begin{exercise}
\label{exercise-solis-null-orthogonal}
In a Lorentzian vector space, two null
vectors that are orthogonal to each other
are collinear.
\end{exercise}

\begin{proof}
Let $u,v$ be non-zero null vectors with
$\left\langle{} u,v \right\rangle{}=0$.
Choose a unit timelike vector $e_0$ and
write $u=ae_0+u'$ and $v=be_0+v'$ with
$u',v'$ orthogonal to $e_0$. The
orthogonal complement of $e_0$ is
positive definite; write $|\cdot|$ for its
norm and $\cdot$ for its scalar product.
The vectors being null and orthogonal
means
$$
a^2=|u'|^2,\qquad b^2=|v'|^2,\qquad
ab=u'\cdot v'.
$$
Hence
$|u'\cdot v'|=|a||b|=|u'||v'|$, which is
the equality case of the Cauchy--Schwarz
inequality. Since $u\neq 0$ we have
$u'\neq 0$, so $v'=\lambda u'$ for some
$\lambda\in\mathbb{R}$. Then
$ab=\lambda|u'|^2=\lambda a^2$, and
$a\neq 0$, so $b=\lambda a$. Therefore
$v=\lambda u$.
\end{proof}

\noindent
Consequently, at each point of a null
hypersurface the tangent space $T_pM$
contains exactly one null direction,
namely the radical. Unlike the Riemannian
case, the tangent bundle of the ambient
space does not split as the direct sum of
$TM$ and the normal bundle, since
$TM^{\perp}\subset TM$.

We write $\xi$ for a local nowhere
vanishing section of
$\operatorname{rad}(TM)$. Given a point of
$M$, the null vectors tangent to $M$ are
the multiples of $\xi$; by
Exercise~\ref{exercise-solis-null-orthogonal}
any other null vector at that point is
\emph{not tangent} to $M$. Choosing such a
null vector at each point smoothly we get
a null line bundle
$\operatorname{tr}(TM)$, the
\emph{transversal bundle}, and
$$
T\widetilde{M}|_M
=TM\oplus\operatorname{tr}(TM).
$$
A section $N$ of $\operatorname{tr}(TM)$
satisfies $\widetilde{g}(N,\xi)\neq 0$,
because $N\notin TM=\xi^{\perp}$. So,
given $\xi$, there is a unique section $N$
of $\operatorname{tr}(TM)$ with
$$
\widetilde{g}(N,\xi)=1,\qquad
\widetilde{g}(N,N)=0.
$$
We call $N$ a \emph{transversal field};
it is null but not tangent to $M$. We also
write $\eta(X)=\widetilde{g}(X,N)$.

\medskip\noindent
{\bf Gauss--Weingarten equations, I.}
For $X,Y\in\Gamma(TM)$ and $N$ as above we
decompose with respect to
$TM\oplus\operatorname{tr}(TM)$. This gives
the \emph{Gauss formula}
$$
\widetilde{\nabla}_XY
=\nabla_XY+h(X,Y),
$$
and the \emph{Weingarten formula}
$$
\widetilde{\nabla}_XN
=-A_NX+\nabla^t_XN,
$$
where $\nabla_XY$ and $A_NX$ are tangent
to $M$ while $h(X,Y)$ and $\nabla^t_XN$ lie
in $\operatorname{tr}(TM)$. Here $h$ is the
second fundamental form, $A_N$ is the
\emph{shape operator} of $M$ with respect
to $N$, and $\nabla^t$ is the induced
linear connection on the line bundle
$\operatorname{tr}(TM)$, the
\emph{transversal connection}. Since the
transversal bundle has rank one we write
$$
h(X,Y)=B(X,Y)N,\qquad
\nabla^t_XN=\tau(X)N,
$$
with $B$ a bilinear form and $\tau$ a
$1$-form on $M$; these are the formulas
in \cite[Section 2]{navarro2020screen}.
Pairing the Gauss formula with $\xi$ and
using $\nabla_XY\perp\xi$ gives
\begin{equation}
\label{equation-solis-b-formula}
B(X,Y)=\widetilde{g}(
\widetilde{\nabla}_XY,\xi).
\end{equation}
In particular
\begin{equation}
\label{equation-solis-b-xi}
B(X,\xi)=\widetilde{g}(
\widetilde{\nabla}_X\xi,\xi)
=\frac{1}{2}X\left(
\widetilde{g}(\xi,\xi)\right)=0.
\end{equation}

\begin{proposition}
\label{proposition-solis-induced-connection}
The induced connection $\nabla$ is torsion
free and $h$ is symmetric. For all
$X,Y,Z\in\Gamma(TM)$,
$$
(\nabla_Xg)(Y,Z)
=B(X,Y)\eta(Z)+B(X,Z)\eta(Y).
$$
In particular $\nabla$ is a metric
connection if and only if $B=0$, that is,
if and only if $M$ is totally geodesic.
\end{proposition}

\begin{proof}
Since $\widetilde{\nabla}$ is torsion
free,
$$
\begin{aligned}
[X,Y]
&=\widetilde{\nabla}_XY
-\widetilde{\nabla}_YX\\
&=\nabla_XY-\nabla_YX
+h(X,Y)-h(Y,X).
\end{aligned}
$$
The bracket is tangent to $M$. Comparing
tangent parts gives
$\nabla_XY-\nabla_YX=[X,Y]$, and comparing
transversal parts gives $h(X,Y)=h(Y,X)$.

For the second claim we start again in
the ambient space. Since
$\widetilde{\nabla}$ is metric,
$$
\begin{aligned}
X\left(g(Y,Z)\right)
&=\widetilde{g}(\widetilde{\nabla}_XY,Z)
+\widetilde{g}(Y,\widetilde{\nabla}_XZ)\\
&=g(\nabla_XY,Z)+g(Y,\nabla_XZ)\\
&\quad+B(X,Y)\eta(Z)+B(X,Z)\eta(Y).
\end{aligned}
$$
By definition
$$
(\nabla_Xg)(Y,Z)=X\left(g(Y,Z)\right)
-g(\nabla_XY,Z)-g(Y,\nabla_XZ),
$$
which gives the formula. If $B=0$ then
$\nabla g=0$. Conversely, take $Z=\xi$.
Since $\eta(\xi)=1$ and
$B(X,\xi)=0$ by
\eqref{equation-solis-b-xi},
$$
(\nabla_Xg)(Y,\xi)=B(X,Y),
$$
so $\nabla g=0$ forces $B=0$.
\end{proof}

\noindent
So, in general, $\nabla$ is torsion free
\emph{but not metric}. The degenerate
metric $g$ has no unique torsion free
metric connection: if $\nabla$ is one, so
is $\nabla+S\otimes\xi$ for any symmetric
$S$. It has one at all only when $B=0$,
since $L_\xi g=-2B$. The formula above is
\cite[Proposition 2.1(5)]{navarro2020screen}.

\begin{definition}
\label{definition-solis-screen}
The \emph{screen distribution} of $M$
(with respect to the chosen transversal
bundle) is
$$
S(TM)=\left(\operatorname{rad}(TM)
\oplus\operatorname{tr}(TM)
\right)^{\perp}.
$$
\end{definition}

\begin{remark}
\label{remark-solis-screen-tangent}
The plane spanned by $\xi$ and $N$ is a
Lorentzian plane, because the Gram matrix
of $(\xi,N)$ has determinant $-1$. Hence
its orthogonal $S(TM)$ is non-degenerate,
in fact Riemannian (spacelike), of rank
$n$. The screen is tangent to $M$ because
it is orthogonal to the radical and
$TM=\operatorname{rad}(TM)^{\perp}$ by
Lemma~\ref{lemma-solis-radical-rank-one}.
Therefore
$$
TM=S(TM)\perp\operatorname{rad}(TM),
$$
and
$$
T\widetilde{M}|_M=S(TM)\perp
\left(\operatorname{rad}(TM)\oplus
\operatorname{tr}(TM)\right).
$$
Conversely, any vector bundle $S(TM)$
complementary to the radical in $TM$
determines a unique transversal bundle
with these properties
\cite{duggal1996lightlike}. One writes
$(M,g,S(TM))$ for a null hypersurface
with a chosen screen.
\end{remark}

\noindent
The following question was raised: which
objects depend on the choice of the
transversal field (equivalently, of the
screen)? By \eqref{equation-solis-b-formula}
the form $B$ depends only on $\xi$, not on
the screen; hence so do the conditions
$B=0$ and $B=\rho g$ below. By contrast
$\nabla$, $A_N$ and $\tau$ depend on the
screen. The dependence on the
normalization of $\xi$ is discussed after
the Gauss--Weingarten equations II.

\medskip\noindent
{\bf Gauss--Weingarten equations, II.}
Let $P\colon TM\to S(TM)$ be the
projection given by the decomposition
$TM=S(TM)\perp\operatorname{rad}(TM)$; so
$X=PX+\eta(X)\xi$ for $X\in\Gamma(TM)$.
Decomposing now with respect to this
splitting we get
$$
\nabla_XPY=\nabla^*_XPY+h^*(X,PY),
$$
$$
\nabla_X\xi
=-A^*_{\xi}X+\nabla^{*t}_X\xi,
$$
where $\nabla^*_XPY$ and $A^*_{\xi}X$ lie
in $S(TM)$ and $h^*(X,PY)$ and
$\nabla^{*t}_X\xi$ lie in
$\operatorname{rad}(TM)$. Here $\nabla^*$
is the induced connection on the screen,
$h^*$ is the second fundamental form of
the screen and $A^*_{\xi}$ is the shape
operator of the screen. Locally
$$
h^*(X,PY)=C(X,PY)\xi,\qquad
\nabla^{*t}_X\xi=-\tau(X)\xi,
$$
with the same $1$-form $\tau$ as before.
Indeed,
$$
\tau(X)
=\widetilde{g}(\nabla^t_XN,\xi)
=-\widetilde{g}(\nabla^{*t}_X\xi,N);
$$
the second equality follows by
differentiating $\widetilde{g}(\xi,N)=1$
along $X$, using that $A_NX$ is tangent to
$M$ and that $A^*_{\xi}X$ lies in the
screen. With this notation,
$$
\widetilde{\nabla}_XN
=-A_NX+\tau(X)N,
$$
$$
\widetilde{\nabla}_X\xi
=-A^*_{\xi}X-\tau(X)\xi,
$$
where the last formula uses
\eqref{equation-solis-b-xi}. Pairing the
last formula with $Y\in\Gamma(TM)$ and
using \eqref{equation-solis-b-formula}
gives, for $X,Y\in\Gamma(TM)$,
\begin{equation}
\label{equation-solis-b-ashape}
B(X,Y)=g(A^*_{\xi}X,Y).
\end{equation}
Indeed
$B(X,Y)=-\widetilde{g}(Y,
\widetilde{\nabla}_X\xi)$, because
$\widetilde{g}(Y,\xi)=0$.

\medskip\noindent
{\bf Rigging.}
One fixes the transversal field by the
conditions
$$
\widetilde{g}(N,\xi)=1,\qquad
\widetilde{g}(N,N)=0,\qquad
\widetilde{g}(N,W)=0,
$$
the last one for every $W\in
\Gamma(S(TM))$. In practice one starts
from a vector field $V$ along $M$ which is
not tangent to $M$, that is, with
$\widetilde{g}(\xi,V)\neq 0$, and puts
\begin{equation}
\label{equation-solis-rigging}
N=\frac{1}{\widetilde{g}(\xi,V)}
\left(V-\frac{\widetilde{g}(V,V)}
{2\widetilde{g}(\xi,V)}\xi\right).
\end{equation}
A direct computation gives
$\widetilde{g}(N,\xi)=1$ and
$\widetilde{g}(N,N)=0$. For $W$ tangent
to $M$ we have
$\widetilde{g}(N,W)
=\widetilde{g}(V,W)/\widetilde{g}(\xi,V)$,
so the corresponding screen is
$$
S(TM)=TM\cap V^{\perp}.
$$
This construction is fundamental in the
\emph{rigging} technique.

\begin{remark}
\label{remark-solis-rigged-metric}
It looks as if we were introducing a
Riemannian metric on $M$ in addition to
the degenerate metric $g$. Indeed, the
tensor
$$
g+\eta\otimes\eta
$$
is a Riemannian metric on $M$: on the
screen it equals $g$, which is positive
definite, it takes the value $1$ on
$(\xi,\xi)$, and the screen is orthogonal
to $\xi$ for both terms, since $\eta$
vanishes on the screen. When the screen
is integrable, each of its leaves is a
spacelike submanifold of codimension two
of $\widetilde{M}$, hence a Riemannian
manifold.
\end{remark}

\noindent
{\bf Rescaling.} The triple
$(M,g,S(TM))$ still allows rescaling in
$\operatorname{rad}(TM)$ and
$\operatorname{tr}(TM)$. If
$\xi'=\alpha\xi$ for a nowhere vanishing
function $\alpha$, the same screen gives
$N'=\alpha^{-1}N$. From the formulas
above,
$$
B'=\alpha B,\qquad
A^*_{\xi'}=\alpha A^*_{\xi},\qquad
A_{N'}=\alpha^{-1}A_N,
$$
$$
\tau'=\tau-d\log|\alpha|.
$$
For instance,
$\widetilde{\nabla}_X\xi'
=X(\alpha)\xi+\alpha\widetilde{\nabla}_X
\xi$, and comparing radical parts gives
the formula for $\tau'$.

\medskip\noindent
{\bf Totally umbilical null
hypersurfaces.}

\begin{definition}
\label{definition-solis-totally-umbilical}
A null hypersurface $(M,g,S(TM))$ is
\emph{totally umbilical} if there is a
smooth function $\mu$ on $M$ such that
$$
A^*_{\xi}PX=\mu PX
$$
for all $X\in\Gamma(TM)$.
\end{definition}

\noindent
Since $A^*_{\xi}\xi=0$
(Exercise~\ref{exercise-solis-a-star-xi}),
this means $A^*_{\xi}=\mu P$. By
\eqref{equation-solis-b-ashape} and the
non-degeneracy of the screen it is
equivalent to $B=\mu g$, the definition
used in \cite{navarro2020screen}. It does
not depend on the screen. If $\mu=0$ the
hypersurface is totally geodesic.

\begin{example}[The light cone]
\label{example-solis-light-cone}
Let
$$
\Lambda^{n+1}_0=\left\{x\in
\mathbb{R}^{n+2}_1:
\left\langle{} x,x \right\rangle{}=0,\
x\neq 0\right\}.
$$
The gradient of
$\left\langle{} x,x \right\rangle{}$ is
$2x$, so the position vector field
$\xi=x$ is normal, hence spans the
radical. Since $\widetilde{\nabla}$ is the
flat connection,
$\widetilde{\nabla}_X\xi=X$ for every
$X$. For $X$ in the screen,
$$
\widetilde{\nabla}_{PX}\xi=PX,\qquad
\nabla_{PX}\xi=PX,\qquad
A^*_{\xi}PX=-PX.
$$
In general
$\widetilde{\nabla}_X\xi=X=PX+\eta(X)\xi$,
and comparing with the Gauss--Weingarten
equations gives $A^*_{\xi}=-P$ and
$\tau=-\eta$, for any screen. So the light
cone is a totally umbilical null
hypersurface, with $\mu=-1$.

The picture drawn in class, for
$n=1$: the radical $\xi$ points along the
generators of the cone; the transversal
$N$ is the other null direction
orthogonal to a horizontal circle on the
cone; and the screen is tangent to the
horizontal circles, parallel to the
$x_1x_2$-plane.
\end{example}

\begin{definition}
\label{definition-solis-screen-umbilical}
The screen $S(TM)$ is \emph{totally
umbilical} if there is a smooth function
$\lambda$ such that
$$
A_NX=\lambda PX
$$
for all $X\in\Gamma(TM)$.
\end{definition}

\noindent
The notes wrote $A_NX=\lambda X$; since
$A_N$ takes values in the screen
(Exercise~\ref{exercise-solis-%
shape-in-screen}),
the right-hand side must be $\lambda PX$.

\begin{example}[The light cone, again]
\label{example-solis-light-cone-screen}
On $\Lambda^{n+1}_0$ take a constant
vector $V$ and write
$a=\left\langle{} \xi,V \right\rangle{}$
and
$c=\left\langle{} V,V \right\rangle{}$,
on the open set where $a\neq 0$. The
rigging \eqref{equation-solis-rigging}
gives
$$
N=\frac{1}{a}\left(V-\frac{c}{2a}\xi
\right).
$$
For $X$ tangent to the cone,
$X(a)=\left\langle{} X,V \right\rangle{}
=a\,\eta(X)$, because $PX\perp V$.
Differentiating, and writing $V$ in terms
of $N$ and $\xi$, one finds
$$
\widetilde{\nabla}_XN
=-\frac{c}{2a^2}PX-\eta(X)N.
$$
Hence
$$
A_NX=\frac{c}{2a^2}PX
=-\frac{c}{2a^2}A^*_{\xi}X,
$$
so the screen is totally umbilical. With
$c=1$ this is the formula
$$
A_NX=-\frac{1}{2\left\langle{} \xi,V
\right\rangle{}^2}A^*_{\xi}X
$$
written in class.
\end{example}

\medskip\noindent
{\bf Robertson--Walker spacetimes.}
These are warped products
$$
\widetilde{M}=-I\times_{\rho}F,\qquad
\widetilde{g}=-dt^2+\rho(t)^2g_F,
$$
where $I\subset\mathbb{R}$ is an
interval, $\rho>0$ and $F$ is a
Riemannian manifold of constant
curvature; the minus sign indicates that
the base carries the negative definite
metric $-dt^2$. If $F$ is only assumed
Riemannian one speaks of generalized
Robertson--Walker (GRW) spacetimes.

\begin{example}
\label{example-solis-rw-space-forms}
\begin{itemize}
\item (De Sitter spacetime.) The map
$(t,y)\mapsto \sinh(t)e_0+\cosh(t)y$, with
$y$ in the unit sphere of the orthogonal
complement of $e_0$, is an isometry
$$
-\mathbb{R}\times_{\cosh}\mathbb{S}^{n+1}
\to\mathbb{S}^{n+2}_1.
$$
\item (Anti de Sitter spacetime.)
Similarly
$(t,y)\mapsto\sin(t)e+\cos(t)y$, with
$e\in\mathbb{R}^{n+3}_2$ a unit timelike
vector and $y$ in the hyperbolic space
$\mathbb{H}^{n+1}$ of the orthogonal
complement of $e$, is an isometric
embedding
$$
-\left(-\tfrac{\pi}{2},\tfrac{\pi}{2}
\right)\times_{\cos}\mathbb{H}^{n+1}
\to\mathbb{H}^{n+2}_1
$$
onto an open subset of anti de Sitter
space. (The notes had
$\mathbb{H}^n_1$ with a question mark.)
\item (Through distance functions.) The
notes record only the start of this
example: a distance function
$d\colon\mathbb{S}^{n+1}\to\mathbb{R}$ and
a subset $S\subset\mathbb{S}^{n+1}$. A
construction of this kind is the
following. In $-I\times_{\rho}F$, the
graph $t=f(x)$ of a function
$f\colon F\to I$ is a null hypersurface
if and only if
$$
|\operatorname{grad}f|=\rho\circ f,
$$
as one checks by computing the norm of
the gradient of $t-f$. If $d_S$ is the
distance to a subset $S\subset F$ and
$\phi'=1/\rho$, then $f=\phi^{-1}\circ
d_S$ satisfies this equation wherever
$d_S$ is smooth and takes values in
$\phi(I)$; the norm of $\operatorname{grad}f$
is taken in $g_F$. Null hypersurfaces of
GRW spacetimes are studied as such graphs
in \cite{navarro2016null}; see also
\cite[Example 3.3]{navarro2020screen}.
\end{itemize}
\end{example}

\begin{remark}
\label{remark-solis-rw-screen}
In a Robertson--Walker space there is a
natural screen: the tangent spaces of the
spacelike submanifolds
$$
M\cap\left(\{t\}\times F\right).
$$
This is the screen used in
\cite[Example 3.3]{navarro2020screen}. It
is the screen given by the rigging
\eqref{equation-solis-rigging} with
$V=\partial_t$, since the orthogonal of
$\partial_t$ is tangent to the slices.
\end{remark}

\medskip\noindent
{\bf Screen quasi-conformal
hypersurfaces.}

\begin{definition}
\label{definition-solis-quasi-conformal}
A null hypersurface $(M,g,S(TM))$ is
\emph{screen quasi-conformal} if there are
smooth functions $\varphi$ and $\psi$
with
$$
A_N=\varphi A^*_{\xi}+\psi P.
$$
\end{definition}

\noindent
The notes left the formula incomplete;
it is taken from
\cite[Definition 3.1]{navarro2020screen}.
When $\psi=0$ and $\varphi$ vanishes
nowhere, these are the screen conformal
hypersurfaces of Atindogbe and Duggal
(not checked), like the light cone of
Example~\ref{example-solis-%
light-cone-screen} for a non-null $V$.
Every null hypersurface of a GRW
spacetime, with the screen of
Remark~\ref{remark-solis-rw-screen}, is
screen quasi-conformal
\cite[Example 3.3]{navarro2020screen}, and
for a screen quasi-conformal hypersurface
the screen is integrable
\cite[Theorem 3.6]{navarro2020screen}.
Applications: classification of null
hypersurfaces which are isoparametric or
Einstein in Lorentzian space forms. A
corollary on the Lorentzian space forms
$\mathbb{K}^{n+2}_1(\overline{c})$
written as Robertson--Walker spacetimes
$-I\times_{\rho}F$ was stated in class,
but the notes do not record it.

\begin{definition}
\label{definition-solis-einstein}
Let $(M,g,S(TM))$ be a null hypersurface
whose induced Ricci tensor
$\operatorname{Ric}$ (the trace of
$Z\mapsto R(Z,X)Y$, with $R$ the curvature
of $\nabla$) is symmetric. Then $M$ is
\emph{Einstein} if there is
$k\in C^{\infty}(M)$ such that
$$
\operatorname{Ric}(X,Y)=k\,g(X,Y)
$$
for all $X,Y\in\Gamma(TM)$.
\end{definition}

\noindent
In \cite[Definition 5.3]{navarro2020screen}
the function $k$ is also required to be
constant along the screen.

\begin{remark}
\label{remark-solis-schur}
In Robertson--Walker spaces there are
null hypersurfaces with
$\operatorname{Ric}=\alpha g$ where
$\alpha$ is not constant. So for null
hypersurfaces Schur's lemma, which would
force $\alpha$ to be constant, does not
hold. An explicit example in de Sitter
space is
\cite[Example 5.4]{navarro2020screen}.
\end{remark}

\noindent
Let $\mathbb{K}^{n+2}_1(\overline{c})$
denote the Lorentzian space form of
curvature $\overline{c}\neq 0$, viewed as
a quadric in a semi-Euclidean space
$\mathbb{R}^{n+3}_s$.

\begin{proposition}
\label{proposition-solis-umbilic-hyperplane}
Let $(M,g,S(TM))\subset
\mathbb{K}^{n+2}_1(\overline{c})$ be a
complete, connected, totally umbilical
null hypersurface. Then $M$ is the
intersection of
$\mathbb{K}^{n+2}_1(\overline{c})$ with a
hyperplane of $\mathbb{R}^{n+3}_s$.
\end{proposition}

\noindent
Null totally umbilical hypersurfaces of
space forms are characterized in
\cite{navarro2016null}; compare
\cite[Example 5.4]{navarro2020screen}.
The statement above was not checked
against that paper (not checked).

\begin{proposition}
\label{proposition-solis-umbilic-section}
Let $(M,g,S(TM))\subset
\mathbb{K}^{n+2}_1(\overline{c})$ be a
totally umbilical null hypersurface and
$S\subset M$ a spacelike hypersurface of
$M$. Then $S$ is totally umbilical in $M$
if and only if $S$ is contained in a
totally geodesic hypersurface of
$\mathbb{K}^{n+2}_1(\overline{c})$
(not checked).
\end{proposition}

\medskip\noindent
{\bf Session 2. Causal theory.}
Lorentz--Minkowski space
$\mathbb{R}^{n+1}_1$ is $\mathbb{R}^{n+1}$
with the usual scalar product of index
one. In a Lorentzian manifold a continuous
choice, at each point, of one of the two
cones of timelike tangent vectors is a
\emph{time orientation}; a
\emph{spacetime} is a connected
time-oriented Lorentzian manifold
$(M,g)$.

\medskip\noindent
{\bf Causality.} We write $x\ll y$ if
there is a future-directed timelike curve
from $x$ to $y$, and $x\leq y$ if there is
a future-directed causal curve from $x$ to
$y$ or $x=y$. The \emph{chronological
future} and the \emph{causal future} of
$x$ are
$$
I^+(x)=\{y:x\ll y\},\qquad
J^+(x)=\{y:x\leq y\},
$$
and $I^-(x)$, $J^-(x)$ are defined
analogously.

\begin{proposition}
\label{proposition-solis-i-open}
In any spacetime $(M,g)$ the sets
$I^+(x)$ and $I^-(x)$ are open in the
manifold topology.
\end{proposition}

\begin{proof}
Let $y\in I^+(x)$, reached by a timelike
curve $\sigma$ from $x$. Take a convex
neighbourhood $U$ of $y$ and a point $z$
of $\sigma$ before $y$, so close to $y$
that the arc of $\sigma$ from $z$ to $y$
lies in $U$. By \cite[Lemma 14.2]{oneill},
via the exponential map at $z$, the set of
points of $U$ reached from $z$ by future
timelike curves in $U$ is the image of an
open cone, so it is an open set containing
$y$, and
each of its points lies in $I^+(x)$:
follow $\sigma$ up to $z$, then a timelike
geodesic. See \cite[Lemma 14.3]{oneill}.
\end{proof}

\begin{definition}
\label{definition-solis-chronological}
A spacetime $(M,g)$ is
\emph{chronological} if there are no
closed timelike curves, that is, $x\ll x$
for no $x$.
\end{definition}

\begin{proposition}
\label{proposition-solis-compact}
If $M$ is compact, then $(M,g)$ is not
chronological.
\end{proposition}

\begin{proof}
Every point $p$ lies in some $I^+(x)$:
take $x$ on a short past-directed
timelike curve from $p$. Hence
$\{I^+(x):x\in M\}$ is an open cover, and
by compactness it has a finite subcover
$I^+(x_1),\ldots,I^+(x_k)$, which we take
minimal. The point $x_1$ lies in some
$I^+(x_i)$. If $i\neq 1$, then
$I^+(x_1)\subset I^+(x_i)$ by
transitivity, contradicting minimality.
So $x_1\in I^+(x_1)$, that is, there is a
closed timelike curve through $x_1$
\cite[Lemma 14.10]{oneill}.
\end{proof}

\begin{definition}
\label{definition-solis-alexandrov}
The \emph{Alexandrov topology}
$\mathcal{A}$ of $(M,g)$ is the topology
generated by the basis
$$
\{I^+(x)\cap I^-(y):x\ll y\}.
$$
\end{definition}

\begin{remark}
\label{remark-solis-alexandrov}
Let $\mathcal{M}$ be the manifold topology.
By Proposition~\ref{proposition-solis-%
i-open},
$\mathcal{A}\subset\mathcal{M}$, but in
general they are different. This was
shown in class with the Lorentz cylinder
$\mathbb{S}^1_1\times\mathbb{R}$ (the
circle carries the negative definite
metric): there $I^+(x)=M$ for every $x$
\cite[p.~402]{oneill}, so $\mathcal{A}$ is
the trivial topology.
\end{remark}

\begin{definition}
\label{definition-solis-strongly-causal}
$(M,g)$ is \emph{strongly causal} if
$\mathcal{A}=\mathcal{M}$.
\end{definition}

\noindent
The usual definition is different: strong
causality holds at $p$ if every
neighbourhood of $p$ contains a
neighbourhood $V$ of $p$ such that every
causal curve segment with endpoints in $V$
stays in the first neighbourhood
\cite[Definition 14.11]{oneill}. That this
is equivalent to $\mathcal{A}=\mathcal{M}$
is a classical theorem of Kronheimer and
Penrose (not checked).

\begin{proposition}[Push-up]
\label{proposition-solis-push-up}
In any spacetime, if $x\ll y$ and
$y\leq z$, or if $x\leq y$ and $y\ll z$,
then $x\ll z$.
\end{proposition}

\noindent
The notes left the statement blank; it is
\cite[Corollary 14.1]{oneill}.

\begin{proposition}
\label{proposition-solis-j-closure}
$J^+(x)\subset\overline{I^+(x)}$; moreover
$I^+(x)$ is the interior of $J^+(x)$.
\end{proposition}

\noindent
See \cite[Lemma 14.6]{oneill}. The set
$J^+(x)$ is not always closed: in
$\mathbb{R}^2_1$ with one point deleted,
points behind the deleted point, on the
null ray from $x$ through it, are in the
closure of $I^+(x)$ but not in $J^+(x)$
\cite[Example 14.4]{oneill}.

\begin{definition}
\label{definition-solis-causally-simple}
$(M,g)$ is \emph{causally simple} if it
has no closed causal curves and $J^+(x)$
and $J^-(x)$ are closed for every $x$.
\end{definition}

\noindent
The notes only required $J^+(x)$ closed;
the definition above is the standard one
\cite{minguzzi2019causality}.

\begin{definition}
\label{definition-solis-globally-hyperbolic}
$(M,g)$ is \emph{globally hyperbolic} if
\begin{itemize}
\item it is strongly causal, and
\item $J^+(p)\cap J^-(q)$ is compact for
all $p,q\in M$.
\end{itemize}
\end{definition}

\noindent
The two conditions were left blank in the
notes; we use
\cite[p.~411]{oneill}.

\begin{proposition}
\label{proposition-solis-gh-equivalent}
The following are equivalent.
\begin{enumerate}
\item $(M,g)$ is globally hyperbolic.
\item $(M,g)$ is strongly causal and the
spaces $C(p,q)$ of continuous causal
curves from $p$ to $q$, with the $C^0$
topology on curves up to
reparametrisation, are compact.
(For piecewise smooth curves this
fails already in Minkowski space.)
\item $(M,g)$ has a Cauchy hypersurface.
\end{enumerate}
\end{proposition}

\noindent
That (3) implies (1) is
\cite[Corollary 14.39]{oneill}; the other
implications (Leray, Geroch) were stated
in class and are not checked here (not
checked). The strong causality in (2) is
added to the notes, following the classical
formulation.

\begin{theorem}
\label{theorem-solis-cauchy-temporal}
If $(M,g)$ is globally hyperbolic, then
it has a Cauchy temporal function (not
checked).
\end{theorem}

\noindent
This is due to Geroch for a continuous
Cauchy time function and to Bernal and
Sánchez for a smooth one; see
\cite{minguzzi2019causality}.

\medskip\noindent
{\bf Lorentzian distance.}
For a future-directed causal curve
$\gamma$ put
$$
L(\gamma)=\int\sqrt{-g(\gamma',\gamma')},
$$
and define
$$
\tau(x,y)=\sup L(\gamma),
$$
the supremum over future-directed causal
curves $\gamma$ from $x$ to $y$, with
$\tau(x,y)=0$ if there is none. (This
$\tau$ is unrelated to the $1$-form $\tau$
of Session 1.)

\begin{remark}
\label{remark-solis-lorentzian-distance}
If $y\notin J^+(x)$ then $\tau(x,y)=0$,
and $\tau$ is not symmetric. Instead of
the triangle inequality there is the
reverse triangle inequality
$\tau(x,y)+\tau(y,z)\leq\tau(x,z)$ for
$x\leq y\leq z$
\cite[Lemma 14.16]{oneill}. The function
$\tau$ is lower semicontinuous
\cite[Lemma 14.17]{oneill}, but it need
not be continuous. The example in class
was $\mathbb{R}^2_1$ with a set $L$
removed; the notes do not say which set.
A closed spacelike segment works
\cite[Example 14.18]{oneill}. Also,
$\tau$ is continuous on globally
hyperbolic spacetimes
\cite[Lemma 14.21]{oneill}.
\end{remark}

\medskip\noindent
{\bf Session 3. General relativity,
singularities and black holes.}
On a spacetime $(M,g)$ the \emph{Einstein
field equations} are
$$
\operatorname{Ric}-\frac{1}{2}Sg
+\Lambda g=T,
$$
where $S$ is the scalar curvature,
$\Lambda$ is the \emph{cosmological
constant} and $T$ is the
\emph{energy-momentum tensor} (in units
with $8\pi G=c=1$).

\begin{quote}
Matter tells spacetime how to curve.
Curvature tells matter how to move.
\hfill (J.~A.~Wheeler, as quoted in
class.)
\end{quote}

\medskip\noindent
{\bf Solutions.}
\begin{itemize}
\item (Lorentz--Minkowski, the setting of
special relativity, 1905.) $T=0$, that
is, there is no matter or energy (at
large scale), $\Lambda=0$,
$M=\mathbb{R}^{n+1}$ and $g$ is the
Minkowski scalar product.

\item (Schwarzschild, 1916.) In dimension
four, $M=\mathbb{R}\times(2m,\infty)\times
\mathbb{S}^2$ with
$$
g=-f(r)\,dt^2+\frac{1}{f(r)}\,dr^2
+r^2g_{\mathbb{S}^2},
$$
where
$$
f(r)=1-\frac{2m}{r}.
$$
The expression degenerates at $r=2m$.
Here $\Lambda=0$ and $T=0$. Schwarzschild
died a few months after publishing it,
during the First World War.

It is interpreted as the exterior of a
black hole. The coordinate changes written
in the notes are garbled; we use the
standard ones. With
$r^*=r+2m\log(r/2m-1)$ one has
$dr^*=dr/f$, and in the coordinate
$v=t+r^*$ the metric becomes
$$
g=-f(r)\,dv^2+2\,dv\,dr
+r^2g_{\mathbb{S}^2},
$$
which extends smoothly to all $r>0$; with
$u=t-r^*$ instead, one gets the analogous
extension with $-2\,du\,dr$. So there are
two families of radial null geodesics,
and for each there is a coordinate system
in which they become straight lines. The
logarithm forces $r>2m$ in the original
coordinates. Depending on the choice, one
finds a region $r<2m$ which future
directed light rays can enter but not
leave (a black hole) or leave but not
enter (a white hole).

\item (De Sitter, 1917.)
$M=\mathbb{S}^{n+1}_1\subset
\mathbb{R}^{n+2}_1$, diffeomorphic to
$\mathbb{R}\times\mathbb{S}^n$, with
$T=0$ and
$$
\Lambda=\frac{n(n-1)}{2}.
$$
This value follows from
$\operatorname{Ric}=ng$ and $S=n(n+1)$.
Writing it as
$-\mathbb{R}\times_{\cosh}\mathbb{S}^n$
(Example~\ref{example-solis-rw-space-forms}),
this is interpreted as a universe that
\emph{is expanding}.

\item (Einstein cylinder, 1917; the
universe that does not expand.)
$M=\mathbb{R}\times\mathbb{S}^n$ with
$g=-dt^2+g_{\mathbb{S}^n}$ and a perfect
fluid
$$
T=(\mu+p)U^{\flat}\otimes U^{\flat}+pg,
$$
where $U=\partial_t$ and $\mu$, $p$ are
the energy density and the pressure. The
field equations reduce to
$$
\mu=\frac{n(n-1)}{2}-\Lambda,\qquad
p=\Lambda-\frac{(n-1)(n-2)}{2};
$$
for $n=3$ this gives $\Lambda=1+p$ (the
notes wrote $1+\mu$). The universe does not
expand because all slices $\{t\}\times
\mathbb{S}^n$ have the same size.

\item (Hubble's expanding universe.)
Hubble worked at the Mount Wilson
Observatory; his law, published in 1929,
says that galaxies recede with velocity
proportional to their distance,
$$
v=H_0d.
$$
Large distance goes with large velocity.
This is related to the Big Bang: going
back in time, the universe was more and
more contracted.

\item (FLRW spacetimes; A.~Friedmann,
G.~Lemaître, H.~Robertson, A.~Walker.)
$M=I\times\mathbb{K}^n(k)$ with
$$
g=-dt^2+\rho(t)^2g_{\mathbb{K}^n(k)},
$$
$T$ a perfect fluid as above. The notes
stop at ``$\rho^2=\ldots$''; for $n=3$ the
field equations give the Friedmann
equation
$$
\left(\rho'\right)^2
=\frac{\mu+\Lambda}{3}\rho^2-k.
$$
\end{itemize}

\medskip\noindent
{\bf Singularities.}

\begin{definition}
\label{definition-solis-singular}
A spacetime is \emph{singular} if it is
inextendible (maximal) and causally
geodesically incomplete.
\end{definition}

\noindent
It is as if particles suddenly
disappeared, at the singularity, at some
point of their trajectory through time.

\medskip\noindent
{\bf Kruskal--Szekeres extension.} In the
null coordinates $u=t-r^*$ and
$v=t+r^*$ the Schwarzschild metric is
$$
g=-\frac{2m\,e^{-r/2m}}{r}\,
e^{(v-u)/4m}\,du\,dv
+r^2g_{\mathbb{S}^2},
$$
since $f(r)=\frac{2m}{r}e^{-r/2m}
e^{r^*/2m}$. With $U=-e^{-u/4m}$ and
$V=e^{v/4m}$,
$$
g=-\frac{32m^3}{r}e^{-r/2m}\,dU\,dV
+r^2g_{\mathbb{S}^2},
$$
where $r$ is determined by
$UV=-(r/2m-1)e^{r/2m}$. This expression
makes sense for all $r>0$ and defines
Kruskal spacetime, which contains the
black hole region; compare
\cite[Definitions 13.22 and 13.25]{oneill},
with other constants.

\begin{theorem}[Birkhoff, 1923]
\label{theorem-solis-birkhoff}
Every spherically symmetric solution of
the vacuum Einstein equations is
(locally) isometric to an open subset of
Schwarzschild spacetime.
\end{theorem}

\noindent
The precise form of the statement
(including the case $m=0$, Minkowski
space, and local versus global) was not
recorded (not checked).

\begin{remark}[Penrose, 1969]
\label{remark-solis-cosmic-censorship}
Penrose's (weak) cosmic censorship
conjecture: let $(S,g,h)$ be generic,
asymptotically flat, regular initial data
for the vacuum Einstein equations (or
suitable matter). If the resulting
spacetime contains singularities, then
these lie behind an event horizon.
Without genericity it fails for the
Einstein--scalar field system
(Christodoulou) (not checked).
\end{remark}

\medskip\noindent
{\bf Singularity theorems.}

\begin{theorem}[Penrose]
\label{theorem-solis-penrose}
Let $(M,g)$ be a spacetime of dimension
$n+1$ such that
\begin{itemize}
\item $M$ is globally hyperbolic,
\item $\operatorname{Ric}(K,K)\geq 0$ for
every null vector $K$,
\item $M$ is future null geodesically
complete, and
\item $M$ contains a compact achronal
spacelike submanifold $S$ of codimension
two which is a trapped surface, that is,
whose mean curvature vector is past
pointing timelike, with O'Neill's
convention that the second fundamental
form is the normal part of
$\nabla_VW$.
\end{itemize}
Then
$E^+(S)=J^+(S)\setminus I^+(S)$ is a
Cauchy hypersurface. In particular, it is
compact, so the Cauchy hypersurfaces of
$M$ are compact
\cite[Corollary 14.32]{oneill}.
\end{theorem}

\noindent
This is \cite[Theorem 14.61]{oneill},
where global hyperbolicity is stated as
the existence of a Cauchy hypersurface
(Proposition~\ref{proposition-solis-%
gh-equivalent}).
The notes dated it 1969 and omitted
``achronal''. As a singularity theorem it
reads: if $M$ satisfies the condition on
the Ricci tensor, has a non-compact
Cauchy hypersurface and contains a
compact achronal trapped surface, then
$M$ is future null geodesically
incomplete \cite[Theorem 14.61,
Corollary A]{oneill}. It is Penrose's
theorem of 1965 \cite{penrose1965collapse}.

\medskip\noindent
{\bf Conformal compactification.}

\begin{example}[Poincaré disc]
\label{example-solis-poincare-disc}
On the unit disc put
$$
g=\frac{4\left(dx^2+dy^2\right)}
{\left(1-x^2-y^2\right)^2}.
$$
This complete metric of curvature $-1$ is
conformal to the Euclidean metric of the
disc, whose closure adds a circle at
infinity. (The notes omitted the factor
$4$ and the square.)
\end{example}

\begin{example}[Lorentzian space forms]
\label{example-solis-einstein-cylinder}
The model is the Einstein cylinder
$M_0=\mathbb{R}\times\mathbb{S}^n$ with
$g_0=-dT^2+g_{\mathbb{S}^n}$. Minkowski
space is conformal to the open subset
$|T|+\chi<\pi$ of $M_0$, where $\chi$ is
the distance on $\mathbb{S}^n$ to a fixed
point, via
$$
T\pm\chi=2\arctan(t\pm r).
$$
The boundary of this region in $M_0$
gives the conformal boundary of Minkowski
space, with future and past null infinity
$\mathcal{I}^{\pm}$ (not checked).
\end{example}

\medskip\noindent
{\bf Asymptotically flat spacetimes.}

\begin{definition}
\label{definition-solis-black-hole}
Suppose that $(M,g)$ admits a conformal
boundary with future null infinity
$\mathcal{I}^+$. The \emph{black hole}
region is
$$
\operatorname{BH}
=M\setminus J^-(\mathcal{I}^+),
$$
what remains after removing the causal
past of $\mathcal{I}^+$; here $J^-$ is
taken in $M\cup\mathcal{I}^+$ and then
intersected with $M$. The \emph{event
horizon} is
$$
\operatorname{HE}=\partial
\operatorname{BH}.
$$
\end{definition}

\noindent
The notes say ``chronological past'', but
the formula uses $J^-$, the causal past.

\medskip\noindent
{\bf Session 4. Lorentzian length
spaces.} Daniel's notes do not record
this session. The notion of Lorentzian
length space was introduced in
\cite{kunzinger2018lorentzian}.

\medskip\noindent
{\bf Problem session.}
Let $(M,g,S(TM))$ be a null hypersurface
as in Session 1.

\begin{exercise}
\label{exercise-solis-shape-in-screen}
The image of the shape operator $A_N$
lies in the screen.
\end{exercise}

\begin{proof}
Recall that
$$
\begin{aligned}
T\widetilde{M}|_M
&=TM\oplus\operatorname{tr}(TM)\\
&=S(TM)\oplus\operatorname{rad}(TM)
\oplus\operatorname{tr}(TM),
\end{aligned}
$$
so $X=X^*+a\xi+bN$ with
$X^*\in S(TM)$. Applying the metric,
$$
\widetilde{g}(X,N)
=\widetilde{g}(X^*,N)
+a\,\widetilde{g}(\xi,N)
+b\,\widetilde{g}(N,N)=a,
$$
since the first and last terms vanish
and $\widetilde{g}(\xi,N)=1$. So a vector
$X$ tangent to $M$ lies in the screen if
and only if $\widetilde{g}(X,N)=0$, and we
must show $\widetilde{g}(A_NX,N)=0$. Now
$$
\begin{aligned}
0&=X\left(\widetilde{g}(N,N)\right)
=2\widetilde{g}(\widetilde{\nabla}_XN,N)\\
&=-2\widetilde{g}(A_NX,N)
+2\widetilde{g}(\nabla^t_XN,N)\\
&=-2\widetilde{g}(A_NX,N),
\end{aligned}
$$
because $\nabla^t_XN=\tau(X)N$ and
$\widetilde{g}(N,N)=0$.
\end{proof}

\begin{exercise}
\label{exercise-solis-a-star-xi}
$A^*_{\xi}\xi=0$.
\end{exercise}

\begin{proof}
The vector $A^*_{\xi}\xi$ lies in the
screen. For $Y\in\Gamma(TM)$, by
\eqref{equation-solis-b-ashape}, the
symmetry of $B$
(Proposition~\ref{proposition-solis-%
induced-connection})
and \eqref{equation-solis-b-xi},
$$
g(A^*_{\xi}\xi,Y)=B(\xi,Y)
=B(Y,\xi)=0.
$$
Since the screen is non-degenerate,
$A^*_{\xi}\xi=0$.
\end{proof}

\begin{exercise}
\label{exercise-solis-a-star-symmetric}
$A^*_{\xi}$ is symmetric with respect to
$g$, that is,
$$
g(A^*_{\xi}X,Y)=g(X,A^*_{\xi}Y)
$$
for all $X,Y\in\Gamma(TM)$.
\end{exercise}

\begin{proof}
We use that
$\widetilde{g}(X,\xi)=0=
\widetilde{g}(Y,\xi)$. Differentiating,
$$
\begin{aligned}
0&=X\left(\widetilde{g}(\xi,Y)\right)
=\widetilde{g}(\widetilde{\nabla}_X\xi,Y)
+\widetilde{g}(\xi,\widetilde{\nabla}_XY),
\\
0&=Y\left(\widetilde{g}(\xi,X)\right)
=\widetilde{g}(\widetilde{\nabla}_Y\xi,X)
+\widetilde{g}(\xi,\widetilde{\nabla}_YX).
\end{aligned}
$$
Subtracting,
$$
0=\widetilde{g}(\widetilde{\nabla}_X\xi,Y)
-\widetilde{g}(\widetilde{\nabla}_Y\xi,X)
+\widetilde{g}(\xi,[X,Y]),
$$
and the last term vanishes because
$[X,Y]$ is tangent to $M$. Hence
$$
\widetilde{g}(\widetilde{\nabla}_X\xi,Y)
=\widetilde{g}(\widetilde{\nabla}_Y\xi,X).
$$
On the other hand,
$$
\begin{aligned}
\widetilde{\nabla}_X\xi
&=\nabla_X\xi+h(X,\xi)\\
&=-A^*_{\xi}X+\nabla^{*t}_X\xi
+h(X,\xi),
\end{aligned}
$$
and $\nabla^{*t}_X\xi$ is a multiple of
$\xi$, so
$$
\widetilde{g}(\widetilde{\nabla}_X\xi,Y)
=-g(A^*_{\xi}X,Y)
+\widetilde{g}(h(X,\xi),Y),
$$
and similarly with $X$ and $Y$
exchanged. It remains to see that
$h(X,\xi)=0$. Since $h(X,\xi)=B(X,\xi)N$
and $\widetilde{g}(N,\xi)=1$, it is enough
to show $\widetilde{g}(h(X,\xi),\xi)=0$.
We use $\widetilde{g}(\xi,\xi)=0$:
$$
\begin{aligned}
0&=X\left(\widetilde{g}(\xi,\xi)\right)
=2\widetilde{g}(\widetilde{\nabla}_X\xi,
\xi)\\
&=2\widetilde{g}(\nabla_X\xi+h(X,\xi),
\xi)
=2\widetilde{g}(h(X,\xi),\xi),
\end{aligned}
$$
because $\nabla_X\xi$ is tangent to $M$,
hence orthogonal to $\xi$. Therefore
$g(A^*_{\xi}X,Y)=g(A^*_{\xi}Y,X)$.
\end{proof}

\begin{exercise}
\label{exercise-solis-ruled}
Every null hypersurface is ruled by null
geodesics: the integral curves of $\xi$
are, after reparametrization, null
geodesics of $\widetilde{M}$.
\end{exercise}

\begin{proof}
By the Gauss--Weingarten equations,
\eqref{equation-solis-b-xi} and
Exercise~\ref{exercise-solis-a-star-xi},
$$
\begin{aligned}
\widetilde{\nabla}_{\xi}\xi
&=\nabla_{\xi}\xi+B(\xi,\xi)N\\
&=-A^*_{\xi}\xi+\nabla^{*t}_{\xi}\xi
=-\tau(\xi)\xi.
\end{aligned}
$$
This is a kind of geodesic equation. Let
$\gamma$ be an integral curve of $\xi$,
so $\widetilde{\nabla}_{\gamma'}\gamma'
=k\gamma'$ with $k=-\tau(\xi)\circ\gamma$.
Choose $s(t)$ with
$$
s'(t)=\exp\int_0^tk,
$$
and let $\sigma$ be defined by
$\sigma(s(t))=\gamma(t)$. Then
$\gamma'=s'\,\sigma'$, so
$$
k\gamma'
=\widetilde{\nabla}_{\gamma'}\gamma'
=s''\sigma'+(s')^2
\widetilde{\nabla}_{\sigma'}\sigma'
=k\gamma'+(s')^2
\widetilde{\nabla}_{\sigma'}\sigma',
$$
hence $\widetilde{\nabla}_{\sigma'}
\sigma'=0$, and $\sigma$ is a null
geodesic contained in $M$.
\end{proof}

\section{Análisis geométrico en superficies}
\label{section-alias-geometric-analysis}

\medskip\noindent
{\bf Curso Avanzado 2.}

Luis Alias,
Universidad de Murcia-Spain.

\medskip\noindent
{\bf Schedule.}

Week 2, Monday to Thursday, 11:00--12:30
(15--18 July 2024).

\medskip\noindent
{\bf Official listing.}

\href{%
https://naturales.ues.edu.sv/%
ecimpasv2024/advanced-curse-2/}{%
Programme entry.}

\medskip\noindent
{\bf Abstract.}

Este curso avanzado va dirigido a estudiantes
que hayan completado al menos un primer curso
introductorio sobre geometría diferencial de
curvas y superficies. En particular, este
curso está pensado como la continuación del
curso previo ``Introducción a la geometría
diferencial de superficies'', que se
impartirá en esta misma escuela CIMPA la
semana anterior y que contiene los
fundamentos de nuestro curso.

Nuestro objetivo es introducir a los
estudiantes en las técnicas avanzadas del
análisis geométrico y en sus aplicaciones al
estudio de la geometría global de superficies
en el espacio euclídeo tridimensional. En
particular, nos proponemos hacer accesibles
al nivel de los estudiantes distintos
resultados de investigación, más o menos
recientes, sobre el comportamiento global de
la curvatura y la topología de superficies en
el espacio euclídeo.

Este curso está basado en un curso anterior
titulado ``Análisis geométrico y geometría
global de superficies: Una introducción
elemental'', que fue impartido por el
conferenciante en la ``XIV Escuela Brasileña
de Geometría Diferencial'' que fue celebrada
en la Universidad Federal de Bahía, Salvador
de Bahía (Brasil), en 2006, en honor de
Shiing-Shen Chern. Los contenidos del curso
son los siguientes: 1) Cálculo intrínseco en
superficies; 2) Superficies convexas. El
teorema de Hadamard; 3) Las fórmulas de
Minkowski y sus aplicaciones; 4) Superficies
de curvatura media constante. El teorema de
Alexandrov; 5) Estimaciones geométricas del
primer valor propio del operador de
Laplace-Beltrami.

\medskip\noindent
{\bf Sources of these notes.}
These notes come from Daniel's notes taken
in Spanish during the course. They cover
three lectures and a problem session.
The course follows the lecturer's
booklet \cite{alias2006analisis}, written
for the XIV Escola de Geometria
Diferencial (Salvador de Bahia, 2006).
Statements and proofs of the first two
lectures were checked against it and
completed from it; page numbers below
refer to that booklet. The third lecture
(minimal and maximal graphs, Bernstein and
Calabi--Bernstein theorems) is not in the
booklet; its computations were completed
and checked by hand here. Daniel's notes
stop before topics 4 and 5 of the abstract
(Alexandrov's theorem and the first
eigenvalue). The last part of these notes
summarizes them from Chapters 4 and 5 of
the booklet, not from the lectures.
Mistakes in the notes taken in class were
corrected silently. Points that could not
be verified are marked ``(not checked)''.

The theorem of Hadamard proved below is
Hadamard's theorem on ovaloids. It is not
the Cartan--Hadamard theorem of the
Differential Geometry chapter, Theorem
\ref{differential-geometry-%
theorem-Hadamard}. Both proofs, however,
show that a local diffeomorphism is a
covering map; compare Proposition
\ref{differential-geometry-%
proposition-local-diffeomorphism-covering}.
Alexandrov's theorem also appears as
Theorem
\ref{xxii-egd-2026-theorem-alexandrov58}
in the XXII Escola de Geometria
Diferencial notes. Rotation surfaces, the
setting of the catenoid in Session 3, are
generalized to spaces of constant
curvature in \cite{docarmo1983rotation}.

\medskip\noindent
{\bf Conventions.}
Throughout, $\Sigma\subset\mathbb{R}^3$ is
a connected regular surface. Regular
surfaces are embedded. The unit normal
is $N$, the Weingarten endomorphism is
$A_p=-dN_p$, and
$$
H=\frac{1}{2}\operatorname{tr}A,
\qquad
K=\det A.
$$
The Laplacian is
$\Delta=\operatorname{div}\circ\nabla$,
so $\lambda$ is an eigenvalue when
$\Delta f+\lambda f=0$. On a compact
$\Sigma$ eigenvalues are nonnegative:
integrate $\operatorname{div}(f\nabla f)=
|\nabla f|^2+f\Delta f$. On the plane
$f=e^x$ gives $\lambda=-1$. For the sphere
of radius $r$
centred at the origin with its inner
normal $N(p)=-p/r$ one gets
$A=\frac{1}{r}\operatorname{id}$, $H=1/r$
and $K=1/r^2$.

\medskip\noindent
{\bf Session 1.}
Intrinsic calculus on surfaces and first
applications.

\medskip\noindent
{\bf Basic notions.}
The Euclidean inner product restricts to
each tangent plane and gives the
\emph{first fundamental form}
$$
\left\langle{} \cdot,\cdot
\right\rangle{}_p:
T_p\Sigma\times T_p\Sigma\to\mathbb{R},
$$
a symmetric positive definite bilinear
form. Everything that depends only on it
belongs to the \emph{intrinsic geometry}
of $\Sigma$. Assume $\Sigma$ orientable
with Gauss map $N:\Sigma\to\mathbb{S}^2$.
Its image $N(\Sigma)$ is the
\emph{spherical image} of $\Sigma$; for a
plane it is a point, and for a compact
surface one expects it to be the whole
sphere (Proposition
\ref{proposition-alias-spherical-image}).

Since $|N|=1$, the vector $dN_p(v)$ is
tangent, and the \emph{Weingarten
endomorphism} (or shape operator)
$$
A_p:T_p\Sigma\to T_p\Sigma,
\qquad
A_p(v)=-dN_p(v),
$$
is self-adjoint. Its eigenvalues
$\kappa_1(p),\kappa_2(p)$ are the
principal curvatures. For a unit vector
$v$, the normal curvature in the direction
$v$ is
$$
\kappa(v)=
\left\langle{} A_p(v),v
\right\rangle{},
$$
and $\kappa_1,\kappa_2$ are its maximum
and minimum. A point is \emph{umbilical}
if $\kappa_1(p)=\kappa_2(p)$, that is, if
the surface bends equally in all
directions. The mean and Gauss curvatures
are
$$
H(p)=\frac{1}{2}\operatorname{tr}A_p
=\frac{\kappa_1(p)+\kappa_2(p)}{2},
\qquad
K(p)=\det A_p=\kappa_1(p)\kappa_2(p).
$$

\begin{lemma}
\label{lemma-alias-h2-geq-k}
For every $p\in\Sigma$ one has
$H^2(p)\geq K(p)$, with equality if and
only if $p$ is umbilical.
\end{lemma}

\begin{proof}
One has
$$
H^2(p)-K(p)
=\frac{(\kappa_1(p)-\kappa_2(p))^2}{4}
\geq 0.
$$
\end{proof}

\medskip\noindent
{\bf Covariant derivative.}
For $f\in C^\infty(\Sigma)$ and
$v\in T_p\Sigma$,
$$
v(f)=df_p(v)
=\frac{d}{dt}\Big|_{t=t_0}f(\alpha(t)),
$$
where $\alpha:I\to\Sigma$ is any curve
with $\alpha(t_0)=p$ and $\alpha'(t_0)=v$.
A tangent vector field is a smooth map
$X:\Sigma\to\mathbb{R}^3$ with
$X(p)\in T_p\Sigma$; the set of them is
$\mathfrak{X}(\Sigma)$. The derivative
$dX_p(v)\in\mathbb{R}^3$ need not be
tangent, so one keeps its tangent part.

\begin{definition}
\label{definition-alias-covariant-derivative}
The \emph{covariant} (or
\emph{intrinsic}) \emph{derivative} of
$X\in\mathfrak{X}(\Sigma)$ with respect to
$v\in T_p\Sigma$ is
$$
\nabla_vX=dX_p(v)
-\left\langle{} dX_p(v),N(p)
\right\rangle{}N(p)\in T_p\Sigma.
$$
For $X,Y\in\mathfrak{X}(\Sigma)$ and
$f\in C^\infty(\Sigma)$ put
$(Xf)(p)=X(p)(f)$ and
$(\nabla_XY)(p)=\nabla_{X(p)}Y$.
\end{definition}

Differentiating $\left\langle{} X,N
\right\rangle{}=0$ gives
$\left\langle{} dX_p(v),N(p)
\right\rangle{}=\left\langle{}
X(p),A_p(v)\right\rangle{}$, hence the
\emph{Gauss formula}
\begin{equation}
\label{equation-alias-gauss-formula}
dX_p(v)=\nabla_vX
+\left\langle{} A_p(v),X(p)
\right\rangle{}N(p).
\end{equation}
The covariant derivative is an affine
connection compatible with the metric:
for $X,Y,Z\in\mathfrak{X}(\Sigma)$ and
$f\in C^\infty(\Sigma)$,
$$
\begin{aligned}
\nabla_{X+Y}Z&=\nabla_XZ+\nabla_YZ,\\
\nabla_{fX}Y&=f\nabla_XY,\\
\nabla_X(Y+Z)&=\nabla_XY+\nabla_XZ,\\
\nabla_X(fY)&=X(f)Y+f\nabla_XY,
\end{aligned}
$$
and
$$
X\left\langle{} Y,Z\right\rangle{}
=\left\langle{} \nabla_XY,Z
\right\rangle{}
+\left\langle{} Y,\nabla_XZ
\right\rangle{}.
$$
It is moreover \emph{symmetric}. Indeed,
since $A_p$ is self-adjoint, the vectors
$dY_p(X(p))$ and $dX_p(Y(p))$ have the
same normal part, so the
\emph{Lie bracket}
$$
[X,Y](p)=dY_p(X(p))-dX_p(Y(p))
$$
is tangent. It satisfies
$[X,Y](f)=X(Y(f))-Y(X(f))$ and
$$
[X,Y]=\nabla_XY-\nabla_YX.
$$
See \cite[Section 1.2,
pp.~3--6]{alias2006analisis}.

\medskip\noindent
{\bf Differential operators. Gradient and
Hessian.}

\begin{definition}
\label{definition-alias-gradient}
Let $f\in C^\infty(\Sigma)$. The
\emph{gradient} of $f$ at $p$ is the
vector $\nabla f(p)\in T_p\Sigma$ such that
$$
\left\langle{} \nabla f(p),v
\right\rangle{}=df_p(v)
$$
for every $v\in T_p\Sigma$.
\end{definition}

Thus $\nabla f\in\mathfrak{X}(\Sigma)$ is
characterized by $\left\langle{} \nabla
f,X\right\rangle{}=X(f)$. For
$f,g\in C^\infty(\Sigma)$ and a smooth
$\phi:\mathbb{R}\to\mathbb{R}$,
$$
\begin{aligned}
\nabla(f+g)&=\nabla f+\nabla g,\\
\nabla(fg)&=g\nabla f+f\nabla g,\\
\nabla\phi(f)&=\phi'(f)\nabla f,
\end{aligned}
$$
and, $\Sigma$ being connected,
$\nabla f\equiv 0$ if and only if $f$ is
constant.

\begin{definition}
\label{definition-alias-hessian}
The \emph{Hessian} of
$f\in C^\infty(\Sigma)$ at $p$ is the map
$$
\nabla^2f_p:
T_p\Sigma\times T_p\Sigma\to\mathbb{R},
\qquad
\nabla^2f_p(v,w)=\left\langle{}
\nabla_v\nabla f,w\right\rangle{}.
$$
\end{definition}

By the Gauss formula
(\ref{equation-alias-gauss-formula}) with
$X=\nabla f$, the Hessian can be computed
with the Euclidean derivative of the
vector valued map $\nabla f$:
\begin{equation}
\label{equation-alias-hessian-euclidean}
\nabla^2f_p(v,w)=\left\langle{}
d(\nabla f)_p(v),w\right\rangle{}.
\end{equation}
The Hessian is bilinear. To see that it
is symmetric, one passes to the global
point of view and considers
$$
\nabla^2f:
\mathfrak{X}(\Sigma)\times
\mathfrak{X}(\Sigma)\to C^\infty(\Sigma),
\qquad
\nabla^2f(X,Y)=\left\langle{}
\nabla_X\nabla f,Y\right\rangle{}.
$$
It is $C^\infty(\Sigma)$-bilinear and
$\nabla^2f(X,Y)(p)=\nabla^2f_p(v,w)$
whenever $X(p)=v$ and $Y(p)=w$.

\begin{lemma}
\label{lemma-alias-hessian-symmetric}
The Hessian is symmetric.
\end{lemma}

\begin{proof}
By compatibility with the metric,
$$
\nabla^2f(X,Y)
=X(Y(f))-\left\langle{} \nabla f,
\nabla_XY\right\rangle{}.
$$
Hence
$$
\begin{aligned}
&\nabla^2f(X,Y)-\nabla^2f(Y,X)\\
&=[X,Y](f)-\left\langle{} \nabla f,
\nabla_XY-\nabla_YX\right\rangle{}=0,
\end{aligned}
$$
by symmetry of $\nabla$. The pointwise
statement follows from the global one.
\end{proof}

\medskip\noindent
{\bf Local extrema and critical points.}

\begin{definition}
\label{definition-alias-extrema}
A function $f:\Sigma\to\mathbb{R}$ attains
a \emph{local minimum} (resp.\
\emph{local maximum}) at $p_0\in\Sigma$ if
there is a neighbourhood $V$ of $p_0$ in
$\Sigma$ with $f(p)\geq f(p_0)$ (resp.\
$f(p)\leq f(p_0)$) for every $p\in V$.
The extremum is \emph{strict} if the
inequality is strict for $p\neq p_0$. It
is a \emph{global} extremum if one can
take $V=\Sigma$. A \emph{critical point}
of $f$ is a point $p_0$ with
$\nabla f(p_0)=0$.
\end{definition}

If $f$ attains a local extremum at
$p_0$, then $p_0$ is a critical point
(the \emph{first derivative test}).

\begin{theorem}[Second derivative test]
\label{theorem-alias-second-derivative-test}
Let $f\in C^\infty(\Sigma)$ and let $p_0$
be a critical point of $f$.
\begin{enumerate}
\item If $f$ attains a local minimum at
$p_0$, then $\nabla^2f_{p_0}(v,v)\geq 0$
for every $v\in T_{p_0}\Sigma$.
\item If $f$ attains a local maximum at
$p_0$, then $\nabla^2f_{p_0}(v,v)\leq 0$
for every $v\in T_{p_0}\Sigma$.
\item If $\nabla^2f_{p_0}(v,v)>0$ for
every $v\neq 0$, then $f$ attains a strict
local minimum at $p_0$.
\item If $\nabla^2f_{p_0}(v,v)<0$ for
every $v\neq 0$, then $f$ attains a strict
local maximum at $p_0$.
\end{enumerate}
\end{theorem}

\begin{proof}
The booklet leaves this as an exercise
\cite[Theorem 1.4, p.~9]{alias2006analisis}.
Let $x:U\to\Sigma$ be a parametrization
with $x(u_0)=p_0$ and put $F=f\circ x$,
$x_i=\partial x/\partial u_i$. Then
$\partial F/\partial u_j=\left\langle{}
\nabla f\circ x,x_j\right\rangle{}$, so
$$
\begin{aligned}
\frac{\partial^2F}{\partial u_i
\partial u_j}
&=\left\langle{}
d(\nabla f)(x_i),x_j\right\rangle{}
+\left\langle{} \nabla f\circ x,
x_{ij}\right\rangle{}.
\end{aligned}
$$
At $u_0$ the last term vanishes because
$\nabla f(p_0)=0$. By
(\ref{equation-alias-hessian-euclidean}),
the matrix of second partial derivatives
of $F$ at the critical point $u_0$ is the
matrix of $\nabla^2f_{p_0}$ in the basis
$x_1(u_0),x_2(u_0)$. Since $x$ is a
homeomorphism onto an open set, local
(strict) extrema of $f$ at $p_0$ and of
$F$ at $u_0$ correspond. The statement
follows from the second derivative test
in $\mathbb{R}^2$.
\end{proof}

Equivalently, at a critical point $p_0$
and for any curve $\alpha$ with
$\alpha(t_0)=p_0$ and $\alpha'(t_0)=v$,
$$
\nabla^2f_{p_0}(v,v)
=\frac{d^2}{dt^2}\Big|_{t=t_0}
f(\alpha(t)).
$$

\medskip\noindent
{\bf Two examples.}
If $f$ is the restriction to $\Sigma$ of a
smooth $F$ defined on an open set of
$\mathbb{R}^3$, then $\nabla f$ is the
tangent part of the Euclidean gradient,
and with $u=\left\langle{}
\operatorname{Grad}F,N\right\rangle{}$,
\begin{equation}
\label{equation-alias-restriction-hessian}
\nabla^2f_p(v,w)
=\operatorname{Hess}F_p(v,w)
+\left\langle{} A_p(v),w
\right\rangle{}u(p)
\end{equation}
\cite[Example 1.5,
p.~10]{alias2006analisis}. Here and below
$v^{\top}$ denotes the tangent part of a
vector $v\in\mathbb{R}^3$.

\begin{example}[Squared distance]
\label{example-alias-squared-distance}
Fix $c\in\mathbb{R}^3$ and let
$$
f(p)=\frac{1}{2}|p-c|^2.
$$
Then $df_p(v)=\left\langle{} p-c,v
\right\rangle{}=\left\langle{}
(p-c)^{\top},v\right\rangle{}$, so
\begin{equation}
\label{equation-alias-grad-distance}
\nabla f(p)=(p-c)^{\top}
=p-c-g(p)N(p),
\end{equation}
where
$g(p)=\left\langle{} p-c,N(p)
\right\rangle{}$. Differentiating,
$$
d(\nabla f)_p(v)=v-v(g)N(p)
+g(p)A_p(v),
$$
whose tangent part gives
\begin{equation}
\label{equation-alias-hess-distance}
\nabla_v\nabla f=v+g(p)A_p(v),
\end{equation}
that is,
$$
\nabla^2f_p(v,w)=\left\langle{} v,w
\right\rangle{}+\left\langle{} A_p(v),w
\right\rangle{}\left\langle{} p-c,N(p)
\right\rangle{}.
$$
In particular, the critical points of
$f$ are the points $p$ where $p-c$ is
normal to $\Sigma$
\cite[Example 1.6,
p.~11]{alias2006analisis}.
\end{example}

\begin{theorem}
\label{theorem-alias-elliptic-point}
Every compact orientable regular surface
$\Sigma\subset\mathbb{R}^3$ has at least
one \emph{elliptic point}, that is, a
point where $K>0$.
\end{theorem}

\begin{proof}
The idea is to enclose $\Sigma$ in a large
sphere and shrink it until it first
touches $\Sigma$; at the contact point the
normal curvatures are bounded by those of
the sphere. Analytically
\cite[Theorem 2.1, p.~15]{alias2006analisis},
let $f(p)=\frac{1}{2}|p|^2$ and let $p_0$
be a global maximum of $f$. Then
$\nabla f(p_0)=0$, so by
(\ref{equation-alias-grad-distance}) with
$c=0$,
$$
N(p_0)=\varepsilon\frac{p_0}{|p_0|},
\qquad \varepsilon=\pm 1,
$$
where $|p_0|>0$. Moreover
$\nabla^2f_{p_0}\leq 0$, so for every
unit vector $v\in T_{p_0}\Sigma$,
$$
1+\varepsilon\kappa(v)|p_0|\leq 0.
$$
If $\varepsilon=1$, then
$\kappa(v)\leq -1/|p_0|<0$ for all $v$; if
$\varepsilon=-1$, then
$\kappa(v)\geq 1/|p_0|>0$ for all $v$.
In both cases $\kappa_1(p_0)$ and
$\kappa_2(p_0)$ have the same sign and
$K(p_0)>0$.
\end{proof}

\begin{remark}
\label{remark-alias-compact-orientable}
Every compact regular surface
$\Sigma\subset\mathbb{R}^3$ is orientable.
By the Jordan--Brouwer separation theorem,
$\Sigma=\partial\Omega$ for a compact
regular domain $\Omega\subset\mathbb{R}^3$,
and $\Sigma$ carries the \emph{inner unit
normal}: the half-line from $p$ in the
direction $N(p)$ starts inside $\Omega$
\cite[Remark 2.2, p.~16]{alias2006analisis}.
With the inner normal, $\varepsilon=-1$ in
the proof above: $\Omega$ lies in the
closed ball of radius $|p_0|$ about $0$ (by
the open and closed argument in the proof
of Proposition
\ref{proposition-alias-spherical-image}),
so the inner normal at $p_0$ is
$-p_0/|p_0|$. Hence all normal curvatures
at $p_0$ are at least $1/|p_0|$, the
normal curvature of the sphere of radius
$|p_0|$ centred at $0$.
\end{remark}

\begin{example}[Height function]
\label{example-alias-height-function}
Let $\Pi$ be the affine plane through
$c\in\mathbb{R}^3$ with unit normal $a$.
The \emph{height function}
$$
h(p)=\left\langle{} p-c,a
\right\rangle{}
$$
is the oriented distance from $p$ to
$\Pi$. Since the Euclidean gradient of
$x\mapsto\left\langle{} x-c,a
\right\rangle{}$ is the constant $a$ and
its Euclidean Hessian vanishes,
(\ref{equation-alias-restriction-hessian})
gives
\begin{equation}
\label{equation-alias-grad-height}
\nabla h=a^{\top}
=a-\left\langle{} a,N\right\rangle{}N
\end{equation}
and
\begin{equation}
\label{equation-alias-hess-height}
\nabla^2h_p(v,w)=\left\langle{} A_p(v),w
\right\rangle{}\left\langle{} a,N(p)
\right\rangle{}.
\end{equation}
The critical points of $h$ are the points
where $N(p)=\pm a$, that is, where the
tangent plane is parallel to $\Pi$
\cite[Example 1.7,
pp.~11--12]{alias2006analisis}.
\end{example}

\begin{proposition}
\label{proposition-alias-spherical-image}
If $\Sigma\subset\mathbb{R}^3$ is a compact
connected regular surface with Gauss map
$N$, then $N(\Sigma)=\mathbb{S}^2$.
\end{proposition}

\begin{proof}
Fix $a\in\mathbb{S}^2$ and let
$h(p)=\left\langle{} p,a\right\rangle{}$.
By compactness $h$ attains its maximum at
some $p_1$ and its minimum at some $p_2$.
These are critical points, so
$N(p_i)=\pm a$. It remains to see that
$N(p_i)=a$ for some $i$.

Write $\Sigma=\partial\Omega$ as in Remark
\ref{remark-alias-compact-orientable} and
let $m=h(p_2)$. The set of points of
$\Omega$ in the open half-space
$\{\left\langle{} x,a\right\rangle{}<m\}$
is open and closed in that half-space,
because $\partial\Omega=\Sigma$ does not
meet it. The half-space is connected and
unbounded, and $\Omega$ is bounded, so
this set is empty. Thus $\Omega$ lies in
$\{\left\langle{} x,a\right\rangle{}\geq
m\}$, and the inner normal at $p_2$ is
$a$. In the same way the inner normal at
$p_1$ is $-a$. Since $\Sigma$ is
connected, $N$ is either the inner or the
outer normal, and in both cases one of
$N(p_1)$, $N(p_2)$ equals $a$.
\end{proof}

The lecture stated this proposition with
a sketch; the role of Jordan--Brouwer in
the last step is an addition here.

\medskip\noindent
{\bf Session 2.}
Convex surfaces, Hadamard's theorem, and
the Minkowski formulas.

\medskip\noindent
{\bf Convex surfaces.}
Let $\Sigma$ be oriented by $N$. For
$p\in\Sigma$ the affine tangent plane
$$
\Pi_p=\{x\in\mathbb{R}^3:
\left\langle{} x-p,N(p)\right\rangle{}=0\}
$$
divides $\mathbb{R}^3$ into the closed
half-spaces
$$
\begin{aligned}
\Pi_p^+&=\{x:\left\langle{} x-p,N(p)
\right\rangle{}\geq 0\},\\
\Pi_p^-&=\{x:\left\langle{} x-p,N(p)
\right\rangle{}\leq 0\}.
\end{aligned}
$$
The half-space $\Pi_p^+$ is the one
into which $N(p)$ points.

\begin{definition}
\label{definition-alias-convexity}
Let $p_0\in\Sigma$.
\begin{enumerate}
\item $\Sigma$ is \emph{locally convex} at
$p_0$ if some neighbourhood $V$ of $p_0$
in $\Sigma$ lies in $\Pi_{p_0}^+$ or in
$\Pi_{p_0}^-$. It is \emph{strictly
locally convex} at $p_0$ if moreover
$V\setminus\{p_0\}$ lies in the interior
of that half-space.
\item $\Sigma$ is \emph{globally convex}
at $p_0$ if all of $\Sigma$ lies in
$\Pi_{p_0}^+$ or in $\Pi_{p_0}^-$. It is
\emph{strictly globally convex} at $p_0$
if moreover $\Sigma\setminus\{p_0\}$ lies
in the interior of that half-space.
\end{enumerate}
\end{definition}

To relate convexity and curvature,
consider the height function with respect
to the affine tangent plane at $p_0$,
$$
h_{p_0}(p)=\left\langle{} p-p_0,N(p_0)
\right\rangle{}.
$$
It vanishes on $\Sigma\cap\Pi_{p_0}$, is
positive on $\Sigma\cap
\operatorname{int}\Pi_{p_0}^+$ and
negative on $\Sigma\cap
\operatorname{int}\Pi_{p_0}^-$. So
$\Sigma$ is (strictly) locally convex at
$p_0$ if and only if $h_{p_0}$ has a
(strict) local extremum at $p_0$, and
(strictly) globally convex at $p_0$ if
and only if $h_{p_0}$ has a (strict)
global extremum at $p_0$. By
(\ref{equation-alias-grad-height}), $p_0$
is always a critical point of $h_{p_0}$,
and by (\ref{equation-alias-hess-height})
$$
\nabla^2(h_{p_0})_{p_0}(v,w)
=\left\langle{} A_{p_0}(v),w
\right\rangle{}.
$$

\begin{proposition}[Convexity and
curvature]
\label{proposition-alias-convexity-curvature}
Let $\Sigma$ be an oriented regular
surface and $p_0\in\Sigma$.
\begin{enumerate}
\item If $\Sigma$ is locally convex at
$p_0$, then $K(p_0)\geq 0$.
\item If $K(p_0)>0$, then $\Sigma$ is
strictly locally convex at $p_0$.
\end{enumerate}
\end{proposition}

\begin{proof}
See \cite[Proposition 2.5,
pp.~17--18]{alias2006analisis}.
If $\Sigma$ is
locally convex at $p_0$, then $h_{p_0}$
has a local extremum there, say a minimum
(otherwise change $N$). By Theorem
\ref{theorem-alias-second-derivative-test},
$\left\langle{} A_{p_0}(v),v
\right\rangle{}\geq 0$ for all $v$. In a
basis $e_1,e_2$ of principal directions
this gives $\kappa_i(p_0)\geq 0$, so
$K(p_0)\geq 0$.

If $K(p_0)>0$, then after possibly
changing $N$ both principal curvatures are
positive, so
$$
\nabla^2(h_{p_0})_{p_0}(v,v)
=\kappa(v)>0
$$
for every unit $v$. By Theorem
\ref{theorem-alias-second-derivative-test},
$h_{p_0}$ has a strict local minimum at
$p_0$.
\end{proof}

\begin{definition}
\label{definition-alias-ovaloid}
An \emph{ovaloid} is a compact connected
regular surface whose Gauss curvature
never vanishes.
\end{definition}

By Remark
\ref{remark-alias-compact-orientable} an
ovaloid is orientable, and by Theorem
\ref{theorem-alias-elliptic-point} and
connectedness its Gauss curvature is
positive everywhere: all its points are
elliptic. By Proposition
\ref{proposition-alias-convexity-curvature}
an ovaloid is strictly locally convex at
every point. Hadamard's theorem upgrades
this to global convexity.

\begin{theorem}[Hadamard]
\label{theorem-alias-hadamard}
Let $\Sigma\subset\mathbb{R}^3$ be an
ovaloid. Then
\begin{enumerate}
\item the Gauss map
$N:\Sigma\to\mathbb{S}^2$ is a
diffeomorphism, and
\item $\Sigma$ is strictly globally convex
at every point.
\end{enumerate}
\end{theorem}

\begin{proof}
We follow \cite[Theorem 2.6,
pp.~19--20]{alias2006analisis}.

(1) Since $\det(dN_p)=\det A_p=K(p)>0$,
$dN_p$ is invertible at every point and
$N$ is a local diffeomorphism. It remains
to show that $N$ is bijective.

Surjectivity: $N(\Sigma)$ is compact,
hence closed in $\mathbb{S}^2$, and open
because $N$ is a local diffeomorphism.
Since $\mathbb{S}^2$ is connected,
$N(\Sigma)=\mathbb{S}^2$.

Injectivity: since $\mathbb{S}^2$ is
simply connected, every covering map
onto it from a connected space is a
homeomorphism, so it suffices to show
that $N$ is a covering map. Fix
$a\in\mathbb{S}^2$. The fibre $N^{-1}(a)$
is finite: otherwise, by compactness, it
would contain a sequence of distinct
points converging to some $p_0\in
N^{-1}(a)$, and $N$ would not be injective
near $p_0$, contradicting that $N$ is a
local diffeomorphism. Write
$N^{-1}(a)=\{p_1,\dots,p_k\}$ and choose
pairwise disjoint neighbourhoods
$\tilde{U}_i$ of $p_i$ such that
$N|_{\tilde{U}_i}$ is a diffeomorphism
onto a neighbourhood $\tilde{V}_i$ of
$a$. Put
$$
V=\bigcap_{i=1}^k\tilde{V}_i,
\qquad
U_i=(N|_{\tilde{U}_i})^{-1}(V).
$$
Then each $N|_{U_i}:U_i\to V$ is a
diffeomorphism. After shrinking $V$ (and
the $U_i$ accordingly), $N^{-1}(V)$ is the
disjoint union of the $U_i$: otherwise
there are
points $q_n\notin\bigcup_iU_i$ with
$N(q_n)\to a$, and a limit point $q$ of
them lies in $N^{-1}(a)$ but outside the
open set $\bigcup_iU_i$, which is
impossible. So $N$ is a covering map.

(2) Fix $p_0\in\Sigma$. The critical
points of $h_{p_0}$ are the points with
$N(p)=\pm N(p_0)$. Since $N$ is
bijective, $h_{p_0}$ has exactly two
critical points, one of which is $p_0$.
By compactness, its global maximum and
minimum are attained at critical points,
and they are strict (otherwise there would
be more than two critical points). Hence
$h_{p_0}$ has a strict global extremum at
$p_0$, and $\Sigma$ is strictly globally
convex at $p_0$.
\end{proof}

The ``shrinking'' step in the covering
argument is spelled out here; the booklet
takes $V$ as the intersection of the
$\tilde{V}_i$ directly.

\medskip\noindent
{\bf Divergence and Laplacian.}

\begin{definition}
\label{definition-alias-divergence}
The \emph{divergence} of
$X\in\mathfrak{X}(\Sigma)$ at $p$ is the
trace of the linear map
$$
(\nabla X)_p:T_p\Sigma\to T_p\Sigma,
\qquad
v\mapsto\nabla_vX.
$$
That is, for an orthonormal basis
$e_1,e_2$ of $T_p\Sigma$,
$$
\operatorname{div}X(p)
=\left\langle{} \nabla_{e_1}X,e_1
\right\rangle{}
+\left\langle{} \nabla_{e_2}X,e_2
\right\rangle{}.
$$
The \emph{Laplacian} of
$f\in C^\infty(\Sigma)$ is
$\Delta f=\operatorname{div}(\nabla f)$,
i.e.\ the trace of the Hessian
$\nabla^2f_p$.
\end{definition}

By the Gauss formula the divergence can
also be computed with the Euclidean
derivative:
$\operatorname{div}X(p)=\sum_i
\left\langle{} dX_p(e_i),e_i
\right\rangle{}$.

The following rules are direct from the
definitions, using
$\nabla_v(fX)=v(f)X+f\nabla_vX$ and
taking traces.

\begin{proposition}
\label{proposition-alias-div-laplacian}
For $X,Y\in\mathfrak{X}(\Sigma)$,
$f,g\in C^\infty(\Sigma)$,
$\lambda\in\mathbb{R}$ and a smooth
$\phi:\mathbb{R}\to\mathbb{R}$,
$$
\begin{aligned}
\operatorname{div}(X+Y)
&=\operatorname{div}X
+\operatorname{div}Y,\\
\operatorname{div}(fX)
&=\left\langle{} \nabla f,X
\right\rangle{}+f\operatorname{div}X,
\end{aligned}
$$
and
$$
\begin{aligned}
\Delta(f+g)&=\Delta f+\Delta g,\\
\Delta(\lambda f)&=\lambda\Delta f,\\
\Delta\phi(f)&=\phi'(f)\Delta f
+\phi''(f)|\nabla f|^2,\\
\Delta(fg)&=f\Delta g+g\Delta f
+2\left\langle{} \nabla f,\nabla g
\right\rangle{}.
\end{aligned}
$$
\end{proposition}

\begin{theorem}[Divergence theorem]
\label{theorem-alias-divergence}
Let $\Sigma\subset\mathbb{R}^3$ be a
compact oriented regular surface and
$X\in\mathfrak{X}(\Sigma)$. Then
$$
\int_\Sigma\operatorname{div}X(p)\,dp=0,
$$
where $dp$ is the area element of
$\Sigma$. In particular
$\int_\Sigma\Delta f\,dp=0$ for every
$f\in C^\infty(\Sigma)$.
\end{theorem}

See \cite[Section 3.1,
pp.~21--22]{alias2006analisis}.

\medskip\noindent
{\bf The Minkowski formulas.}

\begin{theorem}[Minkowski formulas]
\label{theorem-alias-minkowski}
Let $\Sigma\subset\mathbb{R}^3$ be a
compact oriented regular surface with
Gauss map $N$, and fix
$c\in\mathbb{R}^3$. Then
\begin{equation}
\label{equation-alias-minkowski-1}
\int_\Sigma\left(1+H(p)\left\langle{}
p-c,N(p)\right\rangle{}\right)dp=0
\end{equation}
and
\begin{equation}
\label{equation-alias-minkowski-2}
\int_\Sigma\left(H(p)+K(p)\left\langle{}
p-c,N(p)\right\rangle{}\right)dp=0.
\end{equation}
\end{theorem}

These are the \emph{first} and
\emph{second Minkowski formulas}
\cite[Theorem 3.3,
p.~23]{alias2006analisis}. As a check,
for the sphere of radius $r$ centred at
$c$ with inner normal one has
$\left\langle{} p-c,N\right\rangle{}=-r$,
$H=1/r$ and $K=1/r^2$, and both
integrands vanish identically.

\begin{proof}[Proof of the first formula]
Let $f(p)=\frac{1}{2}|p-c|^2$ and
$g(p)=\left\langle{} p-c,N(p)
\right\rangle{}$. Taking the trace of
(\ref{equation-alias-hess-distance}) in a
basis of principal directions,
\begin{equation}
\label{equation-alias-laplacian-f}
\Delta f
=(1+\kappa_1g)+(1+\kappa_2g)
=2(1+Hg).
\end{equation}
Integrate and apply Theorem
\ref{theorem-alias-divergence}.
\end{proof}

The second formula needs the
derivative of the Weingarten
endomorphism, which the lecture
introduced next.

\medskip\noindent
{\bf Appendix: covariant derivative of a
tensor field.}
A \emph{tensor field} (or field of
endomorphisms) on $\Sigma$ assigns to
each $p$ an endomorphism $S_p$ of
$T_p\Sigma$. It is \emph{smooth} if for
every $X\in\mathfrak{X}(\Sigma)$ the field
$SX$, $(SX)(p)=S_p(X(p))$, is smooth. For
example, $A$ is a smooth tensor field.
Such an $S$ gives a map
$S:\mathfrak{X}(\Sigma)\to
\mathfrak{X}(\Sigma)$ that is additive
and $C^\infty(\Sigma)$-linear; conversely,
every such map comes from a unique smooth
tensor field
\cite[Section 1.4,
p.~12]{alias2006analisis}.

\begin{definition}
\label{definition-alias-tensor-derivative}
The \emph{covariant derivative} of a
smooth tensor field $S$ with respect to
$X\in\mathfrak{X}(\Sigma)$ is the tensor
field $\nabla_XS$ given by
$$
(\nabla_XS)(Y)=\nabla_X(SY)
-S(\nabla_XY)
$$
for $Y\in\mathfrak{X}(\Sigma)$.
\end{definition}

The map $\nabla_XS$ is additive and
$C^\infty(\Sigma)$-linear in $Y$, so it is
a tensor field. Its value at $p$ depends
only on $X(p)=v$, which defines
$\nabla_vS:T_p\Sigma\to T_p\Sigma$
\cite[Definition 1.8,
p.~13]{alias2006analisis}.

\begin{proposition}
\label{proposition-alias-trace-commutes}
For a smooth tensor field $S$,
$p\in\Sigma$ and $v\in T_p\Sigma$,
$$
\operatorname{tr}(\nabla_vS)
=v(\operatorname{tr}S)
=\left\langle{} v,
\nabla(\operatorname{tr}S)(p)
\right\rangle{}.
$$
\end{proposition}

This is \cite[Proposition 1.9,
p.~13]{alias2006analisis}, where the
proof is left as an exercise (hint: use
coordinates).

\begin{theorem}[Codazzi--Mainardi]
\label{theorem-alias-codazzi}
The Codazzi--Mainardi equations of a
regular surface $\Sigma$ are equivalent to
$$
(\nabla_XA)(Y)=(\nabla_YA)(X)
$$
for all $X,Y\in\mathfrak{X}(\Sigma)$.
Equivalently, for every $p\in\Sigma$ and
all $v,w\in T_p\Sigma$,
$$
(\nabla_vA)(w)=(\nabla_wA)(v).
$$
\end{theorem}

See \cite[Theorem 1.10,
pp.~13--14]{alias2006analisis}; the proof
is left there as an exercise.

\begin{proof}[Proof of the second
Minkowski formula]
We follow
\cite[pp.~23--25]{alias2006analisis}.
With $f$ and $g$ as above, for
$v\in T_p\Sigma$,
$$
\begin{aligned}
v(g)&=\left\langle{} v,N(p)
\right\rangle{}+\left\langle{} p-c,
dN_p(v)\right\rangle{}\\
&=-\left\langle{} A_p((p-c)^{\top}),v
\right\rangle{},
\end{aligned}
$$
so that
$$
\nabla g=-A(\nabla f).
$$
Differentiating, and using
Codazzi--Mainardi to exchange $v$ and
$\nabla f(p)$,
$$
\begin{aligned}
\nabla_v\nabla g
&=-(\nabla_vA)(\nabla f)
-A(\nabla_v\nabla f)\\
&=-(\nabla_{\nabla f(p)}A)(v)
-A_p(v)-g(p)A_p^2(v),
\end{aligned}
$$
by (\ref{equation-alias-hess-distance}).
Taking traces and using Proposition
\ref{proposition-alias-trace-commutes},
$$
\Delta g=-\left\langle{} \nabla f,
\nabla\operatorname{tr}A\right\rangle{}
-\operatorname{tr}A
-g\operatorname{tr}(A^2).
$$
Since $\operatorname{tr}A=2H$ and
$\operatorname{tr}(A^2)
=\kappa_1^2+\kappa_2^2=4H^2-2K$,
\begin{equation}
\label{equation-alias-laplacian-g}
\Delta g=-2\left\langle{} \nabla f,
\nabla H\right\rangle{}-2H
-2(2H^2-K)g.
\end{equation}
On the other hand, by
(\ref{equation-alias-laplacian-f}),
$$
\begin{aligned}
\operatorname{div}(2H\nabla f)
&=2\left\langle{} \nabla H,\nabla f
\right\rangle{}+2H\Delta f\\
&=2\left\langle{} \nabla H,\nabla f
\right\rangle{}+4H(1+Hg).
\end{aligned}
$$
Adding,
$$
\operatorname{div}(2H\nabla f)+\Delta g
=2(H+Kg).
$$
Integrating and applying Theorem
\ref{theorem-alias-divergence} gives
(\ref{equation-alias-minkowski-2}).
\end{proof}

\begin{remark}
\label{remark-alias-parallel-surfaces}
Following Hsiung, the booklet gives a
second proof of
(\ref{equation-alias-minkowski-2}) from
(\ref{equation-alias-minkowski-1}) with
parallel surfaces
$\Sigma_t=\{p+tN(p):p\in\Sigma\}$, which
are regular for small $|t|$
\cite[Section 3.3,
pp.~26--29]{alias2006analisis}. The map
$\phi_t(p)=p+tN(p)$ has
$d\phi_t=\operatorname{id}-tA$ and
Jacobian $1-2tH+t^2K$, and $\Sigma_t$ has
principal curvatures
$\kappa_i/(1-t\kappa_i)$ and mean
curvature
$$
H_t\circ\phi_t
=\frac{H-tK}{1-2tH+t^2K}.
$$
Writing the first Minkowski formula for
$\Sigma_t$ as an integral over $\Sigma$
gives
$$
\int_\Sigma(1+Hg)\,dp
-t\int_\Sigma(H+Kg)\,dp=0
$$
for all small $t$, which yields both
formulas.
\end{remark}

\begin{theorem}[Liebmann, 1899]
\label{theorem-alias-liebmann}
A compact connected regular surface
$\Sigma\subset\mathbb{R}^3$ with constant
Gauss curvature $K$ is a sphere of radius
$1/\sqrt{K}$.
\end{theorem}

The proof was given in Session 3.

\medskip\noindent
{\bf Session 3.}
Minimal graphs and Bernstein-type
theorems. This lecture is not covered by
the booklet \cite{alias2006analisis}.

\begin{theorem}[Bernstein]
\label{theorem-alias-bernstein}
The only entire minimal graphs in
$\mathbb{R}^3$ are planes. In other
words, the only solutions
$u:\mathbb{R}^2\to\mathbb{R}$ of the
minimal surface equation
$$
\operatorname{Div}\left(
\frac{Du}{\sqrt{1+|Du|^2}}
\right)=0
$$
on the whole plane are the affine
functions $u(x,y)=ax+by+c$.
\end{theorem}

Here $\operatorname{Div}$ and $D$ are
the Euclidean divergence and gradient in
$\mathbb{R}^2$. Bernstein proved this
with PDE methods (the notes say in 1910;
the published proof is usually dated
1915--1917 (not checked)). The lecture
presented a geometric proof attributed to
Chern (1969).

\begin{definition}
\label{definition-alias-minimal}
An oriented regular surface is
\emph{minimal} if $H=0$ everywhere.
\end{definition}

If $H=0$ then $\kappa_2=-\kappa_1$, so
$K=-\kappa_1^2\leq 0$, with $K(p)=0$ if
and only if $A_p=0$. Moreover
$A^2=-K\operatorname{id}$ and
$|A|^2=\kappa_1^2+\kappa_2^2=-2K$.

\medskip\noindent
{\bf Variational characterization.}
For $f\in C_0^\infty(\Sigma)$, the
\emph{normal variation} of $\Sigma$
induced by $f$ is
$$
\Sigma_t=\{p+tf(p)N(p):p\in\Sigma\}
\subset\mathbb{R}^3.
$$
For $|t|<\varepsilon$ small, $\Sigma_t$ is
a regular surface and $\Sigma_0=\Sigma$.
Let $\mathcal{A}(t)$ be the area of
the part of $\Sigma_t$ over the support of
$f$. The \emph{first variation formula}
for the area is
$$
\delta_f\mathcal{A}
=\mathcal{A}'(0)
=-2\int_\Sigma fH\,dp.
$$
On a compact surface and for $f=1$ this
agrees with the Jacobian $1-2tH+t^2K$ of
Remark
\ref{remark-alias-parallel-surfaces}.
Hence $\Sigma$ is minimal if and only if
$\delta_f\mathcal{A}=0$ for every
$f\in C_0^\infty(\Sigma)$.

\medskip\noindent
{\bf Examples.}
The plane is the only minimal surface with
$K\equiv 0$: then $A\equiv 0$, and a
connected totally geodesic surface is an
open piece of a plane. The
\emph{catenoid} is the only non-planar
minimal surface of revolution, and the
\emph{helicoid} is the only non-planar
ruled minimal surface (both up to rigid
motions and homotheties). The helicoid and
the catenoid are locally isometric.

\medskip\noindent
{\bf Minimal graphs.}
Let $\Omega\subset\mathbb{R}^2$ be a
connected open set and
$u\in C^\infty(\Omega)$. The graph of $u$
is called a \emph{graph}; it is
\emph{entire} if $\Omega=\mathbb{R}^2$.
With the upward unit normal
$$
N=\frac{(-Du,1)}{\sqrt{1+|Du|^2}},
$$
the third coordinate of $N$ is positive,
so the spherical image of a graph lies in
the open upper hemisphere. The mean
curvature is given by
$$
2H=\operatorname{Div}\left(
\frac{Du}{\sqrt{1+|Du|^2}}
\right),
$$
so a graph is minimal if and only if $u$
solves the minimal surface equation.
(With $u=(x^2+y^2)/2$ both sides equal
$2$ at the origin.) Affine functions are
solutions, corresponding to planes. The
first non-trivial solution was found by
Scherk in 1835, looking for solutions of
the form $u(x,y)=f(x)+g(y)$:
$$
u(x,y)=\frac{1}{a}\log
\frac{\cos(ax)}{\cos(ay)},
$$
defined where
$\cos(ax)/\cos(ay)>0$.

Indeed, for $f=\log\cos x$ and
$g=-\log\cos y$ one has
$f''=-(1+f'^2)$ and $g''=1+g'^2$, and the
minimal surface equation in
non-divergence form
$$
(1+u_y^2)u_{xx}-2u_xu_yu_{xy}
+(1+u_x^2)u_{yy}=0
$$
holds. The factor $1/a$ is a homothety.

\medskip\noindent
{\bf Parabolicity.}

\begin{definition}
\label{definition-alias-parabolic}
A function $f$ on a Riemannian surface is
\emph{harmonic} if $\Delta f=0$,
\emph{subharmonic} if $\Delta f\geq 0$,
and \emph{superharmonic} if
$\Delta f\leq 0$. The surface is
\emph{parabolic} if the only subharmonic
functions bounded above are the
constants.
\end{definition}

Applying the definition to $-f$, on a
parabolic surface the only superharmonic
functions bounded below are the
constants. Parabolicity in dimension 2
behaves qualitatively differently from
higher dimensions; see below.

Every compact surface is parabolic. If
$\Delta f\geq 0$ on a compact $\Sigma$,
then $\int_\Sigma\Delta f\,dp=0$ forces
$\Delta f=0$, and then
$$
\int_\Sigma|\nabla f|^2dp
=\int_\Sigma\operatorname{div}(f\nabla f)
\,dp=0,
$$
so $f$ is constant
\cite[Lemma 5.2, p.~46]{alias2006analisis}.

\begin{theorem}[Ahlfors,
Blanc--Fiala--Huber]
\label{theorem-alias-ahlfors-huber}
Every complete Riemannian surface
$(\Sigma,g)$ with $K\geq 0$ is parabolic.
\end{theorem}

This was stated without proof (not
checked). In particular every complete
flat surface is parabolic.

\begin{proposition}
\label{proposition-alias-minimal-normal}
Let $\Sigma\subset\mathbb{R}^3$ be an
orientable minimal surface, let
$a\in\mathbb{R}^3$ with $|a|=1$, and put
$\varphi=\left\langle{} N,a
\right\rangle{}$. Then
$\nabla\varphi=-A(a^{\top})$,
$$
|\nabla\varphi|^2=K(\varphi^2-1),
\qquad
\Delta\varphi=2K\varphi.
$$
\end{proposition}

\begin{proof}
This was an exercise. For $v$ tangent,
$v(\varphi)=\left\langle{} dN(v),a
\right\rangle{}=-\left\langle{} A(v),
a^{\top}\right\rangle{}$, so
$\nabla\varphi=-A(a^{\top})$. Since
$A^2=-K\operatorname{id}$ and
$|a^{\top}|^2=1-\varphi^2$,
$$
|\nabla\varphi|^2=\left\langle{}
A^2a^{\top},a^{\top}\right\rangle{}
=-K(1-\varphi^2).
$$
Next,
$a^{\top}=a-\varphi N$ has
$d(a^{\top})(v)=-v(\varphi)N+\varphi A(v)$,
so $\nabla_va^{\top}=\varphi A(v)$ and,
by Codazzi--Mainardi,
$$
\nabla_v\nabla\varphi
=-(\nabla_{a^{\top}}A)(v)
-\varphi A^2(v).
$$
Taking traces with Proposition
\ref{proposition-alias-trace-commutes},
$$
\Delta\varphi
=-2\left\langle{} \nabla H,a
\right\rangle{}-|A|^2\varphi,
$$
which equals $2K\varphi$ when $H=0$.
\end{proof}

\begin{theorem}[Chern; geometric
Bernstein]
\label{theorem-alias-chern}
The only complete orientable minimal
surfaces in $\mathbb{R}^3$ whose spherical
image lies in an open hemisphere are the
planes.
\end{theorem}

Theorem \ref{theorem-alias-bernstein}
follows: the graph of an entire $u$ is
complete, and its spherical image lies in
the open upper hemisphere.

\begin{proof}
Let the hemisphere be centred at $a$, so
that $\varphi=\left\langle{} N,a
\right\rangle{}>0$. By Proposition
\ref{proposition-alias-minimal-normal},
$\Delta\varphi=2K\varphi\leq 0$, so
$\varphi$ is a positive superharmonic
function. If $(\Sigma,g)$ is parabolic,
$\varphi$ is constant, so $K\varphi=0$,
$K\equiv 0$, and $\Sigma$ is a plane.

To prove parabolicity, change the metric
conformally. A metric conformal to $g$ is
$\hat{g}=e^{2u}g$ with $u$ smooth; for
example, the Poincar\'e metric on the unit
disc is
$$
\hat{g}=\frac{4}{(1-x^2-y^2)^2}
(dx^2+dy^2).
$$
In dimension 2,
$$
\hat{\Delta}=e^{-2u}\Delta,
\qquad
\hat{K}=e^{-2u}(K-\Delta u).
$$
By the first identity, subharmonic
functions for $g$ and $\hat{g}$ coincide,
so parabolicity is a conformal
invariant in dimension 2.

Take
$$
\hat{g}=(1+\varphi)^2g,
\qquad
u=\log(1+\varphi).
$$
Since $\varphi>0$, $\hat{g}\geq g$. So
$\hat{g}$ is complete: a divergent curve
$\alpha$ has
$$
\hat{L}(\alpha)=\int\sqrt{\hat{g}
(\alpha',\alpha')}\,dt
\geq\int\sqrt{g(\alpha',\alpha')}\,dt
=\infty.
$$
By Proposition
\ref{proposition-alias-minimal-normal}
and Proposition
\ref{proposition-alias-div-laplacian},
$$
\begin{aligned}
\Delta u
&=\frac{\Delta\varphi}{1+\varphi}
-\frac{|\nabla\varphi|^2}
{(1+\varphi)^2}\\
&=\frac{2K\varphi}{1+\varphi}
-\frac{K(\varphi-1)}{1+\varphi}
=K.
\end{aligned}
$$
Hence $\hat{K}=0$, so $(\Sigma,\hat{g})$
is complete and flat, hence parabolic by
Theorem \ref{theorem-alias-ahlfors-huber},
and so is $(\Sigma,g)$.
\end{proof}

\begin{theorem}[Heinz, 1955]
\label{theorem-alias-heinz}
The only entire graphs of constant mean
curvature in $\mathbb{R}^3$ are planes.
That is, if $H$ is constant, the only
entire solutions of the $H$-graph
equation
$$
\operatorname{Div}\left(
\frac{Du}{\sqrt{1+|Du|^2}}
\right)=2H
$$
are affine functions.
\end{theorem}

\begin{proof}
It suffices to show $H=0$ and apply
Theorem \ref{theorem-alias-bernstein}.
Integrate the equation over a bounded
domain $\Omega\subset\mathbb{R}^2$ and use
the divergence theorem in the plane. Since
the field $Du/\sqrt{1+|Du|^2}$ has norm
less than $1$,
$$
2|H|\operatorname{area}(\Omega)
\leq\operatorname{length}(\partial\Omega).
$$
For a disc of radius $R$ this gives
$|H|\leq 1/R$; letting $R\to\infty$,
$H=0$.
\end{proof}

\medskip\noindent
{\bf The Lorentz--Minkowski space.}
Let $\mathbb{L}^3$ be $\mathbb{R}^3$ with
the metric
$$
\left\langle{} \cdot,\cdot
\right\rangle{}=dx^2+dy^2-dz^2.
$$
A surface $\Sigma\subset\mathbb{L}^3$ is
\emph{spacelike} if every tangent plane is
spacelike, equivalently if the induced
metric is Riemannian. Then there is a
globally defined future-pointing timelike
unit normal $N$, so $\Sigma$ is orientable
and its Gauss map takes values in the
hyperbolic plane
$$
\mathbb{H}^2_+=\{x\in\mathbb{L}^3:
\left\langle{} x,x\right\rangle{}=-1,
\ z>0\}.
$$
With $A=-dN$, a spacelike surface is
\emph{maximal} if $H=0$. Since
$\left\langle{} N,N\right\rangle{}=-1$, the
Gauss equation reads $K=-\det A
=-\kappa_1\kappa_2$, so for a maximal
surface
$$
K=-\kappa_1\kappa_2
=\frac{1}{2}(\kappa_1^2+\kappa_2^2)
=\frac{1}{2}|A|^2\geq 0.
$$

\medskip\noindent
{\bf The Calabi--Bernstein theorem
(1970).}

\begin{theorem}[Calabi--Bernstein,
parametric version]
\label{theorem-alias-calabi-bernstein}
The only complete maximal surfaces in
$\mathbb{L}^3$ are the spacelike planes.
\end{theorem}

\begin{remark}
\label{remark-alias-euclidean-fails}
The Euclidean analogue is false: the
catenoid and the helicoid are complete
non-planar minimal surfaces. In
$\mathbb{R}^3$ the statement holds only
for graphs.
\end{remark}

\begin{theorem}[Calabi--Bernstein,
non-parametric version]
\label{theorem-alias-calabi-bernstein-graphs}
The only entire maximal graphs in
$\mathbb{L}^3$ are the spacelike planes.
That is, the only entire solutions of
$$
\operatorname{Div}\left(
\frac{Du}{\sqrt{1-|Du|^2}}
\right)=0,
\qquad
|Du|<1,
$$
are affine functions.
\end{theorem}

The second theorem does not follow
directly from the first, because entire
spacelike graphs need not be complete
(Remark \ref{remark-alias-entire-graphs}).

\begin{lemma}
\label{lemma-alias-complete-spacelike}
Every complete spacelike surface
$S\subset\mathbb{L}^3$ is an entire
spacelike graph $z=u(x,y)$ with
$u:\mathbb{R}^2\to\mathbb{R}$ and
$|Du|<1$.
\end{lemma}

\begin{proof}
The lecture constructed a covering map.
Let $\pi:S\to\mathbb{R}^2$ be the
projection $\pi(x,y,z)=(x,y)$. For a
tangent vector $v=(v_1,v_2,v_3)$,
$$
|d\pi(v)|^2=v_1^2+v_2^2
\geq v_1^2+v_2^2-v_3^2
=\left\langle{} v,v\right\rangle{},
$$
so $\pi$ does not decrease lengths, and in
particular $d\pi$ is injective and $\pi$ is
a local diffeomorphism. Since $S$ is
complete, $\pi$ is a covering map
(Proposition
\ref{differential-geometry-%
proposition-local-diffeomorphism-covering}).
As $\mathbb{R}^2$ is simply connected,
$\pi$ is a diffeomorphism, so $S$ is the
graph of a function $u$ on
$\mathbb{R}^2$, and spacelike means
$|Du|<1$.
\end{proof}

\begin{remark}
\label{remark-alias-entire-graphs}
In $\mathbb{R}^3$ every entire graph is
complete, since its induced metric
dominates $dx^2+dy^2$. The converse of
Lemma \ref{lemma-alias-complete-spacelike}
fails: there are entire spacelike graphs
in $\mathbb{L}^3$ that are not complete.
For example $u(x,y)=\log\cosh x$ has
$|Du|=|\tanh x|<1$, so its graph is
spacelike, but the ray $y=0$, $x\geq 0$
has induced length
$\int_0^\infty\operatorname{sech}x\,dx
=\pi/2$.
\end{remark}

\begin{proposition}
\label{proposition-alias-maximal-normal}
Let $S\subset\mathbb{L}^3$ be a maximal
surface, let $a\in\mathbb{L}^3$ be
nonzero with
$\left\langle{} a,a\right\rangle{}=c\in
\{0,1,-1\}$, and put
$\varphi=\left\langle{} N,a
\right\rangle{}$. Then
$\nabla\varphi=-A(a^{\top})$,
$$
|\nabla\varphi|^2=K(\varphi^2+c),
\qquad
\Delta\varphi=2K\varphi.
$$
\end{proposition}

\begin{proof}
The proof was done in class; here is a
version. Now $a=a^{\top}-\varphi N$, so
$|a^{\top}|^2=c+\varphi^2$. As in
Proposition
\ref{proposition-alias-minimal-normal},
$\nabla\varphi=-A(a^{\top})$. For a
maximal surface $A^2=K\operatorname{id}$,
which gives the first identity. Next
$\nabla_va^{\top}=-\varphi A(v)$, and by
Codazzi--Mainardi
$$
\nabla_v\nabla\varphi
=-(\nabla_{a^{\top}}A)(v)
+\varphi A^2(v).
$$
Taking traces, with
$\operatorname{tr}A=0$ and
$\operatorname{tr}(A^2)=2K$, gives
$\Delta\varphi=2K\varphi$.
\end{proof}

\begin{proof}[Proof of Theorem
\ref{theorem-alias-calabi-bernstein}]
Since $K=\frac{1}{2}|A|^2\geq 0$ and $S$
is complete, $S$ is parabolic by Theorem
\ref{theorem-alias-ahlfors-huber}. Let
$b=(1,0,1)$, so $c=\left\langle{} b,b
\right\rangle{}=0$, and
$\varphi=\left\langle{} N,b
\right\rangle{}$. Writing
$N=(N_1,N_2,N_3)$ with
$N_3=\sqrt{1+N_1^2+N_2^2}>|N_1|$,
$$
\varphi=N_1-N_3<0.
$$
By Proposition
\ref{proposition-alias-maximal-normal}
and Proposition
\ref{proposition-alias-div-laplacian},
$$
\Delta\frac{1}{\varphi}
=-\frac{\Delta\varphi}{\varphi^2}
+\frac{2|\nabla\varphi|^2}{\varphi^3}
=-\frac{2K}{\varphi}
+\frac{2K}{\varphi}=0.
$$
So $1/\varphi$ is a negative harmonic
function, hence constant by parabolicity.
Then $\nabla\varphi=0$, so
$K\varphi^2=0$, $K\equiv 0$ and
$A\equiv 0$: $S$ is a spacelike plane.
\end{proof}

The notes taken in class only record
that this Laplacian vanishes; the
computation above fills the gap.

\begin{theorem}
\label{theorem-alias-maximal-euclid-complete}
The only maximal surfaces in
$\mathbb{L}^3$ that are complete with
respect to the Euclidean metric of
$\mathbb{R}^3$ are the spacelike planes.
\end{theorem}

\begin{proof}
As in Chern's theorem, one changes the
metric conformally. Let $g$ be the
induced (Riemannian) metric, $g'$ the
metric induced by the Euclidean metric,
$a=(0,0,1)$, so $c=\left\langle{} a,a
\right\rangle{}=-1$, and
$\varphi=\left\langle{} N,a
\right\rangle{}=-N_3\leq -1$. Put
$$
\hat{g}=(1-\varphi)^2g,
\qquad
u=\log(1-\varphi).
$$
By Proposition
\ref{proposition-alias-maximal-normal},
$$
\begin{aligned}
\Delta u
&=-\frac{\Delta\varphi}{1-\varphi}
-\frac{|\nabla\varphi|^2}
{(1-\varphi)^2}\\
&=\frac{-2K\varphi}{1-\varphi}
+\frac{K(1+\varphi)}{1-\varphi}
=K,
\end{aligned}
$$
so $\hat{K}=e^{-2u}(K-\Delta u)=0$.

Next, $\hat{g}(v,v)\geq\frac{1}{2}
g'(v,v)$ for tangent $v$. Indeed, if
$N=(\sinh\theta,0,\cosh\theta)$ after a
rotation about the $z$-axis, the tangent
vectors $e_1=(\cosh\theta,0,\sinh\theta)$
and $e_2=(0,1,0)$ are $g$-orthonormal and
also orthogonal for $g'$, with
$g'(e_1,e_1)=2N_3^2-1$ and
$g'(e_2,e_2)=1$. Thus
$$
g'(v,v)\leq(2N_3^2-1)g(v,v)
\leq 2(1+N_3)^2g(v,v),
$$
i.e.\ $g'\leq 2\hat{g}$. (In Chern's
theorem the constant was $1$; here it is
$\frac{1}{2}$.) Since $g'$ is complete,
so is $\hat{g}$. Hence $(S,\hat{g})$ is
complete and flat, thus parabolic, and so
is $(S,g)$. Now $1/\left\langle{} N,b
\right\rangle{}$ with $b=(1,0,1)$ is a
negative harmonic function, and the end of
the proof of Theorem
\ref{theorem-alias-calabi-bernstein}
applies.
\end{proof}

The notes taken in class stop at
``we find a negative harmonic function'';
the comparison $g'\leq 2\hat{g}$ and the
curvature computation are completed here.
Theorem
\ref{theorem-alias-calabi-bernstein-graphs}
follows, since an entire graph is complete
for the Euclidean metric
(Remark \ref{remark-alias-entire-graphs}).

\medskip\noindent
{\bf Liebmann's theorem, revisited.}
The lecture returned to Theorem
\ref{theorem-alias-liebmann} to prove it
with the Minkowski formulas.

\begin{proof}[Proof of Theorem
\ref{theorem-alias-liebmann}]
We follow \cite[Theorem 3.4,
pp.~25--26]{alias2006analisis}. Write
$\Sigma=\partial\Omega$ with the inner
normal $N$. By Theorem
\ref{theorem-alias-elliptic-point} and
Remark
\ref{remark-alias-compact-orientable}, $K$
is a positive constant and the normal
curvatures at some point are positive.
Since $H^2\geq K>0$, by connectedness
$H>0$ everywhere, and by Lemma
\ref{lemma-alias-h2-geq-k},
$H\geq\sqrt{K}$, with equality exactly at
umbilical points. By Theorem
\ref{theorem-alias-hadamard} and the proof
of Proposition
\ref{proposition-alias-convexity-curvature},
for each $p$ the height function $h_p$
has a strict global minimum at $p$, so
$\operatorname{int}\Omega$ lies in
$\operatorname{int}\Pi_p^+$. Hence, for $c$
in the interior of $\Omega$,
$\left\langle{} c-p,N(p)
\right\rangle{}>0$ for all $p$.

Multiply (\ref{equation-alias-minkowski-1})
by $\sqrt{K}$ and subtract it from
(\ref{equation-alias-minkowski-2}):
$$
\int_\Sigma\left(H-\sqrt{K}\right)
\left(1+\sqrt{K}\left\langle{} c-p,N
\right\rangle{}\right)dp=0.
$$
The integrand is nonnegative, hence zero,
so $H=\sqrt{K}$ and $\Sigma$ is totally
umbilical. A compact connected totally
umbilical surface is a sphere, of radius
$1/H=1/\sqrt{K}$.
\end{proof}

\medskip\noindent
{\bf Problem session.}
The properties of $\nabla$ were proved by
differentiating along curves realizing the
vectors and taking tangent parts. The
gradient, Hessian and Laplacian are
differential operators; for instance, for
a smooth $\phi:\mathbb{R}\to\mathbb{R}$,
$$
\begin{aligned}
\nabla\phi(f)&=\phi'(f)\nabla f,\\
\Delta\phi(f)&=\phi'(f)\Delta f
+\phi''(f)|\nabla f|^2.
\end{aligned}
$$
The second follows from the first and
Proposition
\ref{proposition-alias-div-laplacian}.

\begin{definition}
\label{definition-alias-focal-point}
Fix $p_0\in\Sigma$ and let
$\ell_0=\{p_0+tN(p_0):t\in\mathbb{R}\}$ be
the normal line at $p_0$. A point
$q=p_0+t_0N(p_0)$ of $\ell_0$ is a
\emph{focal point} of $\Sigma$ along
$\ell_0$ if the map
$$
\phi_{t_0}:\Sigma\to\mathbb{R}^3,
\qquad
\phi_{t_0}(p)=p+t_0N(p),
$$
is singular at $p_0$, i.e.\ its
differential at $p_0$ has non-trivial
kernel.
\end{definition}

\begin{exercise}
\label{exercise-alias-focal-points}
Show that the focal points of $\Sigma$
along $\ell_0$ are the points
$$
q=p_0+\frac{1}{\kappa_i(p_0)}N(p_0),
$$
for $i=1,2$ with $\kappa_i(p_0)\neq 0$. In
particular there are at most two focal
points along $\ell_0$.
\end{exercise}

\begin{proof}[Solution]
Let $v\in T_{p_0}\Sigma$ and let $\alpha$
be a curve with $\alpha(0)=p_0$ and
$\alpha'(0)=v$. Then
$$
\begin{aligned}
(d\phi_{t_0})_{p_0}(v)
&=\frac{d}{ds}\Big|_{s=0}
\left(\alpha(s)+t_0N(\alpha(s))\right)\\
&=v-t_0A_{p_0}(v).
\end{aligned}
$$
This vanishes for some $v\neq 0$ if and
only if $t_0\neq 0$ and $1/t_0$ is an
eigenvalue of $A_{p_0}$, i.e.\
$t_0=1/\kappa_i(p_0)$ for some $i$ with
$\kappa_i(p_0)\neq 0$. So the focal points
correspond to the non-zero principal
curvatures, and there are at most two
(at most one if $p_0$ is umbilical, and
none at a planar point).
\end{proof}

\medskip\noindent
{\bf Complement: Alexandrov's theorem and
the first eigenvalue.}
Daniel's notes stop here. The rest of the
abstract (topics 4 and 5) is summarized
from Chapters 4 and 5 of
\cite{alias2006analisis}, not from the
lectures. Throughout, $\Sigma$ is compact,
$\Sigma=\partial\Omega$, and $N$ is the
inner unit normal.

For $Z$ a vector field on $\Omega$, the
divergence theorem reads
$$
\int_\Omega\operatorname{Div}Z\,dx
=-\int_\Sigma\left\langle{} Z,N
\right\rangle{}dp,
$$
the sign coming from the inner normal
\cite[Theorem 4.2, p.~33]{alias2006analisis}.
With $Z(x)=x-c$ it gives
\begin{equation}
\label{equation-alias-volume}
\operatorname{vol}(\Omega)
=-\frac{1}{3}\int_\Sigma\left\langle{}
p-c,N(p)\right\rangle{}dp
\end{equation}
for every $c\in\mathbb{R}^3$
\cite[Lemma 4.3, p.~33]{alias2006analisis}.

\begin{theorem}[Reilly's formula]
\label{theorem-alias-reilly-formula}
For every smooth
$F:\Omega\to\mathbb{R}$, with
$f=F|_\Sigma$ and
$\partial F/\partial N=\left\langle{}
\operatorname{Grad}F,N\right\rangle{}$,
$$
\begin{aligned}
&\int_\Omega\left((\bar{\Delta}F)^2
-|\operatorname{Hess}F|^2\right)dx\\
&=2\int_\Sigma\left(H
\frac{\partial F}{\partial N}
-\Delta f\right)
\frac{\partial F}{\partial N}\,dp\\
&\quad+\int_\Sigma\left\langle{}
A(\nabla f),\nabla f\right\rangle{}dp,
\end{aligned}
$$
where $\bar{\Delta}$ is the Euclidean
Laplacian.
\end{theorem}

This is \cite[Theorem 4.5,
pp.~35--37]{alias2006analisis}, a special
case of \cite{reilly1977applications}. The
proof applies the divergence theorem to
$Z=\bar{\Delta}F\operatorname{Grad}F
-\frac{1}{2}\operatorname{Grad}
|\operatorname{Grad}F|^2$, whose
divergence is
$(\bar{\Delta}F)^2-|\operatorname{Hess}F|^2$
by Bochner's formula in $\mathbb{R}^3$. As
a check, for $F=|x|^2$ on a ball of radius
$r$ both sides equal $32\pi r^3$.

\begin{theorem}[Heintze--Karcher--Ros
inequality]
\label{theorem-alias-heintze-karcher-ros}
If the mean curvature of $\Sigma$ (with
respect to the inner normal) never
vanishes, then $H>0$ and
$$
\operatorname{vol}(\Omega)\leq
\frac{1}{3}\int_\Sigma\frac{1}{H}\,dp,
$$
with equality if and only if $\Sigma$ is a
sphere.
\end{theorem}

\begin{proof}[Sketch]
See \cite[Theorem 4.8,
pp.~40--42]{alias2006analisis}, which
adapts Ros' proof. $H$ is positive
somewhere by Theorem
\ref{theorem-alias-elliptic-point} and
Remark
\ref{remark-alias-compact-orientable}, so
$H>0$. Let $F$ solve $\bar{\Delta}F=1$ in
$\Omega$ with $F=0$ on $\Sigma$. Reilly's
formula gives
$$
\int_\Omega\left(1
-|\operatorname{Hess}F|^2\right)dx
=2\int_\Sigma H\left(
\frac{\partial F}{\partial N}\right)^2dp.
$$
By Cauchy--Schwarz,
$|\operatorname{Hess}F|^2\geq\frac{1}{3}
(\bar{\Delta}F)^2=\frac{1}{3}$, so
$\int_\Sigma H(\partial F/\partial N)^2
\leq\frac{1}{3}\operatorname{vol}(\Omega)$.
Since
$\operatorname{vol}(\Omega)
=-\int_\Sigma\partial F/\partial N$,
Cauchy--Schwarz again gives
$$
\operatorname{vol}(\Omega)^2\leq
\frac{1}{3}\operatorname{vol}(\Omega)
\int_\Sigma\frac{1}{H}\,dp.
$$
In the case of equality,
$\operatorname{Hess}F=\frac{1}{3}
\left\langle{} \cdot,\cdot
\right\rangle{}$, so $F$ is a quadratic
polynomial and $\Sigma=F^{-1}(0)$ is a
sphere.
\end{proof}

\begin{theorem}[Alexandrov]
\label{theorem-alias-alexandrov}
A compact regular (hence embedded) surface
$\Sigma\subset\mathbb{R}^3$ with constant
mean curvature $H$ is a sphere of radius
$1/H$.
\end{theorem}

\begin{proof}
See \cite[Theorem 4.9,
pp.~42--43]{alias2006analisis} and
\cite{alexandrov1958uniqueness}. With the
inner normal, the constant $H$ is positive
by Theorem
\ref{theorem-alias-elliptic-point} and
Remark
\ref{remark-alias-compact-orientable}.
Theorem
\ref{theorem-alias-heintze-karcher-ros}
gives
$$
\operatorname{area}(\Sigma)\geq
3H\operatorname{vol}(\Omega).
$$
The first Minkowski formula with $c=0$
and (\ref{equation-alias-volume}) give
$$
\operatorname{area}(\Sigma)
=-H\int_\Sigma\left\langle{} p,N
\right\rangle{}dp
=3H\operatorname{vol}(\Omega).
$$
So equality holds in the
Heintze--Karcher--Ros inequality, and
$\Sigma$ is a sphere.
\end{proof}

Embeddedness is essential: Wente's tori
are compact immersed surfaces of constant
mean curvature (not checked). Compare
Theorem
\ref{xxii-egd-2026-theorem-alexandrov58},
stated there for boundaries of domains.
The booklet also gives, from Reilly's
formula, Garay's estimate
$\max_\Sigma H\geq 3\operatorname{vol}
(\Omega)/(r^2\operatorname{area}(\Sigma))$
for ovaloids, where $r$ is the radius of
the circumscribed sphere, with equality
only for spheres \cite[Theorem 4.6,
p.~38]{alias2006analisis}, and a second
proof of Theorem
\ref{theorem-alias-liebmann} from
Theorem
\ref{theorem-alias-heintze-karcher-ros}
\cite[Theorem 4.10,
pp.~43--44]{alias2006analisis}.

\medskip\noindent
{\bf The first eigenvalue.}
By the divergence theorem, $\Delta$ is
self-adjoint for the $L^2$ inner product
on $C^\infty(\Sigma)$. A real $\lambda$ is
an \emph{eigenvalue} of $\Delta$ if
$\Delta f+\lambda f=0$ for some
$f\not\equiv 0$. The eigenvalue $0$ has
only the constants as eigenfunctions,
every other eigenvalue is positive, and
the eigenvalues form a sequence
$$
0=\lambda_0<\lambda_1<\lambda_2<\cdots
\to\infty
$$
with finite multiplicities. The first
positive eigenvalue satisfies
\begin{equation}
\label{equation-alias-rayleigh}
\lambda_1\int_\Sigma f^2dp
\leq\int_\Sigma|\nabla f|^2dp
\end{equation}
for every $f$ with $\int_\Sigma f\,dp=0$,
with equality if and only if $f$ is an
eigenfunction of $\lambda_1$
\cite[Section 5.1,
pp.~45--47]{alias2006analisis}. For the
sphere of radius $r$, the eigenvalues
are $k(k+1)/r^2$; in particular
$\lambda_1=2/r^2=2K$, with the linear
functions as eigenfunctions
\cite[Example 5.3,
pp.~47--49]{alias2006analisis}.

\begin{lemma}
\label{lemma-alias-lambda1-centroid}
Let $c=\frac{1}{\operatorname{area}
(\Sigma)}\int_\Sigma p\,dp$ be the centre of
gravity of $\Sigma$. Then
$$
\lambda_1\leq
\frac{2\operatorname{area}(\Sigma)}
{\int_\Sigma|p-c|^2dp},
$$
with equality if and only if $\Sigma$ is a
sphere centred at $c$.
\end{lemma}

\begin{proof}
See \cite[Lemma 5.4,
pp.~49--51]{alias2006analisis}. For a
unit vector $a$, the height function
$h(p)=\left\langle{} p-c,a
\right\rangle{}$ has mean zero and
$|\nabla h|^2=1-\left\langle{} a,N
\right\rangle{}^2$ by
(\ref{equation-alias-grad-height}). Apply
(\ref{equation-alias-rayleigh}) to $h$ for
$a=e_1,e_2,e_3$ and add, using
$\sum_ih_i^2=|p-c|^2$ and
$\sum_i\left\langle{} N,e_i
\right\rangle{}^2=1$:
$$
\lambda_1\int_\Sigma|p-c|^2dp
\leq 3\operatorname{area}(\Sigma)
-\operatorname{area}(\Sigma).
$$
If equality holds, each $h_i$ is an
eigenfunction. Since
$\Delta h_i=2H\left\langle{} N,e_i
\right\rangle{}$, this gives
$2HN+\lambda_1(p-c)=0$, so $p-c$ is
normal everywhere, $|p-c|^2$ is constant
by (\ref{equation-alias-grad-distance}),
and $\Sigma$ is a sphere centred at $c$.
\end{proof}

\begin{theorem}[Reilly]
\label{theorem-alias-reilly-eigenvalue}
Let $\Sigma\subset\mathbb{R}^3$ be a
compact regular surface and
$A_\Sigma=\operatorname{area}(\Sigma)$.
Then
$$
\begin{aligned}
\lambda_1&\leq\frac{2}{A_\Sigma}
\int_\Sigma H^2dp,\\
\lambda_1\left(\int_\Sigma H\,dp
\right)^2&\leq 2A_\Sigma
\int_\Sigma K^2dp,\\
\lambda_1&\leq 2\left(
\frac{A_\Sigma}
{3\operatorname{vol}(\Omega)}\right)^2.
\end{aligned}
$$
Equality holds in any of them (in the
second, when $\int_\Sigma H\,dp\neq 0$) if
and only if $\Sigma$ is a sphere. The
second is written multiplied out because
$\int_\Sigma H\,dp$ need not be positive
for a non-convex surface (not checked).
\end{theorem}

\begin{proof}
See \cite[Theorem 5.5,
pp.~51--53]{alias2006analisis} and
\cite{reilly1977first}. Let $c$ be the
centre of gravity, so that
$2A_\Sigma\geq\lambda_1\int_\Sigma
|p-c|^2dp$ by Lemma
\ref{lemma-alias-lambda1-centroid}. For the
first inequality, multiply by
$\int_\Sigma H^2$ and use Cauchy--Schwarz
and the first Minkowski formula:
$$
\begin{aligned}
2A_\Sigma\int_\Sigma H^2dp
&\geq\lambda_1\left(\int_\Sigma
H\left\langle{} p-c,N\right\rangle{}dp
\right)^2\\
&=\lambda_1A_\Sigma^2.
\end{aligned}
$$
The second is the same with $K$ and the
second Minkowski formula. For the third,
multiply by $A_\Sigma$ and use
(\ref{equation-alias-volume}):
$$
2A_\Sigma^2\geq\lambda_1\left(
\int_\Sigma\left\langle{} p-c,N
\right\rangle{}dp\right)^2
=9\lambda_1\operatorname{vol}(\Omega)^2.
$$
Equality in any of them forces equality
in Lemma
\ref{lemma-alias-lambda1-centroid}.
\end{proof}

For the sphere of radius $r$ all three
right-hand sides equal $2/r^2$. For
surfaces of genus zero there is also
Hersch's inequality
$\lambda_1\operatorname{area}(\Sigma)
\leq 8\pi$, valid for every Riemannian
metric on the sphere; it is not in the
booklet (not checked).

\begin{theorem}
\label{theorem-alias-deshmukh}
Let $\Sigma\subset\mathbb{R}^3$ be an
ovaloid. If $K\leq\frac{1}{2}\lambda_1$
everywhere, then $\Sigma$ is a sphere.
Equivalently, every ovaloid satisfies
$$
\lambda_1\leq 2\max_\Sigma K,
$$
with equality if and only if $\Sigma$ is
a sphere.
\end{theorem}

This is \cite[Theorem 5.6,
pp.~53--55]{alias2006analisis}, an
adaptation of a result of Deshmukh. With
$f=\frac{1}{2}|p-c|^2$ and
$g=\left\langle{} p-c,N\right\rangle{}$,
$c$ the centre of gravity, one combines
(\ref{equation-alias-laplacian-f}) and
(\ref{equation-alias-laplacian-g}) into
$$
\frac{1}{2}\Delta(g^2)
+2\operatorname{div}(Hg\nabla f)
=-3K|\nabla f|^2+2Hg+4Kf,
$$
integrates, and uses the first Minkowski
formula and Lemma
\ref{lemma-alias-lambda1-centroid} to get
$$
\frac{3}{2}\int_\Sigma K|\nabla f|^2dp
\leq\int_\Sigma(2K-\lambda_1)f\,dp,
$$
valid for every compact $\Sigma$. If
$K>0$ and $2K\leq\lambda_1$, then
$\nabla f=0$ and $\Sigma$ is a sphere.

\section{Una introducción a la Geometría de
Finsler}
\label{section-javaloyes-finsler-geometry}

\medskip\noindent
{\bf Curso Avanzado 3.}

Miguel Ángel Javaloyes,
Universidad de Murcia -- España.

\medskip\noindent
{\bf Schedule.}

Week 2, Monday to Thursday, 14:30--16:00
(15--18 July 2024).

\medskip\noindent
{\bf Official listing.}

\href{%
https://naturales.ues.edu.sv/%
ecimpasv2024/advanced-curse-3/}{%
Programme entry.}

\medskip\noindent
{\bf Abstract.}

En este minicurso, intentaremos dar cuenta de
la interacción entre la Geometría de Finsler
y los espacios-tiempos en todas sus múltiples
formas, y las técnicas que permiten manejar
la dependencia de la dirección implícita en
la Geometría de Finsler.

Partiremos de un espacio-tiempo estacionario
mostrando que sus geodésicas luminosas están
controladas por una métrica de Randers
$F(v)=\sqrt{h(v,v)}+\omega(v)$, siendo $h$
una métrica de Riemann y $\omega$ una forma.
Como consecuencia, la causalidad del
espacio-tiempo estacionario puede describirse
en términos de las propiedades de completitud
de esta métrica de Randers. Por otro lado, se
puede utilizar una métrica de Randers para
encontrar soluciones al problema de Zermelo,
cuyo objetivo es encontrar trayectorias que
minimicen el tiempo en presencia de un viento
o corriente (independiente del tiempo).
Además, siempre es posible asociar un
espaciotiempo estacionario a la métrica de
Randers para estudiar el problema de Zermelo.
Lo que es más interesante, cuando el viento
depende del tiempo, el problema de Zermelo ya
no se puede modelar con una métrica de
Randers, pero el modelo de espacio-tiempo
sigue funcionando, que en este caso ya no es
estacionario. La moraleja de la historia es
que todo espacio-tiempo clásico puede
utilizarse para estudiar un determinado
problema de Zermelo. Cuando la velocidad sin
viento depende de la dirección, entonces hay
que considerar los espacio-tiempos de
Finsler. Todos estos resultados se combinan
para dar una sorprendente aplicación del
espacio-tiempo al estudio de la evolución de
un incendio forestal.

El siguiente paso será introducir el cálculo
anisotrópico para manejar los tensores que
dependen de la dirección. Teniendo a nuestra
disposición el cálculo anisotrópico
introduciremos las conexiones anisotrópicas
asociadas con una métrica de pseudo-Finsler y
luego el tensor de curvatura de Chern y la
curvatura bandera. Compararemos las
conexiones anisotrópicas con las conexiones
clásicas sobre el fibrado tangente y luego
derivaremos las identidades de Bianchi. A
continuación, definiremos geodésicas y campos
de Jacobi, mostrando sus propiedades
variacionales. También mostraremos cómo
obtener las ecuaciones de Gauss-Codazzi de
una subvariedad.

En la última parte volveremos a la noción de
espacio tiempo de Finsler y estructura de
conos, estableciendo la relación entre ambas
nociones y mostrando como se pueden definir
ejemplos de manera sencilla. Se introducirá
la noción de campos de Killing y conformes, y
se enfatizará el papel de estos últimos en
las estructuras de conos. Discutiremos las
ecuaciones de Einstein para los
espacios-tiempos de Finsler y cómo se puede
interpretar la divergencia cero del tensor
energía-impulso anisotrópico usando el
transporte paralelo. Cerraremos el minicurso
con algunas palabras sobre cómo hacer que los
espacios-tiempos de Finsler sean compatibles
con el enfoque axiomático de
Ehlers-Pirani-Schild sobre los fundamentos de
la relatividad general.

\medskip\noindent
{\bf Sources of these notes.}
These notes cover only the first of the four
sessions; Daniel's notes taken during the
course stop there. The lecturer pointed to
his papers \cite{javaloyes2014chern} and
\cite{javaloyes2019anisotropic}, where
pseudo-Finsler metrics, their fundamental
tensor and the anisotropic calculus of the
later sessions are developed. For the
material on Minkowski norms we follow and
cite \cite{javaloyes2014definition}, whose
Section 2 contains the statements of this
session (we use the numbering of its arXiv
version). Several statements of the notes
were corrected or made precise (the
fundamental inequality, the comparison of
the fundamental tensor with the second
fundamental form of the indicatrix, the
notions of convexity), the proofs that were
skipped in class were filled in, and the
two gaps of the notes are completed from
the references. The points not verified are
marked ``(not checked)''. Uniformly convex
norms on $\mathbb{R}^{n+1}$ also appear in
the anisotropic surface energies of
Section~%
\ref{xxii-egd-2026-section-%
alexandrov-crystalline}.

\medskip\noindent
{\bf Plan of the course.}
\begin{enumerate}
\item Basic notions, the Zermelo navigation
problem and other applications, interplay
between Randers metrics and stationary
spacetimes.
\item Anisotropic calculus, or how to handle
anisotropic tensors and curvatures.
\item Finsler spacetimes: examples,
applications and foundations.
\end{enumerate}

\medskip\noindent
{\bf Session 1.}
The plan of the session was: definition and
first properties of Finsler metrics; Randers
metrics and stationary spacetimes. The notes
record only the first part.

\medskip\noindent
{\bf The concept of norm in analysis.}
A norm on $\mathbb{R}^n$ is a function
$\|\cdot\| : \mathbb{R}^n \to [0,+\infty)$
such that
\begin{enumerate}
\item $\|\lambda v\| = |\lambda| \|v\|$,
which means that the length of a curve does
not depend on its parametrization;
\item $\|v\| = 0$ if and only if $v = 0$:
nonzero vectors have positive length;
\item $\|v + w\| \leq \|v\| + \|w\|$:
straight segments are the shortest paths.
\end{enumerate}
The \emph{indicatrix} of $\|\cdot\|$ is the
unit sphere
$\Sigma = \{v : \|v\| = 1\}$, and
$B = \{v : \|v\| \leq 1\}$ is its closed
unit ball. In Finsler geometry the first
condition is weakened to
$\|\lambda v\| = \lambda \|v\|$ for
$\lambda > 0$ (positive homogeneity), so
that $\|-v\|$ may differ from $\|v\|$.

\begin{proposition}
\label{proposition-javaloyes-convex-ball}
Let $\|\cdot\| : \mathbb{R}^n \to
[0,+\infty)$ vanish only at $0$ and be
positively homogeneous. Then $B$ is convex
if and only if $\|\cdot\|$ satisfies the
triangle inequality.
\end{proposition}

\begin{proof}
If the triangle inequality holds and
$u, w \in B$, $t \in [0,1]$, then
$$
\|tu + (1-t)w\| \leq t\|u\| + (1-t)\|w\|
\leq 1.
$$
Conversely, assume $B$ convex and let
$u, w \neq 0$ (otherwise there is nothing
to prove). Put $a = \|u\|$, $b = \|w\|$.
Then $u/a$ and $w/b$ lie in $B$, hence so
does the convex combination
$$
\frac{u + w}{a + b} =
\frac{a}{a+b}\frac{u}{a}
+ \frac{b}{a+b}\frac{w}{b},
$$
and positive homogeneity gives
$\|u + w\| \leq a + b$.
\end{proof}

\noindent
This is the argument of
\cite[Proposition 2.3]%
{javaloyes2014definition}.

\begin{definition}
\label{definition-javaloyes-minkowski-norm}
Let $V$ be a finite dimensional real vector
space. A function
$F : V \to [0,+\infty)$ is a
\emph{Minkowski norm} if
\begin{enumerate}
\item $F$ is smooth on $V \setminus \{0\}$;
\item $F(v) = 0$ if and only if $v = 0$;
\item $F$ is positively homogeneous of
degree one: $F(\lambda v) = \lambda F(v)$
for every $v \in V$ and $\lambda > 0$;
\item for every $v \in V \setminus \{0\}$
the \emph{fundamental tensor} of $F$ at $v$,
the symmetric bilinear form
$$
g_v(u,w) = \frac{1}{2}
\frac{\partial^2}{\partial t \, \partial s}
F^2(v + tu + sw) \Big|_{t=s=0},
$$
for $u, w \in V$, is positive definite.
\end{enumerate}
\end{definition}

\noindent
So $g_v$ is half the Hessian of $F^2$ at
$v$, and the triangle inequality has been
replaced by the fourth condition. We will
see that the fourth condition implies the
triangle inequality, and more: it says that
$\Sigma$ is \emph{strongly convex}, which is
stronger than strict convexity. Here $B$ is
\emph{strictly convex} if every tangent
hyperplane of $\Sigma$ meets $B$ at a single
point (equivalently, $\Sigma$ contains no
segment), while $\Sigma$ is strongly convex
if its second fundamental form with respect
to an inward transverse vector is positive
definite (see below). For example, the
$\ell^4$ norm
$F(x,y) = (x^4 + y^4)^{1/4}$ on
$\mathbb{R}^2$ is smooth away from $0$ and
its unit ball is strictly convex, but it is
not a Minkowski norm:
$F^2(1,s) = (1 + s^4)^{1/2}$ has vanishing
second derivative at $s = 0$, so
$g_{(1,0)}$ vanishes on $(0,1)$.

The fundamental tensor can be thought of as
a scalar product attached to each direction
$v$; it depends only on the direction, and
along that direction it reproduces $F$.

\begin{proposition}
\label{proposition-javaloyes-fund-tensor}
Let $F$ be a Minkowski norm and
$v \in V \setminus \{0\}$. Then
\begin{enumerate}
\item $g_{\lambda v} = g_v$ for
$\lambda > 0$;
\item $g_v(v,w) = \frac{1}{2} d(F^2)_v(w)
= F(v)\, dF_v(w)$ for all $w \in V$;
\item $g_v(v,v) = F(v)^2$;
\item $v$ is $g_v$-orthogonal to the tangent
space at $v$ of the level set
$F^{-1}(F(v))$; for $v \in \Sigma$ this is
$T_v\Sigma$.
\end{enumerate}
\end{proposition}

\begin{proof}
Since $F^2$ is positively homogeneous of
degree two, its second derivatives are
homogeneous of degree zero, which is the
first item. For the second, by homogeneity
$$
F^2(v + tv + sw) = (1+t)^2
F^2\left(v + \frac{s}{1+t} w\right)
$$
for $t$ near $0$. Differentiating in $s$ at
$s = 0$ gives
$(1 + t)\, d(F^2)_v(w)$, and then
differentiating in $t$ at $t = 0$ gives
$d(F^2)_v(w)$. Hence
$g_v(v,w) = \frac{1}{2} d(F^2)_v(w)
= F(v)\, dF_v(w)$. For the third item,
Euler's identity for homogeneous functions
gives $dF_v(v) = F(v)$, so
$g_v(v,v) = F(v)^2$. Finally, if $w$ is
tangent to the level set of $F$ through
$v$, then $dF_v(w) = 0$ and
$g_v(v,w) = 0$ by the second item.
\end{proof}

\noindent
See \cite[Proposition 2.2]%
{javaloyes2014definition}. In class the
second item was called a computation and
the third was skipped.

\medskip\noindent
{\bf The second fundamental form.}
Let $D$ be the flat affine connection of
$V$, that is, ordinary differentiation of
vector fields. Let $\Sigma \subset V$ be a
smooth hypersurface and $\xi$ a vector
field along $\Sigma$ transverse to it. For
vector fields $X, Y$ tangent to $\Sigma$,
the Gauss formula
$$
D_X Y = \nabla^\xi_X Y
+ \sigma^\xi(X,Y)\, \xi,
$$
with $\nabla^\xi_X Y$ tangent to $\Sigma$,
defines a torsion free connection
$\nabla^\xi$ on $\Sigma$ and a symmetric
tensor $\sigma^\xi$, the \emph{second
fundamental form} of $\Sigma$ with respect
to $\xi$. No metric is needed. For the
indicatrix $\Sigma = F^{-1}(1)$ of a
Minkowski norm, the field
$\xi(v) = -v$ is transverse to $\Sigma$,
since $dF_v(v) = F(v) = 1$, and it points
into $B$.

\begin{proposition}
\label{proposition-javaloyes-second-form}
Let $F$ be a Minkowski norm with indicatrix
$\Sigma$, and let $\xi(v) = -v$. For every
$v \in \Sigma$ and $u, w \in T_v\Sigma$,
$$
g_v(u,w) = \sigma^\xi_v(u,w).
$$
\end{proposition}

\begin{proof}
Extend $u, w$ to vector fields $X, Y$
tangent to $\Sigma$. Since $F^2 = 1$ on
$\Sigma$, the function $Y(F^2)$ vanishes on
$\Sigma$, and so does its derivative along
$X$. The Hessian of $F^2$ with respect to
$D$ is therefore, at points of $\Sigma$,
$$
\begin{aligned}
\operatorname{Hess}(F^2)(X,Y)
&= X(Y(F^2)) - (D_X Y)(F^2) \\
&= -\sigma^\xi(X,Y)\, \xi(F^2),
\end{aligned}
$$
because the tangent part $\nabla^\xi_X Y$
does not change $F^2$. By Euler's identity
$\xi(F^2)$ at $v$ equals
$d(F^2)_v(-v) = -2F(v)^2 = -2$. Since
$g = \frac{1}{2}\operatorname{Hess}(F^2)$,
we get $g_v(u,w) = \sigma^\xi_v(u,w)$.
\end{proof}

\noindent
This is the computation in the proof of
\cite[Proposition 2.3]%
{javaloyes2014definition}, with
$G = \frac{1}{2}F^2$. The normalization
matters: with $\xi = -v$ and $F(v) = 1$ the
two forms agree exactly; for a general
nonzero $v$ one uses $g_v = g_{v/F(v)}$.

\begin{remark}
\label{remark-javaloyes-change-transversal}
If $\rho$ is another transverse vector field
along $\Sigma$, then
$\rho = c\,\xi + Z$ with $c$ a nowhere
vanishing function and $Z$ tangent to
$\Sigma$. Substituting
$\xi = (\rho - Z)/c$ in the Gauss formula
gives
$$
D_X Y = \left(\nabla^\xi_X Y
- \frac{\sigma^\xi(X,Y)}{c} Z\right)
+ \frac{\sigma^\xi(X,Y)}{c}\, \rho,
$$
so $\sigma^\rho = \sigma^\xi / c$. Thus the
second fundamental forms for different
transverse fields are proportional, and
positive definiteness does not depend on
the choice of an inward pointing
transverse field ($c > 0$).
\end{remark}

\begin{remark}
\label{remark-javaloyes-second-form-curves}
For a curve $\alpha$ in $\Sigma$ with
$\alpha'(0) = u$, the Gauss formula gives
$$
\alpha''(0) = \nabla^\xi_{\alpha'}\alpha'(0)
+ \sigma^\xi(u,u)\, \xi.
$$
So $\sigma^\xi(u,u)$ is the
$\xi$-component of the acceleration of any
curve in $\Sigma$ with velocity $u$. When
$V$ is Euclidean and $\xi$ is the inner
unit normal, this is the normal curvature
of $\Sigma$ in the direction $u$ times
$|u|^2$.
\end{remark}

\begin{corollary}
\label{corollary-javaloyes-strong-convexity}
Let $F$ satisfy the first three conditions
of Definition~%
\ref{definition-javaloyes-minkowski-norm},
and let $v \in \Sigma$. Then $g_v$ is
positive definite if and only if
$\Sigma$ is strongly convex at $v$, that
is, $\sigma^\xi_v$ is positive definite
on $T_v\Sigma$.
\end{corollary}

\begin{proof}
The proofs of
Propositions~%
\ref{proposition-javaloyes-fund-tensor}
and~%
\ref{proposition-javaloyes-second-form}
did not use positive definiteness. Since
$dF_v(v) = 1$, we have
$V = \mathbb{R}v \oplus T_v\Sigma$, the two
summands are $g_v$-orthogonal and
$g_v(v,v) = 1$. Hence $g_v$ is positive
definite if and only if its restriction to
$T_v\Sigma$ is, and that restriction is
$\sigma^\xi_v$.
\end{proof}

\noindent
This is \cite[Proposition 2.3(v)]%
{javaloyes2014definition}.

\begin{proposition}[Fundamental inequality]
\label{proposition-javaloyes-fund-ineq}
Let $F$ be a Minkowski norm. For every
$v \in V \setminus \{0\}$ and $w \in V$,
$$
g_v(v,w) \leq F(v)F(w),
$$
with equality if and only if $w = \lambda v$
for some $\lambda \geq 0$.
\end{proposition}

\begin{proof}
If $w = \lambda v$ with $\lambda \geq 0$
both sides equal $\lambda F(v)^2$. If
$w = \lambda v$ with $\lambda < 0$, the
left-hand side is $\lambda F(v)^2 < 0$ and
the inequality is strict. Assume now that
$w \notin \mathbb{R}v$.

First, for $u \neq 0$ and $x \in V$, write
$g = \frac{1}{2}\operatorname{Hess}(F^2)
= dF \otimes dF + F \operatorname{Hess} F$.
Using $dF_u(x) = g_u(u,x)/F(u)$ and
$g_u(u,u) = F(u)^2$,
$$
F(u) \operatorname{Hess} F_u(x,x) =
g_u(x,x) - \frac{g_u(u,x)^2}{g_u(u,u)}.
$$
By the Cauchy--Schwarz inequality for
$g_u$, this is $\geq 0$, with equality only
for $x \in \mathbb{R}u$.

Let $c(t) = v + t(w - v)$ for
$t \in [0,1]$. This segment avoids $0$,
and $w - v \notin \mathbb{R}c(t)$ for all
$t$, because $w \notin \mathbb{R}v$. So
$f(t) = F(c(t))$ is smooth and
$f''(t) = \operatorname{Hess}
F_{c(t)}(w - v, w - v) > 0$. Strict
convexity of $f$ gives
$f(1) > f(0) + f'(0)$, that is,
$$
F(w) > F(v) + dF_v(w - v) = dF_v(w),
$$
using $dF_v(v) = F(v)$. Multiplying by
$F(v) > 0$ and using
$g_v(v,w) = F(v)\, dF_v(w)$ gives
$g_v(v,w) < F(v)F(w)$.
\end{proof}

\noindent
The statement is
\cite[Remark 2.9, (8)]%
{javaloyes2014definition}.

\begin{corollary}
\label{corollary-javaloyes-triangle}
A Minkowski norm satisfies the triangle
inequality, with equality only when one
vector is a nonnegative multiple of the
other. In particular its unit ball $B$ is
convex, and in fact strictly convex.
\end{corollary}

\begin{proof}
By Proposition~%
\ref{proposition-javaloyes-fund-ineq}
and $g_u(u,x) = F(u)\, dF_u(x)$ we have
$dF_u(x) \leq F(x)$ for all $u \neq 0$ and
$x$, with equality only for
$x = \lambda u$, $\lambda \geq 0$. If
$v + w \neq 0$, Euler's identity gives
$$
\begin{aligned}
F(v + w)
&= dF_{v+w}(v) + dF_{v+w}(w) \\
&\leq F(v) + F(w).
\end{aligned}
$$
Equality forces $dF_{v+w}(v)=F(v)$ and
$dF_{v+w}(w)=F(w)$, so $v=\lambda(v+w)$ and
$w=\mu(v+w)$ with $\lambda,\mu\geq 0$ and
$\lambda+\mu=1$; hence one of $v,w$ is a
nonnegative multiple of the other. If
$v + w = 0$ the inequality is clear, and
equality $0=F(v)+F(-v)$ forces $v=w=0$.
Convexity of $B$ then follows from
Proposition~%
\ref{proposition-javaloyes-convex-ball}.
For strict convexity, if $u\neq u'$ lie in
$\Sigma$ together with their midpoint, then
$F(u+u')=2$, so $u'=\lambda u$ with
$\lambda=1$, a contradiction.
\end{proof}

\noindent
In class this was phrased as: $g_v$
positive definite means that $\Sigma$ has
positive curvature, which implies that $B$
is convex.

\begin{proposition}
\label{proposition-javaloyes-indicatrix}
Let $\dim V \geq 2$. Giving a Minkowski
norm on $V$ is the same as giving a
compact, connected, smooth, strongly
convex hypersurface $\Sigma$ of $V$ that
contains $0$ in the interior of the
bounded region $B$ it encloses. The
norm is recovered as
$$
F(v) = \inf \{\lambda > 0 :
v/\lambda \in B\},
$$
and $\Sigma$ is then its indicatrix, which
is diffeomorphic to a sphere.
\end{proposition}

\noindent
For a Minkowski norm $F$, the indicatrix
is compact because $F$ is bounded below by
a positive multiple of a Euclidean norm, and
it is strongly convex by
Corollary~%
\ref{corollary-javaloyes-strong-convexity};
the map $v \mapsto v/F(v)$ is a
diffeomorphism from any Euclidean unit
sphere onto $\Sigma$. The converse is
\cite[Theorem 2.14]{javaloyes2014definition},
whose hypotheses (each ray from $0$ meets
$\Sigma$ exactly once and transversally)
follow from the convexity of $B$, which in
turn follows from strong convexity by
Hadamard's theorem on ovaloids (not
checked).

\medskip\noindent
{\bf Conic pseudo-Minkowski norms.}
For applications to spacetimes one needs to
relax the definition: the function may be
defined only on a cone, and the fundamental
tensor may be indefinite. The notes stop
after the first line of the definition; we
complete it following
\cite[Section 2]{javaloyes2014chern}.

\begin{definition}
\label{definition-javaloyes-pseudo-norm}
Let $V$ be a finite dimensional real vector
space and $A \subset V \setminus \{0\}$ an
open \emph{conic} subset, that is,
$\lambda v \in A$ for all $v \in A$ and
$\lambda > 0$. A smooth function
$L : A \to \mathbb{R}$ is a
\emph{(conic) pseudo-Minkowski norm} if it
is positively homogeneous of degree two,
$L(\lambda v) = \lambda^2 L(v)$ for
$\lambda > 0$, and for every $v \in A$ its
fundamental tensor
$$
g_v(u,w) = \frac{1}{2}
\frac{\partial^2}{\partial t \, \partial s}
L(v + tu + sw) \Big|_{t=s=0}
$$
is nondegenerate.
\end{definition}

\noindent
A pseudo-Finsler metric on a manifold $M$
is a smooth function $L$ on an open conic
subset $A \subset TM \setminus 0$ whose
restriction to each tangent space is of
this kind \cite[Section 2]%
{javaloyes2014chern}. Particular cases:
$L(v) = g(v,v)$ for a Riemannian metric
$g$, where $g_v = g$ for all $v$;
\emph{Finsler metrics}, where $L > 0$,
$g_v$ is positive definite and
$F = \sqrt{L}$ is a Minkowski norm on each
tangent space (when $A = TM \setminus 0$);
and \emph{Lorentz-Finsler metrics}, where
$g_v$ has Lorentzian signature.

\begin{remark}
\label{remark-javaloyes-time-length}
In the applications announced in the
abstract (the Zermelo navigation problem,
the spread of a wildfire), the time needed
to travel along a path is its length
computed with a Finsler metric.
\end{remark}

\medskip\noindent
{\bf Importance of Finsler metrics.}
Given a Finsler metric $F$ on $M$ and a
piecewise smooth curve
$\gamma : [a,b] \to M$, its
\emph{$F$-length} is
$$
\ell_F(\gamma) =
\int_a^b F(\dot\gamma(s))\, ds.
$$
By positive homogeneity it does not change
under orientation preserving
reparametrizations; reversing the
orientation may change it when
$F(-v) \neq F(v)$. Finsler metrics are
closely related to \emph{anisotropy}: if
the medium is locally anisotropic, so that
the cost of moving depends on the
direction, then one has to consider a
Finsler metric instead of a Riemannian one.

\medskip\noindent
{\bf Applications of Finsler geometry.}
Among the applications, Matsumoto metrics
were mentioned. These are the metrics
$F = \alpha^2/(\alpha - \beta)$, with
$\alpha$ the norm of a Riemannian metric
and $\beta$ a one-form, found by Matsumoto
when measuring the walking time on a slope
\cite[Section 4.2.3]%
{javaloyes2014definition}. Write
$b=|\beta|_\alpha$. The metric is defined
where $\alpha>\beta$, which is all of
$TM\setminus 0$ when $b<1$; it is a Finsler
metric if and only if $b<1/2$, and for
$b\geq 1$ it is defined only on a conic
domain. (This threshold is the standard
strong convexity condition for
$(\alpha,\beta)$-metrics; it was checked
numerically, not against the cited
source.)

\medskip\noindent
{\bf Other concepts related to Finsler
metrics.}
A Finsler metric gives a distance, the
infimum of $F$-lengths of curves joining
two points, which is in general not
symmetric; geodesics, which locally
minimize the $F$-length; and a notion of
curvature, the flag curvature. The notes
of the first session end here.

\bibliography{my}
\bibliographystyle{amsalpha}

\end{document}
